\section{Non-vanishing and injectivity of the lifts}\label{sec:pf:liftinj}

\providecommand{\liftinjstr}{\operatorname{str}}
\providecommand{\liftinjX}{\mathbb{X}}
\providecommand{\liftinjLS}{L^{S}_{\mathrm{tw}}}
\providecommand{\liftinjres}{\operatorname{cf}}
\providecommand{\liftinjcut}{\mathcal{C}}
\providecommand{\liftinjup}{\Theta}
\providecommand{\liftinjP}[1]{\mathrm{#1}}

This section proves \cref{prop:families:liftinj} and three further results of the formal development which
rest on the same circle of ideas: the non-vanishing of theta lifts from a definite unitary group, forced by the
pole of a partial $L$-function, and its consequences for the injectivity of the lifts.

\subsection{Introduction}\label{sec:pf:liftinj:intro}

\paragraph{The results.}
Let $V$ be a totally definite hermitian space and $W$ a hermitian space over the CM field $E$ of a global datum,
with $\dim V=\dim W+t$. The basic principle, due to Rallis and to Kudla--Rallis
\cite{Rallis1984inj,KudlaRallis1994,Yamana2014,GanQiuTakeda2014}, is:
\emph{pole $\Rightarrow$ non-vanishing $\Rightarrow$ injectivity}. If the Satake parameters of a Hecke eigenform
$u$ on $\Uni(V)$ contain a segment $\liftinjstr(\Nm w,\chi(w),t)$ (\cref{def:families:string}), the twisted
partial standard $L$-function of $u$ has a simple pole at a point where the doubling zeta integral of $u$ is
the integral of $u\otimes\bar u$ against a Siegel Eisenstein series on the doubled group; the residue of that
Eisenstein series is a theta integral from the space $W$ of the complementary dimension, so the pole forces
some theta lift of $u$ to $\Uni(W)$ to be non-zero. The unramified theta correspondence then identifies the
parameters of the lift, which gives injectivity on Hecke eigenpackets. We prove:
\begin{enumerate}[label=(\Alph*)]
\item \Cref{pf:liftinj:main} (\cref{prop:families:liftinj}): for every real datum satisfying the standing
hypotheses, the lifts from $\Uni(V_i)$ to $\Uni(W_i)$ are non-zero on eigenvectors with a segment, a good
eigenvector with a segment and vanishing lift is zero, and the lifts are injective in the sense of
\cref{prop:families:liftinj}(c).
\item \Cref{pf:liftinj:real}: the second of these statements for the real data attached to the real places of the
construction (the formal name speaks of the descent of Galois representations at real places).
\item \Cref{pf:liftinj:third}: two bundles of literature statements imply the packages of literature statements used at
the third layer of the construction (an elliptic package, a third-layer literature package, a package of statements on
twist series and Hecke operators, and a local Langlands package).
\item \Cref{pf:liftinj:point}: the theta-local inputs, Harish-Chandra finiteness and rapid decay of cusp forms
imply the point conditions of the development for every global datum with linked inputs.
\end{enumerate}
(C) and (D) are assembly results: their proofs combine many results proved elsewhere and a few substantial
results of their own (Hecke integrality, transform and reading statements for (C); eigen-segment, local type and
row statements and the point statement for (D)), which are proved in \cref{sec:pf:liftinj:further}.

\paragraph{The main argument.}
The heart of the section is the implication
\[
\begin{gathered}
\text{simple pole of }L^S\ (\cref{pf:liftinj:pole})\\
\wedge\ \text{Euler product of the doubling integral}\ (\cref{pf:liftinj:euler})\\
\wedge\ \text{residue vanishing}\ (\cref{pf:liftinj:resvan})\\
\Longrightarrow\ \text{non-vanishing of downward lifts}\ (\cref{pf:liftinj:downne}),
\end{gathered}
\]
followed by the passage from scalar automorphic forms to the vector-valued eigenvectors of real data
(\cref{pf:liftinj:nonvan}). The residue vanishing statement is deduced from an identity of regularised Siegel--Weil
type for Siegel Eisenstein series on $\Uni(n,n)$ (\cref{pf:liftinj:rescore}): at $s=t/2$ the Eisenstein series of a
Schwartz function on $B^n$, $B=W\oplus H^t$ with $H$ a hyperbolic plane, has at most a simple pole, and its residue is
orthogonal to every finite measure on $\Uni(n,n)(\A)$ that is orthogonal to all theta functions of $W$. This is a weak form
of the first term identity of the regularised Siegel--Weil formula of Kudla--Rallis and Ichino
\cite{KudlaRallis1994,Ichino2004}; its formal proof is long (about $70\,000$ lines), and we give its architecture: a
downward induction on the order of the Laurent coefficients by means of a ``cut'' operator at an auxiliary place, together
with an identification of the residue with a theta integral after the cut.

\paragraph{Organisation.}
\Cref{sec:pf:liftinj:conv} fixes conventions and lists the literature inputs.
\Cref{sec:pf:liftinj:L} treats partial $L$-functions and the simple pole, \cref{sec:pf:liftinj:doubling} the
doubling integral. \Cref{sec:pf:liftinj:rsw} is devoted to the residue identity and \cref{sec:pf:liftinj:dn} to
the non-vanishing of downward lifts. \Cref{sec:pf:liftinj:realdata} introduces the real data and proves
\cref{prop:families:liftinj}. \Cref{sec:pf:liftinj:B,sec:pf:liftinj:C,sec:pf:liftinj:D} prove (B), (C), (D). \Cref{sec:pf:liftinj:further} proves
further ingredients: a cut with non-zero eigenvalue, a bound for spherical Eisenstein series, the Siegel--Weil formula in the convergent range, the
Fourier--Jacobi coefficients of Siegel Eisenstein series and residue measures on hermitian cells for the residue identity; the Hecke statements for (C);
and two steps of the point conditions together with the point statement for (D).

\subsection{Conventions and literature inputs}\label{sec:pf:liftinj:conv}

\paragraph{Fields and places.}
A global datum $D$ (\cref{sec:families:theta}) contains a CM extension $E/F$ with $F$ totally real; we write
$x\mapsto\bar x$ for complex conjugation, $\A_F$ for the adeles of $F$ and $\A_E=\A_F\otimes_FE$. Finite places of
$F$ are denoted $v$, finite places of $E$ are denoted $w$, and $\Nm w$ is the absolute norm. For a set $S$ of finite
places of $F$, $S_E$ is the set of places of $E$ above $S$. A \emph{Hecke character} is a continuous character
$\chi:\A_E^\times\to\C^{\times}$ of absolute value one, trivial on $E^\times$; we write $\chi(w)\in\C^\times$ for
its value on a uniformiser at $w$. Let $\eta$ be the quadratic character of $\A_F^\times$ attached to $E/F$; a Hecke
character $\chi$ is \emph{$n$-adapted} if $\chi|_{\A_F^\times}=\eta^n$. The additive character $\psi$ of $\A_F/F$ is always
\emph{standard}: its archimedean component is $x\mapsto\exp(-2\pi i\sum_\tau x_\tau)$. Some statements assume the existence of
an auxiliary Hecke character $\tilde\eta$ of $E$ which is unramified everywhere, $1$-adapted, and has archimedean exponent $1/2$
at every real place.

\paragraph{Hermitian spaces and unitary groups.}
A hermitian space $A$ over $E$ is $E^a$, $a=\dim A$, with a non-degenerate hermitian Gram matrix, a standard lattice and
recorded signatures $\mathrm{sig}_\tau(A)$ at the complex embeddings $\tau$ of $E$; it has \emph{consistent signature data} if at every
$\tau$ the Gram matrix has the recorded signature. We write $G_A=\Uni(A)(\A_F)\subseteq\GL_a(\A_E)$ for the adelic unitary group,
$\Gamma_A=\Uni(A)(F)$ for its rational points, $G_{A,\infty}$ for the archimedean component and $[G_A]=\Gamma_A\backslash G_A$. The
space $A$ is \emph{definite} if its signature is $(\dim A,0)$ at every $\tau$; then $[G_A]$ is compact (\ref{pf:liftinj:D6} below).
$A^\square=A\oplus(-A)$ is the doubled space, $\iota:G_A\times G_A\to G_{A^\square}$ the doubling embedding and
$P^\square\subseteq\Uni(A^\square)$ the Siegel parabolic subgroup (stabiliser of the diagonal). For $n\ge1$ we write
$G_n=\Uni(n,n)(\A_F)$ for the adelic points of the quasi-split unitary group of the standard skew-hermitian form of rank $2n$,
$\Gamma_n$ for its rational points and $P_n=M_nN_n$ for its Siegel parabolic, with $M_n\simeq\Res_{E/F}\GL_n$ and $N_n$ the group
of hermitian $n\times n$ matrices; $|\cdot|$ is the idelic norm.

\begin{definition}[Theta functions, Weil data and theta lifts]\label{pf:liftinj:thetadef}
Let $A,B$ be hermitian spaces and $\liftinjX(A,B)=\mathrm{Mat}_{\dim B\times\dim A}(\A_E)$, with the symplectic form
$\langle\cdot,\cdot\rangle$ attached to the two hermitian forms and the Heisenberg translations
$(\rho(x_0)f)(x)=\psi(\langle x,x_0\rangle)f(x+x_0)$.
\begin{enumerate}[label=(\alph*)]
\item A \emph{theta function} on $\liftinjX(A,B)$ is a continuous function $f$ with $f(\gamma+x)=\psi(-\langle\gamma,x\rangle)f(x)$
for all rational $\gamma\in\mathrm{Mat}_{\dim B\times\dim A}(E)$, fixed by $\rho(x_0)$ for $x_0$ in an open subgroup of finite adelic
matrices, and smooth along archimedean directions.
\item A \emph{Weil datum} for $(A,B)$ with splitting characters $(\chi_A,\chi_B)$ is a representation $\omega$ of $G_A\times G_B$ on
functions on $\liftinjX(A,B)$ which preserves theta functions, intertwines the Heisenberg translations with the natural action of
$G_A\times G_B$ on $\liftinjX(A,B)$, satisfies $(\omega(\gamma,\gamma')f)(0)=f(0)$ for rational $(\gamma,\gamma')$, is jointly continuous
and unitary for the $L^2$-norm on a fundamental domain of $\mathrm{Mat}(E)$, and acts at split places through the splitting
characters; here $\chi_A$ is a $\dim A$-adapted and $\chi_B$ a $\dim B$-adapted Hecke character.
\item The \emph{theta kernel} of a theta function $f$ is
\begin{equation}\label{eq:pf:liftinj:kernel}
\theta_\omega(g,h;f)=(\omega(g,h)f)(0),\qquad g\in G_A,\ h\in G_B,
\end{equation}
which is left $\Gamma_A\times\Gamma_B$-invariant. For a Haar measure on $G_A$ and a left $\Gamma_A$-invariant function $u$ on $G_A$
the \emph{theta lift} of $u$ to $G_B$ is
\begin{equation}\label{eq:pf:liftinj:up}
\liftinjup_\omega(f,u)(h)=\int_{[G_A]}\theta_\omega(g,h;f)\,u(g)\,dg ,
\end{equation}
and $u$ \emph{has no lift} along $\omega$ if $\liftinjup_\omega(f,u)=0$ for every theta function $f$.
\item Given lattices $L_A,L_B$ and a finite set $S$ of finite places of $F$, a theta function is \emph{spherical off $S$} if it is a
lattice vacuum vector at all places outside $S$; given Cartan involutions $s_A,s_B$ of the archimedean hermitian forms, it is a
\emph{Fock vector} if it is killed by a fixed power of the archimedean annihilation operators of the Fock model attached to
$(s_A,s_B)$.
\end{enumerate}
\formal{Formalism.IndefR.Datum.IsThetaFunO}
\formal{Formalism.IndefR.Datum.IsWeilDatumO}
\formal{Formalism.IndefR.Datum.thetaKO}
\formal{GlobalTheory.Core.Rk.IP.up}
\formal{GlobalTheory.Core.Rk.Exh.NoLift}
\formal{GlobalTheory.IndefiniteRowSideFullHecke.SphericalVolumeLift.IsSphOff}
\formal{GlobalTheory.IndefiniteRowSideFullHecke.SphericalVolumeLift.IsFockVec}
\end{definition}

\begin{definition}[Cusp forms, coefficients and Satake parameters]\label{pf:liftinj:coefdef}
Let $A$ be a hermitian space, $\theta$ a family of archimedean embeddings and $c_A$ a Cartan frame of $A$.
\begin{enumerate}[label=(\alph*)]
\item $\mathcal A_0(A)$ is the space of cusp forms on $G_A$ which are finite under the maximal compact subgroup of $G_{A,\infty}$
attached to $c_A$.
\item The \emph{Petersson coefficient} of $u_1,u_2$ is $r(u_1,u_2)(g)=\langle R(g)u_1,u_2\rangle_{[G_A]}$. For a Weil datum $\omega$ for
$(A,B)$ and a measure $\nu$ of covolume one on $\liftinjX(A,B)$ with a fundamental domain $\mathcal F$ for the rational points, the
\emph{Weil coefficient} of theta functions $\Phi_1,\Phi_2$ is
$c_\omega(\Phi_1,\Phi_2)(g)=\int_{\mathcal F}(\omega(g,1)\Phi_1)\overline{\Phi_2}\,d\nu$.
\item A function $u$ \emph{has Satake parameters} $B=(B(w))_w$ off a finite set $S$ if $u$ is spherical at every place $v\notin S$ for the
stabiliser of the standard lattice and, at each place $w$ of $E$ above $v$, the spherical Hecke operators act on $u$ through the
Satake transform of the multiset $B(w)$ of rank $\dim A$ (at split places through the elementary symmetric functions of $B(w)$, at
inert places through the base change of an unramified parameter).
\end{enumerate}
\formal{GlobalTheory.Cusp.A0}
\formal{GlobalTheory.Core.Rk.IP.rcoef}
\formal{GlobalTheory.Core.Rk.IP.wcoef}
\formal{GlobalTheory.LFun.IsParamFun}
\end{definition}

\begin{definition}[Eisenstein series, doubling integrals and Laurent coefficients]\label{pf:liftinj:eisdef}
\begin{enumerate}[label=(\alph*)]
\item For a section $\Phi(s)$ of a degenerate principal series of $G_{A^\square}$ induced from $P^\square$, the Siegel Eisenstein series is
$E(g,s;\Phi)=\sum_{\gamma\in P^\square(F)\backslash\Uni(A^\square)(F)}\Phi(s)(\gamma g)$, which converges for $\Re s>\dim A/2$.
\item The \emph{doubling integral} of a function $\mathcal E$ on $G_{A^\square}$ against $f_1,f_2$ with respect to a character $\mu$ is
\begin{equation}\label{eq:pf:liftinj:dbl}
Z(\mathcal E;f_1,f_2)=\int_{[G_A]\times[G_A]}\mathcal E(\iota(g_1,g_2))\,\overline{f_1(g_1)}\,f_2(g_2)\,
\mu(\det g_2)^{-1}\,dg_1\,dg_2 ,
\end{equation}
and the \emph{zeta integral} is $s\mapsto Z(E(\cdot,s;\Phi);f_1,f_2)$.
\item A function $h$ of one complex variable has \emph{residue $c$ at $s_0$} if it has a meromorphic extension to a neighbourhood of
$s_0$ through a right domain (an open set containing a right half-plane and stable under $z\mapsto z+x$, $x\ge0$) whose Laurent
coefficient of $(s-s_0)^{-1}$ is $c$; it has \emph{a pole of order at most $k$} if this extension has order $\ge-k$ at $s_0$.
\end{enumerate}
\formal{GlobalTheory.Eis.eis}
\formal{GlobalTheory.Eis.dblInt}
\formal{GlobalTheory.Eis.zeta}
\formal{GlobalTheory.Cont.CoeffAt}
\formal{GlobalTheory.Cont.PoleLE}
\end{definition}

\begin{definition}[Partial $L$-functions]\label{pf:liftinj:Ldef}
A \emph{family} is a function $c$ from finite places of $E$ to finite multisets of $\C^\times$. For a family $c$, a set $T$ of finite places
of $E$ and $s\in\C$ put
\[
\mathrm{eul}(B,X)=\prod_{b\in B}(1-bX),\qquad L^T(s,c)=\prod_{w\notin T}\mathrm{eul}\bigl(c(w),\Nm w^{-s}\bigr)^{-1}
\]
(an unconditional product). The family \emph{converges at $s$ off $T$} if $\mathrm{eul}(c(w),\Nm w^{-s})\ne0$ for all $w\notin T$ and
$\sum_{w\notin T}|\mathrm{eul}(c(w),\Nm w^{-s})-1|<\infty$. The twist of $c$ by $z:\{w\}\to\C^\times$ is $(zc)(w)=\{z(w)b:b\in c(w)\}$, and
the \emph{segment family} is
\begin{equation}\label{eq:pf:liftinj:seg}
\mathrm{seg}(x,t)(w)=\liftinjstr(\Nm w,x(w),t)=\bigl\{x(w)\Nm w^{\frac{t-1}2-i}:0\le i<t\bigr\}.
\end{equation}
For a set $S$ of finite places of $F$, $\zeta^S_E(s)=L^{S_E}(s,\{1\})$ is the partial Dedekind zeta function, and for a family $c$ and a
function $m$ with values in $\C^\times$,
\begin{equation}\label{eq:pf:liftinj:LStw}
\liftinjLS(s;c,m)=L^{S_E}(s,m^{-1}c).
\end{equation}
For a Hecke character $\chi$ and the family $\mathrm{par}$ of Satake parameters of a form,
$L^S(s,\mathrm{par}\times\chi)=\liftinjLS(s;\mathrm{par},\chi)$. Finally, for integers $a,b$,
\begin{equation}\label{eq:pf:liftinj:dS}
d^S_{a,b}(s)=\prod_{v\notin S}\ \prod_{i=0}^{a-1}L_v\bigl(2s+a-i,\eta^{b+i}\bigr)
\end{equation}
is the normalising factor of the Siegel Eisenstein series; here $L_v(s,\eta^k)$ is the local $L$-factor of $\eta^k$ at $v$, equal to $1$
when $v$ ramifies in $E$ and $k$ is odd.
\formal{GlobalTheory.LFun.eulerProd}
\formal{GlobalTheory.LFun.Converges}
\formal{GlobalTheory.LFun.segFam}
\formal{GlobalTheory.LFun.LSTw}
\formal{GlobalTheory.LFun.dSfun}
\end{definition}

By \lref{GlobalTheory.Core.Rk.IP.dSfun_analytic_ne_zero}, $d^S_{a,b}$ is holomorphic and nowhere zero on $\Re s>0$.

\paragraph{The literature inputs.}
The main theorem takes five propositions of the development as hypotheses. We state them in the form used.

\begin{littheorem}[{\cite{MoeglinWaldspurger1995,GelbartPiatetskiShapiroRallis1987,PiatetskiShapiroRallis1987,Li1992,Tate1967thesis,Howe1979}}]
\label{pf:liftinj:lit}
\begin{enumerate}[label=(\alph*)]
\item (Rapid decay) Smooth cusp forms of uniform moderate growth on $G_V$ are of rapid decay.
\item (Doubling method) Let $V\ne0$, $a$ a majorant of $V$ for which $G_{V^\square}$ has an Iwasawa decomposition with respect to the
associated maximal compact subgroup $\hat K$, $\mu$ a $b$-adapted Hecke character and $\varphi$ a vector of the degenerate principal
series $I(\mu,s_1)$ of $G_{V^\square}$. Then the Siegel Eisenstein series of the flat section through $\varphi$ continues
meromorphically, with singularities of finite order and polynomial growth near every point; a further conjunct is the true
proposition; and for cusp forms $f_1,f_2\in\mathcal A_0(V)$ the basic identity of the doubling method holds for $\Re s>\dim V/2$: the
doubling integral \eqref{eq:pf:liftinj:dbl} (with $\mu$) of the Eisenstein series converges and equals
$\int_{G_V}\Phi(s)(\iota(g,1))\,\overline{p(g)}\,dg$, where $p(g)=\langle R(g)f_1,f_2\rangle_{[G_V]}$; moreover the zeta integral
continues meromorphically.
\item (Unramified doubling for $\GL_n$) If $\pi$ is a representation of $\GL_n(k)$, $k$ a $p$-adic field of characteristic zero, with
Satake parameter $\alpha$ and zonal spherical function $\omega$, then for $\Re s\gg0$ the doubling zeta integral of the normalised
spherical section of the doubled group against $\omega$ is
\[
L(s+\tfrac12,\alpha)\,L(s+\tfrac12,\alpha^{-1})\prod_{i=1}^{n}(1-q^{-2s-i}).
\]
\item (Partial Hecke $L$-functions) For a number field $K$ and a finite set $T$ of finite places: (i) if $\chi$ is a unitary Hecke character
of $K$ which is not of the form $|\cdot|^{i\tau}$, then $L^T(s,\chi)$ extends to an entire function without zeros on $\Re s\ge1$; (ii)
there is an entire function $Z$ with $Z(s)=(s-1)\zeta_K^T(s)$ for $\Re s>1$ and $Z(s)\ne0$ for $\Re s\ge1$.
\item (Howe's unramified generation) In the unramified local situation (unramified quadratic extension, self-dual lattices, odd residue
characteristic), every vector of the Weil representation fixed by the hyperspecial subgroup of one member of the dual pair lies in the
span of the translates of the vacuum vector under the other member.
\end{enumerate}
\formal{GlobalTheory.Published.CuspRapidDecay}
\formal{GlobalTheory.Published.DoublingMethodLocalZetaIntegrals}
\formal{GlobalTheory.Published.UnramifiedDoublingGL}
\formal{GlobalTheory.Published.HeckeLPartialContinuation}
\formal{GlobalTheory.Hk3.HoweUnram}
\end{littheorem}

\begin{remark}\label{pf:liftinj:litrem}
Inputs (d) and (e) are theorems of the formal development: the continuation and non-vanishing of partial Hecke $L$-functions is
\lref{GlobalTheory.Published.HeckeLPartialContinuation_holds}, and Howe's unramified generation is
\lref{GlobalTheory.Hk3.howeUnram_holds}. The results below nevertheless take them as hypotheses. Of (d) only part (ii) is used in
this section. Statement (b) contains a conjunct which is the true proposition and plays no role.
\end{remark}
% ---------------------------------------------------------------------------
\subsection{Partial $L$-functions and the simple pole}\label{sec:pf:liftinj:L}

The following statement is the source of the pole. Fix number fields $F\subseteq E$ (an arbitrary
$F$-algebra structure on $E$), a set $S$ of finite places of $F$, integers $r,t\ge0$, families $c,P$ and functions
$m,x$ from the finite places of $E$ to $\C^\times$.

\begin{definition}[The simple pole statement]\label{pf:liftinj:SPdef}
The \emph{segment hypothesis} $\mathbf H(S,r,t;c,P;m,x)$ means: $S$ is finite, $r+5\le t$, and for every place $w\notin S_E$ the multiset
$c(w)$ has $r$ elements, each of absolute value at most $\Nm w^{(r-1)/2}$, $|m(w)|=|x(w)|=1$, and
\begin{equation}\label{eq:pf:liftinj:liftfam}
P(w)=x(w)m(w)^{-1}c(w)+\liftinjstr\bigl(\Nm w,x(w),t\bigr).
\end{equation}
For a function $Z$ put $\mathrm{rhs}_Z(s)=\liftinjLS(s;c,m)\prod_{i=0}^{t-1}Z\bigl(s+\tfrac{t-1}2-i\bigr)$. A
\emph{continuation of $\zeta^S_E$} is a function $Z$ meromorphic on $\C$ with $Z=\zeta^S_E$ on $\Re s>1$. A function
$\Lambda$ has a \emph{simple pole} at $s_0$ if it is meromorphic at $s_0$ of order $-1$ and $(s-s_0)\Lambda(s)$
tends to a non-zero limit as $s\to s_0$, $s\ne s_0$. The \emph{conclusion} $\mathbf C(S,r,t;c,P;m,x)$ is the conjunction of:
\begin{enumerate}[label=(\roman*)]
\item for $\Re s>(t+1)/2$: $x^{-1}P$ converges at $s$ off $S_E$, the trivial family converges off $S_E$ at
$s+\frac{t-1}2-i$ for $0\le i<t$, and $\liftinjLS(s;P,x)=\mathrm{rhs}_{\zeta^S_E}(s)$;
\item for $\Re s>(r+1)/2$: $m^{-1}c$ converges at $s$ off $S_E$ and $\liftinjLS(s;c,m)\ne0$; and
$\liftinjLS(\cdot\,;c,m)$ is holomorphic on $\Re s>(r+1)/2$;
\item a continuation of $\zeta^S_E$ exists;
\item for every continuation $Z$ of $\zeta^S_E$: $\mathrm{rhs}_Z$ is meromorphic on $\Re s>(r+1)/2$, one has
$(r+1)/2<(1+t)/2$, $\mathrm{rhs}_Z(s)=\liftinjLS(s;P,x)$ for $\Re s>(t+1)/2$, and $\mathrm{rhs}_Z$ has a simple
pole at $s_0=(1+t)/2$.
\end{enumerate}
The \emph{simple pole statement} (SP) asserts that $\mathbf H$ implies $\mathbf C$ for all data as above.
\formal{GlobalTheory.IndefiniteRowSideFullHecke.SimplePoleTwistedLift.SimplePole}
\formal{GlobalTheory.IndefiniteRowSideFullHecke.SimplePoleTwistedLift.Hyp}
\formal{GlobalTheory.IndefiniteRowSideFullHecke.SimplePoleTwistedLift.Concl}
\end{definition}

We use the following standard facts on Euler products; they are proved in the development by the usual
arguments. Say that $c$ is \emph{$(r,\kappa)$-bounded off $T$} if $|c(w)|\le r$ and $|b|\le\Nm w^{\kappa}$ for
$b\in c(w)$ and $w\notin T$.
\begin{enumerate}[label=(E\arabic*)]
\item\label{pf:liftinj:E1} If $c$ is $(r,\kappa)$-bounded off $T$ and $\Re s>\kappa+1$, then $c$ converges at $s$
off $T$, and $L^T(\cdot,c)$ is holomorphic on $\Re s>\kappa+1$. Indeed each factor $1-b\Nm w^{-s}$ is non-zero
since $|b\Nm w^{-s}|\le\Nm w^{\kappa-\Re s}<1$, one has
$|\mathrm{eul}(c(w),\Nm w^{-s})-1|\le2^r\Nm w^{\kappa-\Re s}$, and $\sum_w\Nm w^{-\sigma}<\infty$ for $\sigma>1$;
the product converges locally uniformly, so its limit is holomorphic.
\item\label{pf:liftinj:E2} If $c$ converges at $s$ off $T$, then $L^T(s,c)\ne0$ (a convergent product of non-zero
factors whose distance to $1$ is summable is non-zero), and if $c,c'$ converge at $s$ off $T$ then
$L^T(s,c+c')=L^T(s,c)L^T(s,c')$, where $c+c'$ is the multiset sum; $L^T(s,c)$ only depends on the values $c(w)$,
$w\notin T$.
\item\label{pf:liftinj:E3} Sums of bounded families are bounded: if $c$ is $(r,\kappa)$- and $c'$ is
$(r',\kappa)$-bounded then $c+c'$ is $(r+r',\kappa)$-bounded; twisting by a function of absolute value one
preserves boundedness, and $(r,\kappa)$-boundedness implies $(r,\kappa')$-boundedness for $\kappa\le\kappa'$.
\item\label{pf:liftinj:E4} If $S$ is finite then so is $S_E$.
\end{enumerate}

\begin{lemma}[Segments and zeta factors]\label{pf:liftinj:segzeta}
Let $T$ be a set of finite places of $E$ and $\Re s>(t+1)/2$. Then
$L^T(s,\mathrm{seg}(1,t))=\prod_{i=0}^{t-1}L^T\bigl(s+\tfrac{t-1}2-i,\{1\}\bigr)$, and each factor on the right
converges.
\formallabel{GlobalTheory.LiftPole.eulerProd_segFam}
\end{lemma}

\begin{proof}
For $0\le i<t$ the argument $s+\frac{t-1}2-i$ has real part $>\frac{t+1}2+\frac{t-1}2-(t-1)=1$, so the trivial
family converges there by \ref{pf:liftinj:E1}, and the product over $i$ of the corresponding convergent products
is the product over $w$ of the products over $i$. At $w$, the local factor of $\mathrm{seg}(1,t)$ is
$\prod_{i<t}(1-\Nm w^{\frac{t-1}2-i}\Nm w^{-s})^{-1}$; reindexing $i\mapsto t-1-i$ turns it into
$\prod_{i<t}(1-\Nm w^{-(s+\frac{t-1}2-i)})^{-1}$.
\end{proof}

\begin{lemma}[Continuations of partial zeta functions]\label{pf:liftinj:zetacont}
\begin{enumerate}[label=(\alph*)]
\item If $Y$ is entire and $Y(s)=(s-1)\zeta^T_K(s)$ for $\Re s>1$, then $s\mapsto(s-1)^{-1}Y(s)$ is a
continuation of $\zeta^T_K$.
\item Two continuations $Z,Z'$ of $\zeta^T_K$ agree on a punctured neighbourhood of every point of $\C$.
\end{enumerate}
\formallabel{GlobalTheory.LiftPole.isZetaCont_of_entire}
\formallabel{GlobalTheory.LiftPole.eventuallyEq_of_isZetaCont}
\end{lemma}

\begin{proof}
(a) is clear ($s-1\ne0$ for $\Re s>1$). (b) $Z-Z'$ is meromorphic on the connected set $\C$ and vanishes on the open
set $\Re s>1$, so its order at $2$, hence (by the identity theorem for meromorphic functions on a connected
domain) at every point, is infinite; that is, $Z-Z'$ vanishes on a punctured neighbourhood of every point.
\end{proof}

\begin{lemma}[Criterion for a simple pole]\label{pf:liftinj:spcrit}
If $G$ is analytic at $s_0$ with $G(s_0)\ne0$ and $\Lambda(s)=(s-s_0)^{-1}G(s)$ on a punctured neighbourhood of
$s_0$, then $\Lambda$ has a simple pole at $s_0$, with $(s-s_0)\Lambda(s)\to G(s_0)$.
\formallabel{GlobalTheory.LiftPole.hasSimplePole_of}
\end{lemma}

\begin{proof}
$\Lambda$ agrees near $s_0$ with the meromorphic function $(s-s_0)^{-1}G$, whose order at $s_0$ is $-1$ because
$G(s_0)\ne0$; and $(s-s_0)\Lambda(s)=G(s)\to G(s_0)$.
\end{proof}

\begin{proposition}[The simple pole]\label{pf:liftinj:pole}
Input (d) of \cref{pf:liftinj:lit} implies the simple pole statement (SP).
\formallabel{GlobalTheory.LiftPoleObj.simplePole_of_hecke}
\formallabel{GlobalTheory.LiftPole.simplePole_of_hecke}
\end{proposition}

\begin{proof}
Assume $\mathbf H(S,r,t;c,P;m,x)$. Since $|m(w)|=1$, the family $m^{-1}c$ is $(r,\frac{r-1}2)$-bounded off $S_E$,
and $\mathrm{seg}(1,t)$ is $(t,\frac{t-1}2)$-bounded; as $r\le t$, $m^{-1}c$ is also $(r,\frac{t-1}2)$-bounded, so
$m^{-1}c+\mathrm{seg}(1,t)$ is $(r+t,\frac{t-1}2)$-bounded off $S_E$ \ref{pf:liftinj:E3}. For $w\notin S_E$,
\eqref{eq:pf:liftinj:liftfam} gives
\begin{equation}\label{eq:pf:liftinj:twP}
x(w)^{-1}P(w)=m(w)^{-1}c(w)+\mathrm{seg}(1,t)(w).
\end{equation}

\emph{(i)} Let $\Re s>(t+1)/2=\frac{t-1}2+1$. By \ref{pf:liftinj:E1} the families $m^{-1}c$, $\mathrm{seg}(1,t)$
and their sum converge at $s$; by \eqref{eq:pf:liftinj:twP} so does $x^{-1}P$, and the convergence of the trivial
family at the shifted points is \cref{pf:liftinj:segzeta}. By \eqref{eq:pf:liftinj:twP}, \ref{pf:liftinj:E2} and
\cref{pf:liftinj:segzeta},
\[
\liftinjLS(s;P,x)=L^{S_E}(s,m^{-1}c)\,L^{S_E}(s,\mathrm{seg}(1,t))
=\liftinjLS(s;c,m)\prod_{i=0}^{t-1}\zeta^S_E\bigl(s+\tfrac{t-1}2-i\bigr).
\]

\emph{(ii)} For $\Re s>\frac{r+1}2=\frac{r-1}2+1$, \ref{pf:liftinj:E1} and \ref{pf:liftinj:E2} give convergence of
$m^{-1}c$, non-vanishing of $\liftinjLS(s;c,m)$ and holomorphy.

\emph{(iii)} By (d)(ii), applied to $K=E$ and the finite set $S_E$ \ref{pf:liftinj:E4}, there is an entire $Z_0$
with $Z_0(s)=(s-1)\zeta^S_E(s)$ for $\Re s>1$ and $Z_0(s)\ne0$ for $\Re s\ge1$; by
\cref{pf:liftinj:zetacont}(a), $(s-1)^{-1}Z_0(s)$ is a continuation of $\zeta^S_E$.

\emph{(iv)} Let $Z$ be a continuation of $\zeta^S_E$. Meromorphy of $\mathrm{rhs}_Z$ on $\Re s>(r+1)/2$ follows from
(ii) and the meromorphy of $Z$. Since $r+5\le t$, $(r+1)/2<(1+t)/2$. For $\Re s>(t+1)/2$ all arguments
$s+\frac{t-1}2-i$ have real part $>1$, where $Z=\zeta^S_E$; so $\mathrm{rhs}_Z(s)=\liftinjLS(s;P,x)$ by (i).
For the pole put $s_0=(1+t)/2$. Then $s_0+\frac{t-1}2-i=t-i$, which is $\ge2$ for $i\le t-2$ and equals $1$ for
$i=t-1$; moreover $s+\frac{t-1}2-(t-1)-1=s-s_0$. Define
\[
G(s)=\liftinjLS(s;c,m)\cdot\prod_{i=0}^{t-2}\zeta^S_E\bigl(s+\tfrac{t-1}2-i\bigr)\cdot
Z_0\bigl(s+\tfrac{t-1}2-(t-1)\bigr).
\]
$G$ is analytic at $s_0$: the first factor by (ii) since $s_0>(r+1)/2$, the middle factors by \ref{pf:liftinj:E1}
since their arguments have real part $>1$ near $s_0$, and $Z_0$ is entire. $G(s_0)\ne0$: the first factor by
(ii), the middle ones by \ref{pf:liftinj:E2} (they are convergent Euler products at $t-i\ge2$), and
$Z_0(1)\ne0$. Finally, on a punctured neighbourhood of $s_0$: $Z=\zeta^S_E$ near the points $s+\frac{t-1}2-i$,
$i\le t-2$ (real part $>1$), and $Z(u)=(u-1)^{-1}Z_0(u)$ for $u$ in a punctured neighbourhood of $1$ by
\cref{pf:liftinj:zetacont}(b) applied to $Z$ and the continuation of (iii); with $u=s+\frac{t-1}2-(t-1)$ we get
$u-1=s-s_0$ and hence $\mathrm{rhs}_Z(s)=(s-s_0)^{-1}G(s)$. \Cref{pf:liftinj:spcrit} concludes.
\end{proof}

\begin{remark}\label{pf:liftinj:polerem}
The proof uses only $r<t$ (and $t\ge1$) from the hypothesis $r+5\le t$, and only part (ii) of input (d). The
point $s_0=(1+t)/2$ for $\liftinjLS(s;P,x)$ corresponds to $s=t/2$ for the doubling zeta integral, whose
$L$-factor is evaluated at $s+\frac12$ (\cref{pf:liftinj:euler}).
\end{remark}

% ---------------------------------------------------------------------------
\subsection{The doubling integral off the level}\label{sec:pf:liftinj:doubling}

\begin{definition}[Doubling pairs and the unramified Euler product]\label{pf:liftinj:ERdef}
Let $A,B$ be hermitian spaces and $t\ge1$.
\begin{enumerate}[label=(\alph*)]
\item A \emph{doubling pair} $(\psi,\theta,A,B,\chi_A,\chi_B,s_A,s_B,t,c_A,a,\omega,\nu,\mathcal F)$ consists of: $1\le\dim A$, $\dim B=\dim A+t$,
a family $\theta$ of archimedean embeddings, consistent signature data for $A$ with Cartan frame $c_A$, Cartan involutions $s_A,s_B$ of the
archimedean forms of $A,B$ compatible with $c_A$, a majorant $a$ of $A$ normalised by $s_A$, a Weil datum $\omega$ for $(A,B)$ with characters
$(\chi_A,\chi_B)$, and a measure $\nu$ of covolume one on $\liftinjX(A,B)$ with a fundamental domain $\mathcal F$ for the rational points.
\item Put $s_0(A,B)=(\dim B-\dim A)/2$. For a lattice $L$ of $A$, $\mathrm{ht}_L(g)$ is the height of $\iota(g,1)\in G_{A^\square}$ attached to $L$
and $a$. For theta functions $\Phi_1,\Phi_2$, functions $u_1,u_2$ on $G_A$ and $s\in\C$ put
\begin{equation}\label{eq:pf:liftinj:integrand}
I_s(g)=c_\omega(\Phi_1,\Phi_2)(g)\,\mathrm{ht}_L(g)^{\,s-s_0(A,B)}\,r(u_1,u_2)(g)\qquad(g\in G_A).
\end{equation}
For a finite set $S_f$ of places of $F$, $G_A^{S_f}$ denotes the set of $g\in G_A$ whose components at $v\notin S_f$ lie in the stabiliser of
$L_v$.
\item Theta functions $\Phi_1,\Phi_2$ are \emph{unramified off $S_f$} (for self-dual lattices $L_A,L_B$) if the places above $2$ outside
$S_f$ split in $E$, $\chi_B$ is unramified off $S_f$, and $\Phi_1,\Phi_2$ are spherical off $S_f$ for $L_A,L_B$.
\item The \emph{unramified Euler product} (ER) for the global datum $D$ asserts: for every doubling pair, self-dual lattices $L_A,L_B$, a finite set
$S_f$, theta functions $\Phi_1,\Phi_2$ unramified off $S_f$, a Haar measure on $G_A$, cusp forms $u_1,u_2\in\mathcal A_0(A)$ such that $u_1$ has Satake
parameters $\mathrm{par}$ off $S_f$ and $u_2$ is spherical off $S_f$, and every $s$ with $\Re s>\dim A+1$ for which $I_s$ (with $L=L_A$) is
integrable:
\begin{equation}\label{eq:pf:liftinj:ER}
\int_{G_A}I_s(g)\,dg=\frac{L^{S_f}(s+\frac12,\mathrm{par}\times\chi_B)}{d^{S_f}_{\dim A,\dim B}(s)}
\int_{G_A^{S_f}}I_s(g)\,dg .
\end{equation}
\end{enumerate}
\formal{GlobalTheory.Core.Rk.IP.Pair}
\formal{GlobalTheory.Core.Rk.IP.integrand}
\formal{GlobalTheory.Core.Rk.IP.openOff}
\formal{GlobalTheory.Core.Rk.IP.Unram}
\formal{GlobalTheory.Core.Rk.IP.EulerRight}
\end{definition}

\begin{proposition}[Unramified computation of the doubling integral]\label{pf:liftinj:euler}
Inputs (a) and (c) of \cref{pf:liftinj:lit} imply (ER) for every global datum $D$.
\formallabel{GlobalTheory.T.doublingEulerProductFarRight_of_rapidDecay_unramified}
\formallabel{GlobalTheory.Core.Rk.IP.eulerRight_of_unramifiedDoubling}
\end{proposition}

\begin{proof}
The proof combines results proved elsewhere. By \lref{GlobalTheory.Core.Rk.EulerProd.eulerRight}, input (a) together with three local
statements implies (ER); the local statements are:
\begin{itemize}
\item \emph{spherical factorisation}: at every $v\notin S_f$ and for every $s$ there is a spherical section $\varphi_v$ of $I(\chi_B^{-1},s)$ at $v$
such that right translation of $g\mapsto c_\omega(\Phi_1,\Phi_2)(g)\mathrm{ht}_L(g)^{s-s_0}$ by $y\in\Uni(A)(F_v)$ multiplies it by
$\varphi_v(\iota(y,1))$;
\item \emph{local unramified zeta integrals}: for a self-dual lattice and an adapted Hecke character $\chi$ unramified above $v$, the local doubling
integral of the spherical section against a spherical matrix coefficient with parameters $\mathrm{par}$ equals
$L_v(s+\frac12,\mathrm{par}\times\chi)/d_v(s)$ times its value at $1$, for $\Re s>\dim A$;
\item \emph{parameter bounds}: the spherical parameters of a non-zero spherical coefficient have $\dim A$ elements of absolute value at most
$\Nm w^{\dim A/2}$.
\end{itemize}
Spherical factorisation is \lref{GlobalTheory.T.doublingSphericalFactorization_of_weilDualPair}, which derives it from the existence of
Weil representations for unitary dual pairs (\lref{GlobalTheory.weilRepresentationUnitaryPair}). At split places the local zeta integrals are
computed from input (c) in \lref{GlobalTheory.Core.Rk.IP.localZetaOn_split}, and the parameter bounds are
\lref{GlobalTheory.Core.Rk.BdSplit.paramBound_split}. At non-split places the local unramified doubling integral for unitary groups
(\lref{GlobalTheory.InertZeta.unramDoublingInert}, a closed formula for the spherical doubling integral of $\Uni(n)$ over an unramified
quadratic extension) gives the parameter bounds (\lref{GlobalTheory.T.sphericalParameterBound_nonsplit_places}) and then the local zeta
integrals (\lref{GlobalTheory.LocZetaInertAd.localZetaOn_inert}). Splitting the places into split and non-split ones gives the local
statements at all places.
\end{proof}

% ---------------------------------------------------------------------------
\subsection{The residue identity}\label{sec:pf:liftinj:rsw}

\subsubsection{Siegel Eisenstein series on $\Uni(n,n)$ in the Schr\"odinger model}\label{sec:pf:liftinj:schr}

Fix $n\ge1$ and let $G_n,\Gamma_n,P_n=M_nN_n$ be as in \cref{sec:pf:liftinj:conv}; $m(a)=\mathrm{diag}(a,{}^t\bar a^{-1})$
($a\in\GL_n(\A_E)$), $n(b)$ ($b$ hermitian) and $w_n=\bigl(\begin{smallmatrix}0&1\\-1&0\end{smallmatrix}\bigr)$. Let
$K\subseteq G_n$ be the stabiliser of the standard lattice at all finite places, $K_v=K\cap G_{n,v}$ its component at a
finite place $v$ (viewed in $G_n$), and for a finite set $S$ of finite places let $K^S\subseteq G_n$ be the subgroup of
elements whose components at all finite $v\notin S$ lie in the standard compact subgroup.

\begin{definition}[Schr\"odinger model, heights and sections]\label{pf:liftinj:schrdef}
Let $V$ be a hermitian space of dimension $r_V$ and $\chi$ a Hecke character.
\begin{enumerate}[label=(\alph*)]
\item $\mathcal S(V^n)$ is the space of functions on $\mathrm{Mat}_{r_V\times n}(\A_E)$ spanned by products of a Schwartz function of the
archimedean components and a Schwartz--Bruhat function of the finite components.
\item A \emph{Schr\"odinger Weil representation} for $(\psi,\chi)$ is a representation $\omega$ of $G_n$ on $\mathcal S(V^n)$ with
\begin{equation}\label{eq:pf:liftinj:schr}
\begin{aligned}
\omega(m(a))\Phi(x)&=\chi(\det a)\,|\det a|^{r_V/2}\,\Phi(xa),\qquad
\omega(n(b))\Phi(x)=\psi\bigl(\tfrac12\tr(b\,\langle x,x\rangle_V)\bigr)\Phi(x),\\
\omega(w_n)\Phi(x)&=\int\Phi(y)\,\psi\bigl(\tr\langle x,y\rangle_V\bigr)\,d\mu(y),
\end{aligned}
\end{equation}
where $\mu$ is a measure of covolume one for $\mathrm{Mat}_{r_V\times n}(E)$ and $\langle x,y\rangle_V={}^t\bar x\,J_Vy$. Its theta function is
\begin{equation}\label{eq:pf:liftinj:thetaS}
\theta_\omega(\Phi)(g,h)=\sum_{\xi\in\mathrm{Mat}_{r_V\times n}(E)}(\omega(g)\Phi)(h^{-1}\xi)\qquad(g\in G_n,\ h\in G_V).
\end{equation}
\item A \emph{height} on $G_n$ is a continuous function $H:G_n\to\R_{>0}$ with $H(qg)=|\det a(q)|H(g)$ for $q\in P_n$ with Levi component
$m(a(q))$; it \emph{has level} if it is right invariant under an open subgroup of the finite adelic points.
\item A function $\varphi$ on $G_n$ has \emph{law $(\chi,s)$} if $\varphi(qg)=\chi(\det a(q))|\det a(q)|^{s+n/2}\varphi(g)$ for $q\in P_n$. If
$\varphi$ has law $(\chi,s_0)$, then $\varphi_z=H^{z-s_0}\varphi$ has law $(\chi,z)$, and
$E(g,\varphi_z)=\sum_{\gamma\in(P_n\cap\Gamma_n)\backslash\Gamma_n}\varphi_z(\gamma g)$.
\end{enumerate}
\formal{GlobalTheory.IndefiniteRowSideFullHecke.ModerateGrowthSchwartz.schwartzA}
\formal{GlobalTheory.StdUnn.IsWeilSch}
\formal{GlobalTheory.StdUnn.theta}
\formal{PositivityTheory.Rg.Rsw.IsHt}
\formal{PositivityTheory.Rg.Rsw.HasLevel}
\formal{PositivityTheory.Rg.Rsw.IsLaw}
\formal{PositivityTheory.Rg.Rsw.fam}
\end{definition}

\begin{definition}[Continuation data and polar coefficients]\label{pf:liftinj:contdef}
Let $\varphi$ have law $(\chi,s_0)$ and let $H$ be a height. A \emph{continuation datum} $(U,d,\mathfrak E,\sigma)$ consists of an open
right domain $U\ni s_0$, a polynomial $d\ne0$, and a function $\mathfrak E(z,g)$ on $U\times G_n$, holomorphic in $z$, jointly continuous and
left $\Gamma_n$-invariant in $g$, with $\mathfrak E(z,g)=d(z)E(g,\varphi_z)$ for $\Re z>\sigma$. A \emph{radius} is an $\eps>0$ such that the
closed disc $\bar D(s_0,\eps)$ lies in $U$ and $d$ has no zero on $\bar D(s_0,\eps)\smallsetminus\{s_0\}$. The \emph{polar coefficients} are
\begin{equation}\label{eq:pf:liftinj:cf}
\liftinjres_k(g)=\frac1{2\pi i}\oint_{|z-s_0|=\eps}(z-s_0)^{k-1}\,\frac{\mathfrak E(z,g)}{d(z)}\,dz\qquad(k\ge1),
\end{equation}
i.e.\ the coefficients of $(z-s_0)^{-k}$ in the Laurent expansion at $s_0$ of the continuation $\mathfrak E/d$ of
$z\mapsto E(g,\varphi_z)$. The residue is $\liftinjres_1$.
\formal{PositivityTheory.Rg.Rsw.IsCont}
\formal{PositivityTheory.Rg.Rsw.IsRad}
\formal{PositivityTheory.Rg.Rsw.cf}
\end{definition}

\begin{definition}[The residue identity]\label{pf:liftinj:rescoredef}
A \emph{hyperbolic decomposition} $B=W\oplus H^t$ of hermitian spaces consists of matrices $e,e'\in\mathrm{Mat}_{\dim B\times t}(E)$ spanning
isotropic subspaces with ${}^t\bar e'J_Be=1$ and an isometric embedding $i_W$ of $W$ orthogonal to them, with $\dim W+2t=\dim B$. The
\emph{residue identity} (RC) is the following statement. Let $D$ be a global datum, $n,t\ge0$, and $B,W$ hermitian spaces with a hyperbolic
decomposition $B=W\oplus H^t$. Assume
\[
\dim W+t=n,\qquad \dim W+2<2t,\qquad \dim W<t,\qquad 1\le\dim W,
\]
that $\psi$ is a character of $\A_F/F$, $\chi$ a $\dim W$-adapted Hecke character, $\theta$ a family of archimedean embeddings, that
$\omega_B,\omega_W$ are Schr\"odinger Weil representations of $G_n$ on $\mathcal S(B^n)$, $\mathcal S(W^n)$ for $(\psi,\chi)$, $H$ a height with
level, $\Phi\in\mathcal S(B^n)$, and that $(U,d,\mathfrak E,\sigma)$ is a continuation datum at $s_0=t/2$ for $\varphi(g)=(\omega_B(g)\Phi)(0)$ and
$H$, with radius $\eps$. Then:
\begin{enumerate}[label=(\roman*)]
\item $\liftinjres_k=0$ for all $k\ge2$;
\item for every finite measure space $(X,\pi)$, measurable $c:X\to G_n$ with relatively compact image and
$\kappa\in L^1(\pi)$ such that $\int_X\theta_{\omega_W}(\Psi)(c(x),h)\,\kappa(x)\,d\pi(x)=0$ for all $\Psi\in\mathcal S(W^n)$ and
$h\in G_W$, one has $\int_X\liftinjres_1(c(x))\,\kappa(x)\,d\pi(x)=0$.
\end{enumerate}
\formal{GlobalTheory.MixModel.HypDec}
\formal{PositivityTheory.Rg.Rsw.ResCore}
\end{definition}

Thus at $s=t/2$ the Eisenstein series has at most a simple pole, and its residue lies in the closure, for the weak
topology defined by finite measures with relatively compact support, of the span of the theta functions of $W$. This
is an orthogonality form of the statement that the residue is a theta integral from $W$, the first term identity of
the regularised Siegel--Weil formula \cite{KudlaRallis1994,Ichino2004}.

\subsubsection{The ingredients}\label{sec:pf:liftinj:ingr}

\begin{definition}[Fourier coefficients, cuts, idempotents and partial integration]\label{pf:liftinj:cutdef}
\begin{enumerate}[label=(\alph*)]
\item A \emph{unipotent coordinate} is a map $x\mapsto n(x)$ from adelic hermitian matrices to $N_n$ as above; \emph{coefficient data}
$(n(\cdot),\lambda,\mathcal G)$ consist of such a map, a measure $\lambda$ of covolume one on the adelic hermitian matrices and a fundamental
domain $\mathcal G$ for the rational ones. For a hermitian $\beta\in\mathrm{Herm}_n(E)$ and a function $F$ on $G_n$,
$\mathrm{fc}_\beta(F)(g)=\int_{\mathcal G}F(n(x)g)\,\overline{\psi(\frac12\tr(x\beta))}\,d\lambda(x)$ is the $\beta$-th \emph{Fourier coefficient},
and $F$ has \emph{rank less than $k$} if $\mathrm{fc}_\beta(F)=0$ for all $\beta$ of rank $\ge k$ and all coefficient data.
\item A \emph{cut at $v$} is a finite set $\liftinjcut$ of hermitian matrices whose entries have vanishing archimedean components and vanishing
components at all finite places other than $v$, with coefficients $(c_x)_{x\in\liftinjcut}$; it acts on functions by
\begin{equation}\label{eq:pf:liftinj:cut}
(\liftinjcut F)(g)=\sum_{x\in\liftinjcut}c_x\,F(g\,n(x)),\qquad \widehat{\liftinjcut}(T)=\sum_{x\in\liftinjcut}c_x\,\psi\bigl(\tfrac12\tr(xT)\bigr),
\end{equation}
and it \emph{kills rank $<r$} if $\widehat{\liftinjcut}(\beta)=0$ for every rational hermitian $\beta$ of rank $<r$ which is integral at $v$
(that is, $\psi(\frac12\tr(x\beta))=1$ for all $x$ integral at $v$).
\item For a finite set $S$ of places, an \emph{$S$-idempotent} is an $e\in\A_E$ with archimedean components $1$ and component at $v$ equal to
$1$ if $v\in S$ and $0$ otherwise; $\Phi$ is \emph{$(S,e)$-good} if $\Phi(x)=\Phi(ex)$ when $x$ is integral outside $S$ and $\Phi(x)=0$
otherwise.
\item For a hyperbolic decomposition $B=W\oplus H^t$ and \emph{auxiliary data} $I$ (a level of $B$ with a compact group $K_I$ and its Haar
probability measure, a measure $\mu_1$ of covolume one on $\mathrm{Mat}_{t\times n}(\A_E)$, and a Haar measure $m_W$ on $G_W$), the \emph{partial
integration} $\Phi\mapsto\Phi^I$ averages $\Phi$ over $K_I$ and then integrates over the $t$ hyperbolic coordinates:
\[
\Phi^I(a)=\int_{\mathrm{Mat}_{t\times n}(\A_E)}\int_{K_I}\Phi\bigl(k^{-1}(e'x+i_Wa)\bigr)\,dk\,d\mu_1(x).
\]
\end{enumerate}
\formal{PositivityTheory.PfSWUS.CoefData}
\formal{PositivityTheory.PfSWUS.fc}
\formal{PositivityTheory.Rg.Rsw.RankLT}
\formal{PositivityTheory.Rg.Rsw.L.Cut}
\formal{PositivityTheory.Rg.Rsw.L.cutR}
\formal{PositivityTheory.Rg.Rsw.L.Cut.Kills}
\formal{PositivityTheory.Rg.Rsw.L.IsIdem}
\formal{PositivityTheory.Rg.Rsw.L.Good}
\formal{PositivityTheory.Rg.Rsw.L.IkD}
\formal{PositivityTheory.Rg.Rsw.L.ikp}
\end{definition}

Coefficient data exist (\lref{GlobalTheory.CompanionTheory.SWCoef.coefDataExists}: for every $n$ there are a unipotent coordinate, a
measure of covolume one and a fundamental domain). The proof of (RC) combines the following fifteen statements, each a theorem of the development;
in each, the data range over all global data and all objects of the indicated kinds.
\begin{enumerate}[label=(R\arabic*)]
\item\label{pf:liftinj:R1} \emph{Low rank of the polar coefficients} (\lref{PositivityTheory.Rg.Rsw.polar}). If $n=r+t$, $r+2<2t$, $\chi$ is
$r$-adapted, $\chi=1$ when $r=0$, $H$ is a height, $\varphi$ continuous of law $(\chi,t/2)$, with continuation datum and radius at $t/2$, then
$\liftinjres_k$ has rank less than $r+1$ for $k\ge1$ and rank less than $r$ for $k\ge2$; moreover the rank-$r$ coefficients of $\liftinjres_1$
are given by a residue measure applied to Fourier--Jacobi sections.
\item\label{pf:liftinj:R2} \emph{Regularity and finiteness of the polar coefficients} (\lref{PositivityTheory.Rg.Rsw.cfFJ}). Under a
continuation datum and radius, each $\liftinjres_k$ is continuous and left $\Gamma_n$-invariant, there is $N$ with $\liftinjres_k=0$ for $k>N$,
and Fourier--Jacobi coefficients commute with the contour integral \eqref{eq:pf:liftinj:cf}.
\item\label{pf:liftinj:R3} \emph{Idempotents} (\lref{PositivityTheory.Rg.Rsw.L.idemEx}). $S$-idempotents exist, and they are idempotent
and fixed by complex conjugation.
\item\label{pf:liftinj:R4} \emph{Goodness for large $S$} (\lref{PositivityTheory.Rg.Rsw.L.goodEv}). For every
$\Phi\in\mathcal S(\mathrm{Mat}_{a\times b})$ there is a finite $S_0$ such that $\Phi$ is $(S,e)$-good for all $S\supseteq S_0$ and all
$S$-idempotents $e$.
\item\label{pf:liftinj:R5} \emph{Invariance of good functions} (\lref{PositivityTheory.Rg.Rsw.L.fixW}). For a Schr\"odinger Weil
representation on $\mathcal S(V^n)$ there is a finite $S_0$ such that for $S\supseteq S_0$, every $S$-idempotent $e$ and every $(S,e)$-good
$\Phi$, $\omega(u)\Phi=\Phi$ for all $u\in K^S$.
\item\label{pf:liftinj:R6} \emph{Majorants} (\lref{PositivityTheory.Rg.Rsw.L.wMaj}). For a Schr\"odinger Weil representation,
$(g,x)\mapsto(\omega(g)\Phi)(x)$ is continuous, and for every compact $C\subseteq G_n$ there is $M\in\mathcal S(V^n)$ with
$|(\omega(g)\Phi)(x)|\le|M(x)|$ for $g\in C$.
\item\label{pf:liftinj:R7} \emph{A cut with non-zero eigenvalue} (\cref{pf:liftinj:jpos}). Let $n=r+t$,
$r\ge1$, $\chi$ $r$-adapted, and let $n(\cdot)$ be a unipotent coordinate. For every finite set $S$ there are a place $v\notin S$, a cut $\liftinjcut$
at $v$ killing rank $<r$ and an entire function $c$ with $c(t/2)\ne0$ such that for every Haar probability measure on $K_v$, every $z$ and every
continuous right $K_v$-invariant $f$ of law $(\chi,z)$: $\int_{K_v}(\liftinjcut f)(gk)\,dk=c(z)f(g)$.
\item\label{pf:liftinj:R8} \emph{Cuts on polar coefficients} (\cref{pf:liftinj:cfrel} below). Under a continuation datum and
radius at $s_0$ for $\varphi$ of law $(\chi,s_0)$, a cut $\liftinjcut$ at $v$ and an entire $c$ with
$\int_{K_v}(\liftinjcut\varphi_z)(gk)\,dk=c(z)\varphi_z(g)$ for all $z,g$: if $k\ge1$ and $\liftinjres_j=0$ for all $j>k$, then
$\int_{K_v}(\liftinjcut\,\liftinjres_k)(gu)\,du=c(s_0)\liftinjres_k(g)$.
\item\label{pf:liftinj:R9} \emph{Invariance of polar coefficients} (\lref{PositivityTheory.Rg.Rsw.L.cfInv}). If $H$ and $\varphi$ are right
invariant under $u$, so is every $\liftinjres_k$.
\item\label{pf:liftinj:R10} \emph{Cuts kill low rank} (\lref{PositivityTheory.Rg.Rsw.L.cutLow}). If $\liftinjcut$ is a cut at $v$ killing
rank $<r$ and $F$ is continuous, left $\Gamma_n$-invariant and right $K_v$-invariant, then: $\liftinjcut F=0$ if $F$ has rank less than $r$;
$\liftinjcut F$ is right invariant under an open subgroup at $v$; and for $g$ with $g_v=1$,
$\mathrm{fc}_\beta(\liftinjcut F)(g)=\widehat{\liftinjcut}(\beta)\,\mathrm{fc}_\beta(F)(g)$, which vanishes if $\rank\beta<r$.
\item\label{pf:liftinj:R11} \emph{The residue after a cut} (\lref{PositivityTheory.Rg.Rsw.L.cutId}). In the setting of (RC) with auxiliary
data $I$, a cut $\liftinjcut$ at $v$ killing rank $<\dim W$ such that $\omega_B(k)\Phi=\Phi$ and $H(gk)=H(g)$ for $k\in K_v$, a fundamental domain
$\mathcal F_W$ for $\Gamma_W$ in $G_W$, $\Psi=\Phi^I$ and $\Psi'=\sum_{x\in\liftinjcut}c_x\,\omega_W(n(x))\Psi$, there is $\kappa'\in\C$ with
\begin{equation}\label{eq:pf:liftinj:cutid}
(\liftinjcut\,\liftinjres_1)(g)=\kappa'\int_{\mathcal F_W}\theta_{\omega_W}(\Psi')(g,h)\,dm_W(h)\qquad(g\in G_n).
\end{equation}
\item\label{pf:liftinj:R12} \emph{Bounds for cut theta functions} (\lref{PositivityTheory.Rg.Rsw.L.thBdd}). For $\Psi$ fixed by $K_v$ and
$\Psi'$ as in \ref{pf:liftinj:R11}: $\theta_{\omega_W}(\Psi')$ is jointly continuous, left $\Gamma_n$-invariant in $g$, invariant under an open
subgroup at $v$, bounded on $C\times G_W$ for every compact $C\subseteq G_n$, and given by an explicit sum over rank-$\dim W$ matrices along $n(x)g$.
\item\label{pf:liftinj:R13} \emph{Volume and Fourier coefficients of theta integrals} (\lref{PositivityTheory.Rg.Rsw.L.thCoef}). For a Haar
measure $m_W$ and a fundamental domain $\mathcal F_W$ for $\Gamma_W$: $m_W(\mathcal F_W)<\infty$; and if $\vartheta(\Psi)$ denotes the integral over
$\mathcal F_W$ of the sum of $\Psi(h^{-1}\xi)$ over the rational $\xi$ of rank $\dim W$, then the $\beta$-th Fourier coefficient of
$x\mapsto\vartheta(\psi(\frac12\tr(x\langle\cdot,\cdot\rangle_W))\Psi)$ vanishes if $\beta$ is not of the form $\langle\xi,\xi\rangle_W$ with
$\rank\xi=\dim W$, and is an orbital integral of $\Psi$ otherwise.
\item\label{pf:liftinj:R14} \emph{Partial integration preserves Schwartz spaces} (\lref{PositivityTheory.Rg.Rsw.L.ikSch}):
$\Phi^I\in\mathcal S(W^n)$ for $\Phi\in\mathcal S(B^n)$.
\item\label{pf:liftinj:R15} \emph{Partial integration intertwines} (\lref{PositivityTheory.Rg.Rsw.L.ikEq2}):
$\omega_W(g)\Phi^I=(\omega_B(g)\Phi)^I$.
\end{enumerate}
In the development \ref{pf:liftinj:R1}, \ref{pf:liftinj:R2}, \ref{pf:liftinj:R3}, \ref{pf:liftinj:R6}, \ref{pf:liftinj:R9},
\ref{pf:liftinj:R14}, \ref{pf:liftinj:R15} are proved directly; \ref{pf:liftinj:R4} from \ref{pf:liftinj:R3}; \ref{pf:liftinj:R7} is
\cref{pf:liftinj:jpos} below;
\ref{pf:liftinj:R5} from \ref{pf:liftinj:R3}, \ref{pf:liftinj:R6} and a generation statement for $K$; \ref{pf:liftinj:R8} is
\cref{pf:liftinj:cfrel} below (its formal statement takes the existence of heights as a hypothesis, which is not used); \ref{pf:liftinj:R10} from a
Fourier inversion statement and weak approximation at one place; \ref{pf:liftinj:R12} from
\ref{pf:liftinj:R6}, a lattice bound and weak approximation; \ref{pf:liftinj:R13} from the lattice bound. The bulk of the proof of (RC) lies in
\ref{pf:liftinj:R11}, which is deduced from \ref{pf:liftinj:R1}, \ref{pf:liftinj:R2}, \ref{pf:liftinj:R9}, \ref{pf:liftinj:R10},
\ref{pf:liftinj:R12}--\ref{pf:liftinj:R15}, a coefficient identity (\lref{GlobalTheory.CompanionTheory.RkCoef.target}, expressing the
rank-$\dim W$ Fourier coefficients of the residue through a residue measure) and the identification of that residue measure with the
theta-side measure (\lref{PositivityTheory.Rg.Rsw.L.muId}). The latter rests on the Siegel--Weil formula for hermitian forms in the
convergent range (\cref{pf:liftinj:sw}), the fibration over a Gram matrix (\cref{pf:liftinj:fib}), a bound for spherical Eisenstein series
(\cref{pf:liftinj:sph}), the Fourier--Jacobi coefficients of Siegel Eisenstein series (\cref{pf:liftinj:fj}) and the description of residue measures on
hermitian cells and of Eisenstein measures (\cref{pf:liftinj:mucell}, \cref{pf:liftinj:eismeas}); these and \ref{pf:liftinj:R7} are proved in
\cref{sec:pf:liftinj:further}, and \ref{pf:liftinj:R8} is proved next. The remaining items are cited.

\begin{lemma}[Cuts on polar coefficients]\label{pf:liftinj:cfrel}
Statement \ref{pf:liftinj:R8} holds: let $H$ be a height, $\varphi$ continuous of law $(\chi,s_0)$, $(U,d,\mathfrak E,\sigma)$ a continuation datum with
radius $\eps$, $\liftinjcut$ a cut at $v$, $\nu_v$ a Haar probability measure on $K_v$ and $c$ an entire function with
$\int_{K_v}(\liftinjcut\varphi_z)(gk)\,d\nu_v(k)=c(z)\varphi_z(g)$ for all $z,g$. If $k\ge1$ and $\liftinjres_j=0$ for all $j>k$, then
$\int_{K_v}(\liftinjcut\,\liftinjres_k)(gu)\,d\nu_v(u)=c(s_0)\liftinjres_k(g)$ for all $g$.
\formallabel{PositivityTheory.Rg.Rsw.L.cfRel}
\end{lemma}

\begin{proof}
\emph{Step 1: the cut commutes with the continuation.} For $\Re z>\max(\sigma,n/2)$ the Eisenstein series converges absolutely and locally
uniformly, dominated by a majorant along the finitely many translates by $n(x)$, $x\in\liftinjcut$, and $K_v$; hence the cut and the average over $K_v$
may be taken term by term, and the hypothesis gives
$\int_{K_v}(\liftinjcut\,E(\cdot,\varphi_z))(gu)\,d\nu_v(u)=c(z)E(g,\varphi_z)$. Multiplying by $d(z)$:
$\int_{K_v}(\liftinjcut\,\mathfrak E(z,\cdot))(gu)\,d\nu_v(u)=c(z)\mathfrak E(z,g)$ for $\Re z$ large. Both sides are holomorphic in $z\in U$ (the left side
by joint continuity of $\mathfrak E$, compactness of $K_v$ and holomorphy of $\mathfrak E(\cdot,h)$), and $U$ is connected and contains a right half-plane;
so the identity holds for all $z\in U$.

\emph{Step 2: the contour integral.} By linearity of the finite sum \eqref{eq:pf:liftinj:cut},
$(\liftinjcut\,\liftinjres_k)(gu)=\frac1{2\pi i}\oint(z-s_0)^{k-1}(\liftinjcut\,\mathfrak E(z,\cdot))(gu)/d(z)\,dz$. The integrand is jointly continuous and
bounded on the circle $\times K_v$, so the integral over $K_v$ may be taken inside the contour integral, and Step 1 gives
\[
\int_{K_v}(\liftinjcut\,\liftinjres_k)(gu)\,d\nu_v(u)=\frac1{2\pi i}\oint_{|z-s_0|=\eps}(z-s_0)^{k-1}c(z)\frac{\mathfrak E(z,g)}{d(z)}\,dz .
\]
\emph{Step 3: expanding $c$.} Write $c(z)=\sum_{m\ge0}c_m(z-s_0)^m$, uniformly convergent on the circle. The term $m$ contributes
$c_m\oint(z-s_0)^{k+m-1}\mathfrak E(z,g)/d(z)\,dz=2\pi i\,c_m\liftinjres_{k+m}(g)$, which vanishes for $m\ge1$ by hypothesis. Hence the right side equals
$c_0\liftinjres_k(g)=c(s_0)\liftinjres_k(g)$.
\end{proof}

\subsubsection{The assembly}

\begin{theorem}[Residue identity]\label{pf:liftinj:rescore}
The residue identity (RC) holds.
\formallabel{PositivityTheory.Rg.Rsw.resCore}
\end{theorem}

The theorem is the combination of the ingredients \ref{pf:liftinj:R1}--\ref{pf:liftinj:R15} with the following
proposition, whose proof we give in full.

\begin{proposition}\label{pf:liftinj:rescoreassembly}
The statements \ref{pf:liftinj:R1}--\ref{pf:liftinj:R15} imply (RC).
\formallabel{PositivityTheory.Rg.Rsw.L.resCore}
\end{proposition}

\begin{proof}
Let the data be as in \cref{pf:liftinj:rescoredef}, and write $r=\dim W$. Since $\dim B=r+2t=n+t$ and $\dim B\equiv r\pmod 2$, $\chi$ is also
$\dim B$-adapted.

\emph{Step 1: the section.} By \ref{pf:liftinj:R6}, $\varphi(g)=(\omega_B(g)\Phi)(0)$ is continuous, and by \eqref{eq:pf:liftinj:schr} it has law
$(\chi,t/2)$: for $q\in P_n$ with Levi component $m(a)$, $\varphi(qg)=\chi(\det a)|\det a|^{\dim B/2}\varphi(g)$, and $\dim B/2=t/2+n/2$. Since
$n=r+t$, $r+2<2t$ and $r\ge1$, \ref{pf:liftinj:R1} applies and gives
\begin{equation}\label{eq:pf:liftinj:P2}
\liftinjres_k\ \text{has rank less than } r\qquad(k\ge2).
\end{equation}

\emph{Step 2: an auxiliary place.} Let $U_H$ be an open subgroup of finite adelic points under which $H$ is right invariant, and choose a finite set
$S_1$ with $K^{S_1}\subseteq U_H$ (\lref{GlobalTheory.PlaceGen.exists_levelOff_le}: every open subgroup contains the level subgroup off some
finite set). Let $S_2$ be given by \ref{pf:liftinj:R5} for $\omega_B$ and $S_3$ by \ref{pf:liftinj:R4} for $\Phi$, and put $S=S_1\cup S_2\cup S_3$.
Fix coefficient data $(n(\cdot),\lambda,\mathcal G)$. By \ref{pf:liftinj:R7} (applicable as $n=r+t$, $r\ge1$) there are $v\notin S$, a cut
$\liftinjcut$ at $v$ killing rank $<r$ and an entire $c$ with $c(t/2)\ne0$ having the averaging property of \ref{pf:liftinj:R7}. Let $e$ be an
$S$-idempotent \ref{pf:liftinj:R3}. Then $\Phi$ is $(S,e)$-good by \ref{pf:liftinj:R4}, so $\omega_B(k)\Phi=\Phi$ for $k\in K_v\subseteq K^S$ by
\ref{pf:liftinj:R5}; and $H(gk)=H(g)$ for $k\in K_v\subseteq K^{S_1}\subseteq U_H$. Hence $\varphi(gk)=\varphi(g)$ and $\varphi_z(gk)=\varphi_z(g)$
for $k\in K_v$.

\emph{Step 3: the cut on the coefficients.} Let $\nu_v$ be the Haar probability measure of $K_v$. For each $z$ the function $\varphi_z$ is
continuous, of law $(\chi,z)$ and right $K_v$-invariant, so $\int_{K_v}(\liftinjcut\varphi_z)(gk)\,d\nu_v(k)=c(z)\varphi_z(g)$. By
\ref{pf:liftinj:R8},
\begin{equation}\label{eq:pf:liftinj:rel}
k\ge1,\ \liftinjres_j=0\ \text{for all } j>k\ \Longrightarrow\
\int_{K_v}(\liftinjcut\,\liftinjres_k)(gu)\,d\nu_v(u)=c(t/2)\,\liftinjres_k(g)\quad(g\in G_n).
\end{equation}
By \ref{pf:liftinj:R2} each $\liftinjres_k$ is continuous and left $\Gamma_n$-invariant and there is $N$ with $\liftinjres_k=0$ for $k>N$; by
\ref{pf:liftinj:R9} each $\liftinjres_k$ is right $K_v$-invariant.

\emph{Step 4: no higher poles.} We show by induction on $m\ge0$: $\liftinjres_k=0$ for all $k\ge2$ with $N<k+m$. For $m=0$ this is the choice of $N$.
Let $m\ge0$, and let $k\ge2$ with $N<k+m+1$. For $j>k$ we have $N<j+m$, so $\liftinjres_j=0$ by induction, and \eqref{eq:pf:liftinj:rel} gives
$\int_{K_v}(\liftinjcut\,\liftinjres_k)(gu)\,d\nu_v(u)=c(t/2)\liftinjres_k(g)$. By \eqref{eq:pf:liftinj:P2} and \ref{pf:liftinj:R10} (applied to
$F=\liftinjres_k$), $\liftinjcut\,\liftinjres_k=0$; hence $c(t/2)\liftinjres_k=0$ and $\liftinjres_k=0$. Taking $m=N+1$ proves (i).

\emph{Step 5: the residue after the cut.} Choose a good level of $B$ with compact group $K_I$ and its Haar probability measure
(\lref{GlobalTheory.RegSch.exists_good_frameLat}: such levels exist for every hermitian space), a measure $\mu_1$ of covolume one on
$\mathrm{Mat}_{t\times n}(\A_E)$ (\lref{GlobalTheory.MixModel.exists_isCovolOne}: such measures exist) and a Haar measure $m_W$ on $G_W$; call
these data $I$, and let $\mathcal F_W$ be a fundamental domain for $\Gamma_W$ in $G_W$ (\lref{GlobalTheory.FundDomDiscrete.exists_fundDom}:
the discrete subgroup $\Gamma_W$ has a measurable fundamental domain). Put $\Psi=\Phi^I\in\mathcal S(W^n)$ \ref{pf:liftinj:R14} and
$\Psi'=\sum_{x\in\liftinjcut}c_x\,\omega_W(n(x))\Psi$. By \ref{pf:liftinj:R11} there is $\kappa'$ with \eqref{eq:pf:liftinj:cutid}. By
\ref{pf:liftinj:R15}, $\omega_W(k)\Psi=(\omega_B(k)\Phi)^I=\Phi^I=\Psi$ for $k\in K_v$. By \ref{pf:liftinj:R12}, $\theta_{\omega_W}(\Psi')$ is
continuous and bounded on $C\times G_W$ for compact $C$, and by \ref{pf:liftinj:R13}, $m_W(\mathcal F_W)<\infty$.

\emph{Step 6: orthogonality.} Let $(X,\pi,c_X,\kappa)$ be as in (ii). The set $C=\overline{c_X(X)}\cdot K_v$ is compact; let $b$ bound
$|\theta_{\omega_W}(\Psi')|$ on $C\times G_W$. By Step 4, \eqref{eq:pf:liftinj:rel} holds for $k=1$, and with \eqref{eq:pf:liftinj:cutid} and
Fubini on $K_v\times\mathcal F_W$,
\[
\liftinjres_1(c_X(x))\,\kappa(x)=\frac{\kappa'}{c(t/2)}\int_{K_v\times\mathcal F_W}
\theta_{\omega_W}(\Psi')\bigl(c_X(x)u,h\bigr)\,\kappa(x)\,d(\nu_v\times m_W)(u,h).
\]
The integrand on $X\times K_v\times\mathcal F_W$ is measurable and bounded by $b|\kappa(x)|$, which is integrable for the finite measure
$\pi\times\nu_v\times m_W|_{\mathcal F_W}$. By Fubini,
\[
\int_X\liftinjres_1(c_X(x))\kappa(x)\,d\pi(x)=\frac{\kappa'}{c(t/2)}\int_{K_v\times\mathcal F_W}
\Bigl(\int_X\theta_{\omega_W}(\omega_W(u)\Psi')(c_X(x),h)\,\kappa(x)\,d\pi(x)\Bigr)\,d(u,h),
\]
using $\theta_{\omega_W}(\Psi')(gu,h)=\theta_{\omega_W}(\omega_W(u)\Psi')(g,h)$, which follows from \eqref{eq:pf:liftinj:thetaS}. The inner integral
vanishes by the hypothesis of (ii) applied to $\omega_W(u)\Psi'$ and $h$. This proves (ii).
\end{proof}

\begin{proof}[Proof of \cref{pf:liftinj:rescore}]
Apply \cref{pf:liftinj:rescoreassembly}; its fifteen hypotheses are the theorems cited in \ref{pf:liftinj:R1}--\ref{pf:liftinj:R15}, instantiated
as described after the list.
\end{proof}

\subsubsection{Residue vanishing on the doubled group}\label{sec:pf:liftinj:resvansec}

We now transport (RC) to the doubled group $G_{A^\square}$ of a definite space $A$.

\begin{definition}[Siegel--Weil sections and residue vanishing]\label{pf:liftinj:RVdef}
\begin{enumerate}[label=(\alph*)]
\item For a Weil datum $\omega^\square$ for $(A^\square,W)$ with characters $(1,\chi)$, a Haar measure on $G_A$ and a function $u$ on $G_A$, $u$
\emph{kills the theta kernels} of $\omega^\square$ if
\begin{equation}\label{eq:pf:liftinj:kills}
Z\bigl(\theta_{\omega^\square}(\cdot,h;\varphi);\bar u,\bar u\bigr)=0
\qquad\text{for every theta function }\varphi\text{ on }\liftinjX(A^\square,W)\text{ and }h\in G_W,
\end{equation}
the doubling integral \eqref{eq:pf:liftinj:dbl} being taken with respect to $\chi^{-1}$.
\item For a Weil datum $\omega_B^\square$ for $(A^\square,B)$ with characters $(1,\chi)$, a measure $\nu$ of covolume one on $\liftinjX(A,B)$ with
fundamental domain $\mathcal F$ and a theta function $\varphi$ on $\liftinjX(A^\square,B)$, the \emph{Siegel--Weil section} is
\begin{equation}\label{eq:pf:liftinj:swsec}
\varphi^{\mathrm{SW}}(g)=\int_{\mathcal F}\bigl(\omega_B^\square(g,1)\varphi\bigr)(\ell(x))\,d\nu(x)\qquad(g\in G_{A^\square}),
\end{equation}
where $\ell:\liftinjX(A,B)\to\liftinjX(A^\square,B)$ is the embedding along the diagonal Lagrangian. For a majorant $a$ of $A$, $\hat K$ denotes
the associated maximal compact subgroup of $G_{A^\square}$, and $E(g,s;\varphi^{\mathrm{SW}})$ is the Siegel Eisenstein series of the flat section
through $\varphi^{\mathrm{SW}}$, normalised at $s_0(A,B)=(\dim B-\dim A)/2$.
\item \emph{Residue vanishing} (RV) asserts: let $A$ be definite with definite data $(\theta,c_A,a)$ (\cref{pf:liftinj:Defdef}), assume that the
places above $2$ split in $E$, and let $B,W$ have consistent signature data, with a hyperbolic decomposition $B=W\oplus H^t$, $\dim W>0$, $t>0$,
$\dim A=\dim W+t$, $\dim W+2<2t$ and $\dim W<t$. Let $\chi$ be a character, $\omega_B^\square$ and $\omega_W^\square$ Weil data for $(A^\square,B)$
and $(A^\square,W)$ with characters $(1,\chi)$, $\nu$ of covolume one on $\liftinjX(A,B)$, $m$ a Haar measure on $G_A$, $u$ continuous, left
$\Gamma_A$-invariant and killing the theta kernels of $\omega_W^\square$, and $\varphi$ a $\hat K$-finite theta function on $\liftinjX(A^\square,B)$.
Then there is a function $E_1$ on $G_{A^\square}$ such that at every point $g$ the function $s\mapsto E(g,s;\varphi^{\mathrm{SW}})$ has residue
$E_1(g)$ at $s_0=s_0(A,B)$, and $Z(E_1;\bar u,\bar u)=0$ (with respect to $\chi^{-1}$).
\end{enumerate}
\formal{GlobalTheory.Core.Rk.Exh.Kills}
\formal{GlobalTheory.IndefiniteRowSideFullHecke.SphericalZero.swSec}
\formal{GlobalTheory.Eis.EisCoeffAt}
\formal{PositivityTheory.Rg.Rsw.Kplug.ResidueVanishing}
\end{definition}

The proof uses four auxiliary statements, which we prove, and some transport statements.

\begin{lemma}[Laurent coefficients from contour integrals]\label{pf:liftinj:K2}
Let $U$ be an open right domain, $e$ holomorphic on $U$, $d\ne0$ a polynomial, $\eps$ a radius at $s_0$ for $(U,d)$,
and $h$ a function with $e(z)=d(z)h(z)$ for $\Re z>\sigma$. If
$\frac1{2\pi i}\oint_{|z-s_0|=\eps}(z-s_0)^{k-1}\frac{e(z)}{d(z)}dz=0$ for all $k\ge2$, then $h$ has a pole of order at most one at
$s_0$ with residue $\frac1{2\pi i}\oint_{|z-s_0|=\eps}\frac{e(z)}{d(z)}dz$.
\formallabel{PositivityTheory.Rg.Rsw.Kplug.continuation_pole_order}
\end{lemma}

\begin{proof}
$e/d$ is meromorphic on $U$ and equals $h$ on a right half-plane, so it is the meromorphic extension of $h$ through $U$
(unique since $U$ is connected). On the punctured disc $0<|z-s_0|\le\eps$ it is holomorphic, so by Laurent's theorem its
coefficient of $(z-s_0)^{-k}$ is $\frac1{2\pi i}\oint(z-s_0)^{k-1}e/d\,dz$. These vanish for $k\ge2$, so the order at $s_0$ is
at least $-1$ and the residue is the $k=1$ integral.
\end{proof}

\begin{lemma}[Rescaling a simple pole]\label{pf:liftinj:K2b}
If $h$ has a pole of order at most one at $s_0$ with residue $c$, $\rho>0$, and $h'(s)=\rho^{s_0-s}h(s)$ for $\Re s>\sigma$, then
$h'$ has residue $c$ at $s_0$.
\formallabel{PositivityTheory.Rg.Rsw.Kplug.simple_pole_rescaling}
\end{lemma}

\begin{proof}
$s\mapsto\rho^{s_0-s}=\exp((s_0-s)\log\rho)$ is entire with value $1$ at $s_0$; the product of a function with at most a
simple pole and residue $c$ by a function holomorphic at $s_0$ with value $1$ has residue $c$, and the product is the
meromorphic extension of $h'$.
\end{proof}

\begin{lemma}[Regularity of the polar coefficients]\label{pf:liftinj:K2c}
Under a continuation datum and a radius, every $\liftinjres_k$ is continuous and left $\Gamma_n$-invariant.
\formallabel{PositivityTheory.Rg.Rsw.Kplug.continuous_invariant_coefficients}
\end{lemma}

\begin{proof}
On the circle $|z-s_0|=\eps$ the polynomial $d$ does not vanish, and $(z,g)\mapsto\mathfrak E(z,g)$ is jointly continuous, so
\eqref{eq:pf:liftinj:cf} is a parametric integral over a compact interval of a jointly continuous function, hence
continuous in $g$; invariance follows from the invariance of $\mathfrak E(z,\cdot)$.
\end{proof}

For the transport we use the following objects, constructed elsewhere in the development: the \emph{standard unitary chart}, an isomorphism
$\kappa:G_{A^\square}\to G_n$, $n=\dim A$, attached to a skew element of $E$, which maps the Siegel parabolic of $A^\square$ onto $P_n$ and rational
points onto rational points, together with $p\in P^\square$ conjugating $\hat K$ into $K$ under $\kappa$
(\lref{GlobalTheory.T.exists_siegel_parabolic_conj_doubled_compact} and \lref{GlobalTheory.T.exists_standardUnitaryChart_doubledSpace}:
such a chart and such a $p$ exist); a shifted character $\psi_\lambda(x)=\psi(\lambda x)$, $\lambda\in F$, with the archimedean parameters of the
chart (\lref{GlobalTheory.T.shiftedAdditiveCharacter_archimedeanParameters}); and, for a Weil datum $\omega^\square$ for $(A^\square,V)$, a
Schr\"odinger Weil representation $\omega_S$ of $G_n$ on $\mathcal S(V^n)$ for $(\psi_\lambda,\chi^{-1})$ transported through $\kappa$
(\lref{GlobalTheory.T.exists_weilDatum_transport_schwartzBruhat_standard}: such a transport exists;
\lref{GlobalTheory.T.transportedSchrodinger_standardLeviAction}, \lref{GlobalTheory.T.transportedSchrodinger_siegelUnipotentAction}
and \lref{GlobalTheory.T.weyl_formula_standardUnitaryGroup_weilDatum}: it satisfies the three formulas \eqref{eq:pf:liftinj:schr}). The
\emph{standard height} of the chart is $H_\kappa(g)=\mathrm{ht}(\kappa^{-1}(g)p^{-1})\,|a(p)|$, where $\mathrm{ht}$ is the Iwasawa height of
$G_{A^\square}$ for $\hat K$ and $a(p)$ is the Levi determinant of $p$.

\begin{lemma}[Theta kernels through the chart]\label{pf:liftinj:K4}
In this situation let $u$ be continuous, left $\Gamma_A$-invariant and kill the theta kernels of $\omega^\square_W$, and let $F$ be a
continuous left $\Gamma_n$-invariant function on $G_n$ such that $\int_XF(c(x))\kappa(x)\,d\pi(x)=0$ for every finite measure space
$(X,\pi)$, measurable $c$ with relatively compact image and $\kappa\in L^1(\pi)$ with
$\int_X\theta_{\omega_S}(\Psi)(c(x),h)\kappa(x)\,d\pi(x)=0$ for all $\Psi\in\mathcal S(W^n)$, $h\in G_W$. Then
$Z(g\mapsto F(\kappa(gp));\bar u,\bar u)=0$ with respect to $\chi^{-1}$.
\formallabel{PositivityTheory.Rg.Rsw.Kplug.theta_continuation}
\end{lemma}

\begin{proof}
Since $A$ is definite, there are a measurable fundamental domain $X_A$ for $\Gamma_A$ in $G_A$ and a compact $K_A\supseteq X_A$
\ref{pf:liftinj:D6}. For continuous left-invariant $\mathcal E$ on $G_{A^\square}$ the doubling integrand
$w_{\mathcal E}(x)=\mathcal E(\iota(x))\overline{\bar u(x_1)}\bar u(x_2)\chi(\det x_2)$ is continuous, hence integrable on $K_A\times K_A$, and the
doubling integral is $\int_{X_A\times X_A}w_{\mathcal E}\,d(m\times m)$. Write $w_{\mathcal E}(x)=\mathcal E(\iota(x))w(x)$ with $w$ continuous, put
$\pi=(m\times m)|_{X_A\times X_A}$ (a finite measure) and $c(x)=\kappa(\iota(x)p)$ for $x\in K_A\times K_A$, $c(x)=1$ otherwise; $c$ is measurable
with relatively compact image, and $w\in L^1(\pi)$. For $\Psi\in\mathcal S(W^n)$ there is a theta function $f$ on $\liftinjX(A^\square,W)$ with
$\theta_{\omega_S}(\Psi)(\kappa(gp),h)=\theta_{\omega_W^\square}(g,h;f)$ for all $g,h$: take $f=\omega^\square_W(p,1)f_0$, where $f_0$ corresponds to
$\Psi$ under the Schr\"odinger transform (\lref{GlobalTheory.T.schrodingerTheta_eq_schwartzBruhat}: theta functions and Schwartz functions
correspond under this transform; \lref{GlobalTheory.T.thetaKernel_eq_tsum_schrodingerFunction}: the theta kernel is the lattice sum
\eqref{eq:pf:liftinj:thetaS} of the transform). Hence $\int\theta_{\omega_S}(\Psi)(c(x),h)w(x)\,d\pi(x)=Z(\theta_{\omega_W^\square}(\cdot,h;f);\bar u,\bar u)=0$
by \eqref{eq:pf:liftinj:kills}. The hypothesis on $F$ gives $\int F(c(x))w(x)\,d\pi(x)=0$, and this integral is
$Z(F(\kappa(\cdot\,p));\bar u,\bar u)$.
\end{proof}

\begin{proposition}[Residue vanishing]\label{pf:liftinj:resvan}
The continuation of Siegel Eisenstein series (the first part of input (b) of \cref{pf:liftinj:lit}) and (RC) imply (RV).
\formallabel{PositivityTheory.Rg.Rsw.Kplug.res_killed'}
\formallabel{PositivityTheory.Rg.Rsw.Kplug.residue_vanishes_of_siegel_eisenstein}
\end{proposition}

\begin{proof}
The continuation statement is extracted from input (b) by \lref{ResidueKilled.sec_of_dm} (the first conjunct of (b)). Let the data be as in
\cref{pf:liftinj:RVdef}(c); $\dim B=\dim W+2t$. Since $\omega_W^\square$ is a Weil datum, $\chi^{-1}$ is a $\dim W$-adapted Hecke character
(\lref{GlobalTheory.T.isHeckeChar_restrictionEta_inv}: the inverse of an adapted Hecke character is an adapted Hecke character), and
$G_{A^\square}$ has an Iwasawa decomposition for $\hat K$ (\lref{GlobalTheory.Core.Rk.IP.hasIwasawa_Khat}: this holds for every majorant).

\emph{Continuation.} The Siegel--Weil section $\varphi^{\mathrm{SW}}$ of the $\hat K$-finite theta function $\varphi$ lies in the degenerate principal
series $I(\chi^{-1},s_0)$ (\lref{GlobalTheory.SwSecMem.swSec_mem_degPS}: Siegel--Weil sections of $\hat K$-finite theta functions are
sections of this degenerate principal series). By the continuation hypothesis, $E(g,s;\varphi^{\mathrm{SW}})$ continues meromorphically with
singularities of finite order. From this one obtains (\lref{PositivityTheory.Rg.Rsw.Kplug.local_representation_smoothness}: a polynomial
clearing the singularities near a point yields a continuation datum) an open right domain $U\ni s_0$, a polynomial $d\ne0$, a radius $\eps$ and a
function $\mathfrak E(z,g)$ holomorphic in $z\in U$, jointly continuous and left invariant, with $\mathfrak E(z,g)=d(z)E(g,z;\varphi^{\mathrm{SW}})$ for
$\Re z>\dim A/2$.

\emph{Transport.} Fix the chart $\kappa$, $p$, $\lambda$ and the transported Schr\"odinger representations $\omega_{S,B}$ on $\mathcal S(B^n)$ and
$\omega_{S,W}$ on $\mathcal S(W^n)$ as above, and let $\Phi\in\mathcal S(B^n)$ correspond to $\omega_B^\square(p^{-1},1)\varphi$ under the
Schr\"odinger transform. Then (\lref{PositivityTheory.Rg.Rsw.Kplug.continuation_through_chart}, which identifies the flat section with a
translate of the standard section and compares the coset sums through $\kappa$): $\varphi_S(g)=(\omega_{S,B}(g)\Phi)(0)$ is continuous;
$(U,d,\mathfrak E_S,\dim A/2)$ with $\mathfrak E_S(z,g)=|a(p)|^{z-s_0}\mathfrak E(z,\kappa^{-1}(g)p^{-1})$ is a continuation datum for $\varphi_S$ and the
height $H_\kappa$; and
\begin{equation}\label{eq:pf:liftinj:transp}
E(g,z;\varphi^{\mathrm{SW}})=|a(p)|^{s_0-z}\,E\bigl(\kappa(gp),\varphi_{S,z}\bigr)\qquad(\Re z>\dim A/2).
\end{equation}
Moreover $H_\kappa$ is a height with level: it is continuous and positive; $H_\kappa(qg)=|\det a(q)|H_\kappa(g)$ for $q\in P_n$, because
$\kappa^{-1}$ maps $P_n$ to the Siegel parabolic compatibly with Levi determinants and the Iwasawa height is left equivariant; and $H_\kappa$ is right
invariant under the standard level subgroup, because $p\kappa^{-1}(k)p^{-1}\in\hat K$ for $k$ in it. By \cref{pf:liftinj:K2c} the polar coefficients
$\liftinjres_k$ of $(U,d,\mathfrak E_S)$ at $s_0$ are continuous and left $\Gamma_n$-invariant.

\emph{Application of (RC).} Since $s_0=(\dim B-\dim A)/2=t/2$ and $\dim A=\dim W+t$, \cref{pf:liftinj:rescoredef} applies with $n=\dim A$, the
character $\psi_\lambda$ and $\chi^{-1}$: $\liftinjres_k=0$ for $k\ge2$, and $\liftinjres_1$ is orthogonal to every finite measure orthogonal to the
theta functions of $\omega_{S,W}$.

\emph{The residue.} Put $E_1(g)=\liftinjres_1(\kappa(gp))$. For fixed $g$, apply \cref{pf:liftinj:K2} to $h(z)=E(\kappa(gp),\varphi_{S,z})$ and
$e(z)=\mathfrak E_S(z,\kappa(gp))$: $h$ has at most a simple pole at $s_0$ with residue $E_1(g)$. By \eqref{eq:pf:liftinj:transp} and
\cref{pf:liftinj:K2b} (with $\rho=|a(p)|$), $E_1(g)$ is the residue of $z\mapsto E(g,z;\varphi^{\mathrm{SW}})$ at $s_0$. Finally
$Z(E_1;\bar u,\bar u)=0$ by \cref{pf:liftinj:K4} applied to $F=\liftinjres_1$.
\end{proof}

% ---------------------------------------------------------------------------
\subsection{Non-vanishing of downward lifts}\label{sec:pf:liftinj:dn}

\begin{definition}[Definite data, level places, forms of segment type, occurrence]\label{pf:liftinj:Defdef}
Let $A$ be a hermitian space.
\begin{enumerate}[label=(\alph*)]
\item $(\theta,c_A,a)$ are \emph{definite data} for $A$ if $\theta$ is a family of archimedean embeddings, $A$ is definite with consistent
signature data, $c_A$ is a Cartan frame of $A$ for which the identity is a Cartan involution at every real place, and $a$ is a majorant of $A$
normalised by $c_A$.
\item $(\beta,Z)$ are \emph{level places} for $A$ if $\beta$ is a frame of the standard lattice of $A$, $Z$ is a finite set of places of $E$ for which an
idempotent of $\A_E$ supported on $Z$ exists, and every place of $F$ above $2$ splits in $E$. The associated \emph{principal congruence subgroup}
$\mathrm{prin}(A,\beta,Z)\subseteq G_A$ is the principal congruence subgroup of level $Z$ with respect to $\beta$, and $\mathrm{prin}_\infty(A,\beta,Z)$
is the set of $g\in G_A$ whose finite part lies in it.
\item Let $S$ be the set of places of $F$ below $Z$. A function $u$ is a \emph{form of segment type} $(r,t)$ with parameters $\mathrm{par}$ and
character $\chi_W$ if $u\in\mathcal A_0(A)$, $u$ is right invariant under $\mathrm{prin}(A,\beta,Z)$, $u$ has Satake parameters $\mathrm{par}$ off $S$,
$\chi_W$ is unramified off $S$, and $\mathbf H(S,r,t;c,\mathrm{par};\chi_W,\chi_W)$ holds for some family $c$ (\cref{pf:liftinj:SPdef}; here $\chi_W$
stands for $w\mapsto\chi_W(w)$); that is, $r+5\le t$ and $\mathrm{par}(w)=c(w)+\liftinjstr(\Nm w,\chi_W(w),t)$ with $|c(w)|=r$ and
$|b|\le\Nm w^{(r-1)/2}$ for $b\in c(w)$, $w$ off $S$.
\item For a Weil datum $\omega$ for $(A,W)$, a theta function $\Phi$ \emph{occurs} for $u$ if there is a Cartan involution $s_W$ of $W$ such that $\Phi$
is a Fock vector for $(1,s_W)$, $\Phi$ is spherical off $S$ for the frame lattices, $\omega(\kappa,1)\Phi=\Phi$ for $\kappa\in\mathrm{prin}(A,\beta,Z)$, and
for every measure $\nu$ of covolume one on $\liftinjX(A,W)$ and every Haar measure on the archimedean group $G_{A,\infty}$,
\begin{equation}\label{eq:pf:liftinj:realint}
\int_{G_{A,\infty}}c_\omega(\Phi,\Phi)(b)\,r(u,u)(b)\,db\ne0 .
\end{equation}
\item The \emph{non-vanishing of downward lifts} (DN) for a global datum $D$ asserts: let $(\theta,c_A,a)$ be definite data for $A$, let $W$ have consistent
signature data with $\dim W>0$, let $\dim A=\dim W+t$, assume that $\tilde\eta$ exists, let $\omega$ be a Weil datum for $(A,W)$ with characters
$(\chi_A,\chi_W)$, $(\beta,Z)$ level places for $A$, $m$ a Haar measure on $G_A$, $u$ a form of segment type $(\dim W,t)$ with character $\chi_W$, and
$\Phi$ a theta function occurring for $u$. Then $u$ has a lift along $\omega$: $\liftinjup_\omega(f,u)\ne0$ for some theta function $f$ on
$\liftinjX(A,W)$.
\end{enumerate}
\formal{GlobalTheory.Core.Rk.Exh.Defin}
\formal{GlobalTheory.Core.Rk.Exh.Places}
\formal{GlobalTheory.Core.Rk.Exh.Form}
\formal{GlobalTheory.Core.Rk.Exh.Occ}
\formal{GlobalTheory.Core.Rk.Exh.DownNe}
\end{definition}

\begin{lemma}[Vanishing lifts kill theta kernels]\label{pf:liftinj:kills}
Let $A$ be definite, $\omega$ a Weil datum for $(A,W)$ with characters $(\chi_A,\chi_W)$ and $\omega^\square$ one for $(A^\square,W)$ with
characters $(1,\chi_W)$, $m$ a Haar measure on $G_A$ and $u$ continuous and left $\Gamma_A$-invariant. If $u$ has no lift along
$\omega$, then $u$ kills the theta kernels of $\omega^\square$.
\formallabel{GlobalTheory.T.doubling_integrals_vanish_of_all_lifts_vanish}
\formallabel{GlobalTheory.Core.Rk.Kill.kills_of_noLift}
\end{lemma}

\begin{proof}
Let $X_A\subseteq K_A$ be a fundamental domain inside a compact set \ref{pf:liftinj:D6}; $u$ is bounded, say by $C$. Fix $h\in G_W$ and a theta
function $\varphi$ on $\liftinjX(A^\square,W)$. The doubling integrand of $\theta_{\omega^\square}(\cdot,h;\varphi)$ against $\bar u,\bar u$ is
continuous, so the doubling integral is the integral over $X_A\times X_A$, and there it equals
\[
\int_{X_A\times X_A}c(x)\,(\omega^\square(\iota(x),h)\varphi)(0)\,dx,\qquad c(x_1,x_2)=u(x_1)\overline{u(x_2)}\chi_W(\det x_2),
\]
a bounded continuous weight. By the seesaw identity for the doubling map
(\lref{GlobalTheory.T.doublingMap_thetaIdentity_any_weilData}), for theta functions $f_1,f_2$ on $\liftinjX(A,W)$ the function
$\delta(f_1,f_2)(x)=f_1(x_L)\overline{f_2(-x_R)}$ is a theta function on $\liftinjX(A^\square,W)$ and
\[
\theta_{\omega^\square}(\iota(g_1,g_2),h;\delta(f_1,f_2))=\theta_\omega(g_1,h;f_1)\,\overline{\theta_\omega(g_2,h;f_2)}\,\chi_W(\det g_2)^{-1}.
\]
For $\varphi=\delta(f_1,f_2)$ the integrand is therefore $\theta_\omega(x_1,h;f_1)u(x_1)\cdot\overline{\theta_\omega(x_2,h;f_2)u(x_2)}$, and the
integral factorises as $\liftinjup_\omega(f_1,u)(h)\overline{\liftinjup_\omega(f_2,u)(h)}=0$ (the integral over $X_A$ of a left-invariant
function being the integral over $[G_A]$). By linearity the integral vanishes for $\varphi$ in the span $M$ of the $\delta(f_1,f_2)$. Now $M$
consists of theta functions, is stable under Heisenberg translations, and every theta function on $\liftinjX(A^\square,W)$ is a uniform limit of
elements of $M$ (\lref{GlobalTheory.T.doublingMap_span_dense_heisenbergStable}); and for a finite measure, a continuous map $x\mapsto(\iota(x),h)$
with values in a compact set and a bounded continuous weight, the vanishing of $\int c(x)(\omega^\square(p(x))g)(0)$ for all $g\in M$ extends to every
theta function which is a uniform limit of elements of a Heisenberg-stable space of theta functions
(\lref{GlobalTheory.T.weilRep_average_kills_theta_of_dense_family}).
\end{proof}

\begin{lemma}[The zeta integral of a Fock pair]\label{pf:liftinj:zetaint}
Assume input (b) of \cref{pf:liftinj:lit}. Let $(\psi,\theta,A,B,\chi_A,\chi_B,s_A,s_B,t,c_A,a,\omega,\nu,\mathcal F)$ be a doubling pair,
$\omega^\square$ a Weil datum for $(A^\square,B)$ with characters $(1,\chi_B)$, $\Phi$ a Fock vector for $(A,B)$, $m$ a Haar measure and
$u\in\mathcal A_0(A)$. Put $\varphi=\delta(\Phi,\Phi)$. Then the Siegel Eisenstein series $E(\cdot,s;\varphi^{\mathrm{SW}})$ (flat at
$s_0(A,B)$) continues meromorphically with singularities of finite order and is holomorphic on $\Re s>\dim A/2$, and for $\Re s>\dim A/2$
the function $I_s$ of \eqref{eq:pf:liftinj:integrand} (with $\Phi_1=\Phi_2=\Phi$, $u_1=u_2=u$ and the frame lattice) is integrable and
\[
Z\bigl(E(\cdot,s;\varphi^{\mathrm{SW}});\bar u,\bar u\bigr)=\int_{G_A}I_s(g)\,dg\qquad(\text{with respect to }\chi_B^{-1}).
\]
\formallabel{GlobalTheory.Core.Rk.Exh.zeta_eq_integral_integrand_nc}
\end{lemma}

\begin{proof}
By the seesaw identity (\lref{GlobalTheory.T.doublingMap_thetaIdentity_any_weilData}) $\varphi$ is a theta function, and
$\varphi^{\mathrm{SW}}(\iota(g,1))=c_\omega(\Phi,\Phi)(g)$ (\lref{GlobalTheory.T.siegelWeilSection_doubling_eq_matrixCoefficient}: the
Siegel--Weil section of $\delta(\Phi_1,\Phi_2)$ restricted to $G_A\times1$ is the Weil coefficient). Since $\Phi$ is a Fock vector, $\varphi$ is
$\hat K$-finite (\lref{GlobalTheory.T.doublingVector_fockPair_isCompactFinite}); $G_{A^\square}$ has an Iwasawa decomposition for $\hat K$
(\lref{GlobalTheory.T.adelicIwasawa_doubledGroup_selfDualLattice}), so $\varphi^{\mathrm{SW}}\in I(\chi_B^{-1},s_0)$
(\lref{GlobalTheory.T.doublingSiegelWeilSection_mem_degeneratePrincipalSeries}). Input (b), applied to $\varphi^{\mathrm{SW}}$ and the cusp forms
$\bar u$ (which lie in $\mathcal A_0(A)$ by \lref{GlobalTheory.T.cuspFormsArchimedeanKFinite_star_comp_mem}: the complex conjugate of a cusp form
in $\mathcal A_0(A)$ lies in $\mathcal A_0(A)$), gives the continuation statements and, for $\Re s>\dim A/2$, the basic identity
$Z(E(\cdot,s);\bar u,\bar u)=\int_{G_A}\varphi^{\mathrm{SW}}_s(\iota(g,1))\overline{p(g)}\,dg$ with convergence, where $p$ is the Petersson coefficient
of $(\bar u,\bar u)$. Finally $\varphi^{\mathrm{SW}}_s(\iota(g,1))=c_\omega(\Phi,\Phi)(g)\,\mathrm{ht}(g)^{s-s_0}$ (flat section) and
$\overline{p(g)}=r(u,u)(g)$, so the integrand is $I_s$.
\end{proof}

\begin{lemma}[The residue of the zeta integral is non-zero]\label{pf:liftinj:resne}
Assume (SP). Let $S$ be a finite set of places of $F$, $\chi$ a Hecke character and $\mathrm{par}$ a family with
$\mathbf H(S,r,t;c,\mathrm{par};\chi,\chi)$. Let $a,b\in\N$, $C\ne0$, and let $\mathcal Z$ be a function with
$\mathcal Z(s)=\frac{L^S(s+\frac12,\mathrm{par}\times\chi)}{d^S_{a,b}(s)}C$ for $\Re s$ large. If $\mathcal Z$ has residue $z$ at $t/2$, then $z\ne0$.
\formallabel{GlobalTheory.T.zeta_residue_ne_zero_of_simple_pole_lift_L_function}
\formallabel{GlobalTheory.Core.Rk.Exh.residue_ne_zero_of_simplePole}
\end{lemma}

\begin{proof}
By (SP) (\cref{pf:liftinj:SPdef}, parts (iii), (iv)) there is a continuation $Z_0$ of $\zeta^S_E$ such that $R=\mathrm{rhs}_{Z_0}$ is meromorphic on
$\Re s>(r+1)/2$, equals $L^S(s,\mathrm{par}\times\chi)=\liftinjLS(s;\mathrm{par},\chi)$ for $\Re s>(t+1)/2$, and has a simple pole at $(1+t)/2$
with $(s-\frac{1+t}2)R(s)\to\rho\ne0$. The function $d=d^S_{a,b}$ is holomorphic and non-vanishing on $\Re s>0$. Hence $s\mapsto R(s+\frac12)C/d(s)$
is meromorphic on the right domain $\Re s>r/2$, which contains $t/2$ (as $r+5\le t$), and it agrees with $\mathcal Z$ for $\Re s$ large; so it is the
meromorphic extension of $\mathcal Z$ at $t/2$. As $s\to t/2$, $(s-\frac t2)R(s+\frac12)C/d(s)\to\rho C/d(t/2)$. By uniqueness of Laurent
coefficients, $z=\rho C/d(t/2)\ne0$.
\end{proof}

\begin{lemma}[Hyperbolic sums]\label{pf:liftinj:hypsum}
Let $W$ have consistent signature data and $t\ge0$. Let $H$ be the standard hyperbolic space of dimension $2t$ and $B=W\oplus H$. Then $H$ and $B$ have
consistent signature data, $H$ has signature $(t,t)$ at every complex embedding, $B=W\oplus H^t$ is a hyperbolic decomposition, and the Gram matrix
$J_H$ of $H$ (at each real place) is a Cartan involution of $H$.
\formallabel{GlobalTheory.T.exists_sum_hyperbolicPlanes_orthogonal_wittDecomposition}
\formallabel{GlobalTheory.HypSumWitt.exists_hypSum_orthDec_hypDec}
\end{lemma}

\begin{proof}
The signature statements hold for the standard hyperbolic space and are additive under orthogonal sums; the orthogonal and hyperbolic decompositions
of $B$ are the obvious ones. The Gram matrix satisfies $J_H^2=1$, so $s_H=J_H$ is an involution with $J_Hs_H=1$ positive definite, i.e.\ a Cartan
involution.
\end{proof}

\begin{lemma}[The level part of the zeta integral]\label{pf:liftinj:levelconst}
Let $(\theta,c_A,a)$ be definite data for $A$, $B=W\oplus H$ an orthogonal sum, $\omega_B,\omega_W$ Weil data for $(A,B)$, $(A,W)$ with characters
$(\chi_A,\chi_W)$ and $\omega_H$ one for $(A,H)$ with characters $(\chi_A,1)$, $(\beta,Z)$ level places for $A$ with $S$ the places of $F$ below $Z$, $m$
a Haar measure on $G_A$, $u\in\mathcal A_0(A)$ right invariant under $\mathrm{prin}(A,\beta,Z)$, and $\Phi$ a theta function occurring for $u$ along
$\omega_W$. Let $f\ne0$ be a theta function on $\liftinjX(A,H)$, fixed by $\omega_H(\kappa,1)$ for $\kappa\in\mathrm{prin}(A,\beta,Z)$ and by $\omega_H(b,1)$ for
$b\in G_{A,\infty}$, such that $c_{\omega_H}(f,f)(g)=0$ for every $g\in G_A^S$ outside $\mathrm{prin}_\infty(A,\beta,Z)$, and let $\nu_B$ be a measure of
covolume one on $\liftinjX(A,B)$. Then there is $C\ne0$ such that $\int_{G_A^S}I_s(g)\,dg=C$ for every $s\in\C$, where $I_s$ is formed with
$\omega_B$, $\Phi_1=\Phi_2=\Phi\boxtimes f$, $u_1=u_2=u$ and the frame lattice.
\formallabel{GlobalTheory.T.levelPlaces_integral_nonzero_constant}
\formallabel{GlobalTheory.Core.Rk.Exh.LocalConst.integral_const}
\end{lemma}

\begin{proof}
Choose measures of covolume one on $\liftinjX(A,W)$ and $\liftinjX(A,H)$. The Weil coefficient of a box product is the product of the Weil coefficients
(\lref{GlobalTheory.T.weilCoefficient_externalProduct_eq_mul_orthogonalSum}: for an orthogonal sum $B=W\oplus H$,
$c_{\omega_B}(\Phi\boxtimes f,\Phi\boxtimes f)=c_{\omega_W}(\Phi,\Phi)\,c_{\omega_H}(f,f)$). Put $F=c_{\omega_W}(\Phi,\Phi)\,c_{\omega_H}(f,f)\,r(u,u)$; it is
continuous, and right invariant under $\mathrm{prin}(A,\beta,Z)$ since $\Phi$, $f$ and $u$ are. The integration lemma for principal congruence subgroups
(\lref{GlobalTheory.T.principalCongruence_realPoints_integration_heightOne}) gives: $\mathrm{prin}_\infty(A,\beta,Z)$ is a measurable subset of
$G_A^S$, the height $\mathrm{ht}$ equals $1$ on it, and there is a Haar measure $m_\infty$ on $G_{A,\infty}$ with
$\int_{\mathrm{prin}_\infty}F'\,dm=\int_{G_{A,\infty}}F'\,dm_\infty$ for every continuous $F'$ right invariant under $\mathrm{prin}(A,\beta,Z)$. On
$G_A^S$ the integrand $I_s$ vanishes outside $\mathrm{prin}_\infty$ (because $c_{\omega_H}(f,f)$ does) and equals $F$ on $\mathrm{prin}_\infty$ (because the
height is $1$ there). Hence $\int_{G_A^S}I_s=\int_{G_{A,\infty}}F\,dm_\infty$, independently of $s$. For $b\in G_{A,\infty}$,
$c_{\omega_H}(f,f)(b)=\langle\omega_H(b,1)f,f\rangle=\langle f,f\rangle=:c$, and $c\ne0$ since $f\ne0$ is a theta function
(\lref{GlobalTheory.RawLift.VecNonzero.thetaInner_self_ne_zero}: the theta inner product of a non-zero theta function with itself is non-zero).
Therefore $C=c\int_{G_{A,\infty}}c_{\omega_W}(\Phi,\Phi)(b)\,r(u,u)(b)\,dm_\infty(b)$, which is non-zero by the occurrence condition
\eqref{eq:pf:liftinj:realint}.
\end{proof}

The proof of the next theorem also uses the following results of this section, which we cite without proof.
\begin{enumerate}[label=(D\arabic*)]
\item\label{pf:liftinj:D1} (\lref{GlobalTheory.T.exists_vacuumVector_principalCongruenceSupport}) Let $\omega_H$ be a Weil datum for
$(A,H)$ with characters $(\chi_A,1)$, $s_A,s_H$ Cartan involutions, $\dim A\le\dim H$, $(\beta,Z)$ level places, and assume that
$\tilde\eta$ exists. Then there is a theta function $f\ne0$ on $\liftinjX(A,H)$ which is a Fock vacuum vector at every real place, spherical
off $Z$, fixed by $\omega_H(\kappa,1)$ for $\kappa\in\mathrm{prin}(A,\beta,Z)$, and such that $c_{\omega_H}(f,f)(g)=0$ for every $g\in G_A^{S}$
($S$ the places below $Z$) outside $\mathrm{prin}_\infty(A,\beta,Z)$.
\item\label{pf:liftinj:D2} (\lref{GlobalTheory.T.definite_real_group_fixes_split_space_weil_vacuum}) If moreover the identity is a Cartan
involution of $A$, $\dim H=2t$ and $H$ has signature $(t,t)$ at every real place, then every Fock vacuum vector $f$ as in \ref{pf:liftinj:D1} is fixed
by $\omega_H(b,1)$ for $b\in G_{A,\infty}$.
\item\label{pf:liftinj:D3} (\lref{GlobalTheory.T.boxProducts_orthogonalSum_target_cartan_fock}) If $B=W\oplus H$ is an orthogonal sum and
$s_W,s_H$ are Cartan involutions, there are a Cartan involution $s_B$ and a self-dual lattice $L_B$ of $B$ such that the box product
$\Phi\boxtimes\Phi'$ on $\liftinjX(A,B)=\liftinjX(A,W)\oplus\liftinjX(A,H)$ of Fock vectors (resp.\ of vectors spherical off $S_f$) is a Fock
vector (resp.\ spherical off $S_f$ for $L_B$).
\item\label{pf:liftinj:D4} (\lref{GlobalTheory.T.zeta_integral_eisenstein_residue}) Let $\Phi(s)$ be a family of sections whose
Eisenstein series continues, is holomorphic on $\Re s>\dim A/2$ and has singularities of finite order with polynomial growth, let $f_1,f_2$ be of
rapid decay with convergent zeta family, and let $E_1$ be the residue of the Eisenstein series at $s_0$. Then the zeta integral
$s\mapsto Z(E(\cdot,s;\Phi);f_1,f_2)$ has residue $Z(E_1;f_1,f_2)$ at $s_0$.
\item\label{pf:liftinj:D6} (\lref{GlobalTheory.T.definite_unitary_quotient_compact_fundamentalDomain}) If $A$ is definite, $[G_A]$ is
compact, and $\Gamma_A$ has a measurable fundamental domain contained in a compact set. (Compactness is the classical criterion for anisotropic groups
\cite{Borel1963}, $\Uni(A)$ being anisotropic over $F$ when $A$ is definite.)
\end{enumerate}

\begin{theorem}[Non-vanishing of downward lifts]\label{pf:liftinj:downne}
Input (b) of \cref{pf:liftinj:lit} and (SP) imply: (ER) $\Rightarrow$ (DN) for every global datum $D$.
\formallabel{GlobalTheory.T.downNe_of_eulerRight_doubling}
\formallabel{GlobalTheory.Core.Rk.Exh.Top.downNe_k_of_hypotheses'}
\formallabel{GlobalTheory.Core.Rk.Exh.Top.downNe_of_core}
\end{theorem}

\begin{proof}
By \cref{pf:liftinj:rescore} we may use (RC), and by \cref{pf:liftinj:resvan} (with the continuation part of input (b)) (RV). Let the data be as in
\cref{pf:liftinj:Defdef}(e), and suppose that $u$ has no lift along $\omega$. Write $r=\dim W$ and let $S$ be the set of places of $F$ below $Z$.

\emph{Step 1: the doubling space.} Let $H$ and $B=W\oplus H$ be as in \cref{pf:liftinj:hypsum}, with the Cartan involution $s_H$ of $H$. The segment
hypothesis gives $c$ with $\mathbf H(S,r,t;c,\mathrm{par};\chi_W,\chi_W)$; in particular $r+5\le t$, so $\dim B=r+2t$, $\dim A=r+t\le2t=\dim H$,
$r<t$ and $r+2<2t$.

\emph{Step 2: Weil data and test vectors.} By the existence of Weil representations for unitary pairs
(\lref{GlobalTheory.weilRepresentationUnitaryPair}: for adapted Hecke characters there is a Weil datum for every pair of hermitian spaces) there are
Weil data $\omega_H$ for $(A,H)$ with characters $(\chi_A,1)$ and $\omega_B$ for $(A,B)$ with characters $(\chi_A,\chi_W)$ (the character $1$ is
$2t$-adapted and $\chi_W$ is $(r+2t)$-adapted, by parity), and doubled Weil data $\omega_B^\square$ for $(A^\square,B)$ and $\omega_W^\square$ for
$(A^\square,W)$ with characters $(1,\chi_W)$ (\lref{GlobalTheory.Core.Rk.exists_doubled_weilDatum}: a Weil datum for $(A,B)$ yields one for
$(A^\square,B)$). Let $f$ be as in \ref{pf:liftinj:D1} for $(A,H)$; by \ref{pf:liftinj:D2} it is fixed by $G_{A,\infty}$. By \ref{pf:liftinj:D3}
there are $s_B$ and a self-dual lattice $L_B$ such that $\Phi^+=\Phi\boxtimes f$ is a Fock vector for $(A,B)$ and spherical off $S$. Choose $\nu$ of
covolume one on $\liftinjX(A,B)$ (\lref{GlobalTheory.MixModel.exists_isCovolOne}: such measures exist); then
$(\psi,\theta,A,B,\chi_A,\chi_W,1,s_B,t,c_A,a,\omega_B,\nu,\mathcal F)$ is a doubling pair.

\emph{Step 3: the residue vanishes.} By \cref{pf:liftinj:kills}, $u$ kills the theta kernels of $\omega_W^\square$. The function
$\varphi=\delta(\Phi^+,\Phi^+)$ is a theta function on $\liftinjX(A^\square,B)$ (\lref{GlobalTheory.Core.Rk.DeltaTheta.main}: $\delta$ maps pairs
of theta functions to theta functions) and is $\hat K$-finite (\lref{GlobalTheory.Core.Rk.IP.delta_isKFinVec}: $\delta$ of Fock vectors is
$\hat K$-finite). Since $t>0$, $r>0$, $\dim A=r+t$, $r+2<2t$, $r<t$ and the places above $2$ split, (RV) gives the residue $E_1$ of
$E(\cdot,s;\varphi^{\mathrm{SW}})$ at $s_0(A,B)$ with $Z(E_1;\bar u,\bar u)=0$.

\emph{Step 4: the residue of the zeta integral.} Let $\mathcal Z(s)=Z(E(\cdot,s;\varphi^{\mathrm{SW}});\bar u,\bar u)$. By \cref{pf:liftinj:zetaint},
$E(\cdot,s;\varphi^{\mathrm{SW}})$ has the regularity required in \ref{pf:liftinj:D4}, and $\mathcal Z(s)=\int_{G_A}I_s$ for $\Re s>\dim A/2$. By
\ref{pf:liftinj:D6} and the continuity of $u$, $\bar u$ is of rapid decay (\lref{GlobalTheory.Unf.contAut_rapidDec_bounded_integrable}:
continuous automorphic functions on a compact quotient are bounded and of rapid decay), and $\chi_W^{-1}$ is a Hecke character. By
\ref{pf:liftinj:D4}, $\mathcal Z$ has residue $Z(E_1;\bar u,\bar u)=0$ at $s_0(A,B)$.

\emph{Step 5: the Euler product.} By \cref{pf:liftinj:levelconst} (with $f$ from Step 2) there is $C\ne0$ with $\int_{G_A^S}I_s=C$ for all $s$. The form $u$ is spherical off $S$ (it has
Satake parameters there), the places above $2$ outside $S$ split, $\chi_W$ is unramified off $S$, the frame lattice of $A$ is self-dual
(\lref{GlobalTheory.WeilAtPlace.isSelfDualC_frameLat}) and $\Phi^+$ is spherical off $S$, so $\Phi^+$ is unramified off $S$. For $\Re s>\dim A+1$,
\cref{pf:liftinj:zetaint} and (ER) give
\[
\mathcal Z(s)=\int_{G_A}I_s=\frac{L^S(s+\frac12,\mathrm{par}\times\chi_W)}{d^S_{\dim A,\dim B}(s)}\int_{G_A^S}I_s
=\frac{L^S(s+\frac12,\mathrm{par}\times\chi_W)}{d^S_{\dim A,\dim B}(s)}\,C .
\]

\emph{Step 6: contradiction.} Since $s_0(A,B)=(\dim B-\dim A)/2=t/2$, Step 4 says that $\mathcal Z$ has residue $0$ at $t/2$, whereas
\cref{pf:liftinj:resne} (with the segment hypothesis of Step 1 and Step 5) says that this residue is non-zero. Hence $u$ has a lift along $\omega$.
\end{proof}

% ---------------------------------------------------------------------------
\subsection{Real data and the proof of \cref{prop:families:liftinj}}\label{sec:pf:liftinj:realdata}

\subsubsection{Real data}
The statement of \cref{prop:families:liftinj} concerns the \emph{real data} of the development. Fix a global datum $D$ with ranks
$r=\dim W_i$ and $t$ (\cref{sec:families:theta}), a finite group $G$, index types $\mathrm{Pl}_0$ and $T$ (the latter indexes the
\emph{tuples} $\mu$, with a base tuple $\mu_0$), a family input $\mathcal A$, a number field $K_s$ with $K_s\to\mathit{LL}$, a real datum $R$ over
$(D,\mathcal A,K_s)$ and a face $F$ of $(D,\mathcal A)$. For $0\le i<m$ let $V_i,W_i$ be the hermitian spaces of $D$ attached to $i$. The components of $R$
which enter are (for $0\le i<m$ and tuples $\mu$):
\begin{itemize}
\item finite-dimensional Hilbert spaces $H_V(i,\mu)$ of vector-valued functions on the finite adelic points of $\Uni(V_i)$, with values in a
coefficient space $W_V(i)$, and extensions $y\mapsto\mathrm{ext}_V(y):G_{V_i}\to W_V(i)$;
\item a Haar measure $\nu_V(i,\mu)$ on $G_{V_i}$ with a fundamental domain $\mathfrak F_V(i,\mu)$ for $\Gamma_{V_i}$, a frame $\beta_V(i)$, and finite sets
$P_g$ (level primes) and $Q(\mu)$ of places of $E$;
\item a standard additive character $\psi$, a family $\theta$ of archimedean embeddings, and Weil data $\omega_P(i,\mu)$ for $(V_i,W_i)$ with characters
$(\chi_W(i,\mu),\chi_V(i,\mu))$, where $\chi_W(i,\mu),\chi_V(i,\mu)$ are the two Hecke characters attached to $(\mu,i+1)$ by the theta data of $R$; we
write $\chi_V(w)=\chi_V(i,\mu)(w)$ and $\chi_W(w)$ for their values at a place $w$;
\item a family of Fock vectors $\mathrm{pt}(i,\mu,P)$ on $\liftinjX(V_i,W_i)$ attached to polynomials $P$;
\item a space $\mathrm{MMap}(i)$ of maps from the discrete series models of $G_{W_i,\infty}$ to functions on $G_{W_i}$, a finite-dimensional space
$H_W(i)$ with an embedding $e_W:H_W(i)\to\mathrm{MMap}(i)$, and two linear \emph{theta maps} $\theta^F_{i,\mu}:H_V(i,\mu)\to\mathrm{MMap}(i)$ and
$\theta^G_{i,\mu}:H_V(i,\mu)\to H_W(i)$.
\end{itemize}

A place $w$ of $E$ is \emph{good} for $\mu$ if $w$ is distinct from, and not conjugate to, every place of $P_g\cup Q(\mu)$; equivalently, $w$ does not lie
above the finite set $S(\mu)$ of places of $F$ below $P_g\cup Q(\mu)$.

\begin{definition}[Hecke eigenvectors, segments, lifts and the three properties]\label{pf:liftinj:realdefs}
Fix $i$ and a tuple $\mu$.
\begin{enumerate}[label=(\alph*)]
\item A \emph{Hecke family} is a function $h$ assigning to every good $w$ and every finitely supported function $c$ on labels $d\in\mathrm{List}(\Z)$
an endomorphism $h(w,c)$ of $H_V(i,\mu)$ compatible with the classical Hecke operators (sums of right translates over double coset
representatives). A \emph{system} is a function $s$ on pairs (place of $E$, label); $y\in H_V(i,\mu)$ is an \emph{eigenvector with system $s$} if
$h(w,\delta_d)y=s(w,d)y$ for all good $w$ and all Hecke labels $d$ at $w$. A family $B$ is a \emph{Satake family of rank $n$ for $s$} off a set $S$ of
places of $F$ if at the places off $S$ the values of $s$ are the Satake transforms of the multisets $B(w)$ of rank $n$.
\item A system $s$ \emph{has a segment} if there is a Satake family $B$ of rank $\dim V_i$ for $s$ off $S(\mu)$ such that for every good $w$,
$B(w)=B'(w)+\liftinjstr(\Nm w,\chi_V(w),t)$ for a multiset $B'(w)$ all of whose elements have absolute value at most $\Nm w^{(r-1)/2}$
(\cref{def:families:string}).
\item For $y\in H_V(i,\mu)$, a theta function $f$ on $\liftinjX(V_i,W_i)$, a linear form $\ell$ on $W_V(i)$ and $g\in G_{W_i}$, the \emph{downward theta
integral} is
\begin{equation}\label{eq:pf:liftinj:downint}
\mathrm{down}(y,f,\ell,g)=\int_{\mathfrak F_V(i,\mu)}\theta_{\omega_P(i,\mu)}(h,g;f)\,\ell\bigl(\mathrm{ext}_V(y)(h)\bigr)\,d\nu_V(i,\mu)(h),
\end{equation}
and $y$ \emph{has some lift} if $\mathrm{down}(y,f,\ell,g)\ne0$ for some $(f,\ell,g)$.
\item $y$ is \emph{good} if for every $\ell$ the function $h\mapsto\ell(\mathrm{ext}_V(y)(h))$ is fixed by the group algebra operator of the
idempotent $e_C$ of the face (\cref{prop:families:liftinj}(b)).
\item $\liftinjP{NV}(i,\mu)$: for every Hecke family $h$ and every system $s$ with a segment, every non-zero eigenvector with system $s$ has some lift.
\item $\liftinjP{TNZ}(i,\mu)$: for every Hecke family $h$ and every system $s$ with a segment, every good eigenvector $y$ with system $s$ and
$\theta^G_{i,\mu}(y)=0$ is zero.
\item $\liftinjP{INJ}(R)$: for every $i$, every Hecke family $h$ at $\mu_0$, every system $s$ and all eigenvectors $x,y\in H_V(i,\mu_0)$ with system $s$:
if $\theta^G_{i,\mu_0}(x)\ne0$, $y$ is good and $\theta^G_{i,\mu_0}(y)=0$, then $y=0$.
\end{enumerate}
The \emph{standing hypotheses} (G) on $(R,F)$ are: the core conditions of $R$ (level, group, family, models, chain and theta-compatibility
conditions), the existence of size data with a parameter choice, absolute irreducibility of the characters in the family input, good types at the
primes $g_0,g_1$, the guard condition on the sets $Q(\mu)$, and the gap condition $\lambda_n-\lambda_{n+1}\ge100r^4$ on the Harish-Chandra
parameters. The \emph{fullness} condition (F) on $R$ says that each $\omega_P(i,\mu)$ is a full Weil datum for $(V_i,W_i)$ with the characters above.
\formal{Formalism.IndefR.Links.RealDat}
\formal{GlobalTheory.Hk3.IsOpsV}
\formal{GlobalTheory.DownLift.IsEigV}
\formal{GlobalTheory.LFun.IsSysParamD}
\formal{GlobalTheory.DownLift.HasString}
\formal{GlobalTheory.DownLift.downInt}
\formal{GlobalTheory.DownLift.SomeLift}
\formal{GlobalTheory.DownLift.GoodVt}
\formal{GlobalTheory.DownLift.NonVan}
\formal{GlobalTheory.DownLift.TypeNZ}
\formal{GlobalTheory.LiftRows.LiftInj}
\formal{GlobalTheory.LiftRows.Given}
\formal{Formalism.IndefR.Links.RealDat.FullP}
\end{definition}

\subsubsection{Non-vanishing on segments}

\begin{proposition}[Non-vanishing on segments]\label{pf:liftinj:nonvan}
If (DN) holds for $D$, then $\liftinjP{NV}(i,\mu)$ holds for every real datum $R$ over $D$ satisfying (G) and (F), every $i$ and every tuple $\mu$.
\formallabel{GlobalTheory.T.thetaLift_nonvanishing_onString_of_definiteDown}
\formallabel{GlobalTheory.Core.Rk.NonVanAt.nonVan_of_downNe}
\end{proposition}

\begin{proof}
Let $h$ be a Hecke family, $s$ a system with a segment and $y\ne0$ an eigenvector with system $s$. We use five facts about real data, proved
from (G) and (F) in this section of the development:
\begin{enumerate}[label=(N\arabic*)]
\item\label{pf:liftinj:N1} (\lref{GlobalTheory.T.thetaLiftPair_definite_place_character_data}) $V_i$ carries definite data
$(\theta,c_A,a)$; $(\beta_V(i),P_g\cup Q(\mu))$ are level places for $V_i$, and the places of $F$ below $P_g\cup Q(\mu)$ form $S(\mu)$;
$\chi_V(i,\mu)$ is unramified off $S(\mu)$; $\tilde\eta$ exists; $W_i$ has consistent signature data, $\dim W_i=r>0$, $\dim V_i=r+t$ and $r+5\le t$.
\item\label{pf:liftinj:N2} (\lref{GlobalTheory.T.realIntegral_thetaLift_definiteForm_ne_zero}) For every Haar measure $m$ on $G_{V_i}$
there are a linear form $\ell$ on $W_V(i)$ and a polynomial $P$ such that, with $u_y=\ell\circ\mathrm{ext}_V(y)$ and $\Phi_P=\mathrm{pt}(i,\mu,P)$,
the archimedean integral \eqref{eq:pf:liftinj:realint} for $(\Phi_P,u_y)$ is non-zero for every $\nu$ of covolume one and every Haar measure on
$G_{V_i,\infty}$.
\item\label{pf:liftinj:N3} (\lref{GlobalTheory.T.linearForm_extV_mem_A0_and_prin_invariant}) For every $\ell$, $u_y\in\mathcal A_0(V_i)$ and
$u_y$ is right invariant under $\mathrm{prin}(V_i,\beta_V(i),P_g\cup Q(\mu))$.
\item\label{pf:liftinj:N4} (\lref{GlobalTheory.T.satake_string_gives_simple_pole_hypotheses}) If $s$ has a segment, there are a Satake
family $B$ of rank $\dim V_i$ for $s$ off $S(\mu)$ and a family $c$ with $\mathbf H(S(\mu),r,t;c,B;\chi_V,\chi_V)$; this rewrites the segment
condition, using that the good places are the places off $S(\mu)$, that $B'(w)$ has $\dim V_i-t=r$ elements, and that $|\chi_V(w)|=1$.
\item\label{pf:liftinj:N5} (\lref{GlobalTheory.T.scalarComponent_parameters_contain_string} and
\lref{GlobalTheory.T.thetaLiftFockVectors_spherical_principalFixed}) If $u_y\ne0$, then $u_y$ has Satake parameters $B$ off $S(\mu)$ for every
Satake family $B$ of $s$; and each $\Phi_P$ is a Fock vector for $(1,s_W)$ with $s_W$ the Cartan involution of the discrete series model of $W_i$,
spherical off $S(\mu)$ and fixed by the principal congruence subgroup.
\end{enumerate}
Let $m=\nu_V(i,\mu)$, and choose $\ell,P$ by \ref{pf:liftinj:N2}. Then $u_y\ne0$: otherwise $r(u_y,u_y)=0$ and the integral
\eqref{eq:pf:liftinj:realint} would vanish (for some choice of $\nu$ and of a Haar measure on $G_{V_i,\infty}$, which exist). Take $B,c$ from
\ref{pf:liftinj:N4}. By \ref{pf:liftinj:N1}, \ref{pf:liftinj:N3}, \ref{pf:liftinj:N4} and \ref{pf:liftinj:N5}, $u_y$ is a form of segment type
$(r,t)$ for $(V_i,\beta_V(i),P_g\cup Q(\mu))$ with parameters $B$ and character $\chi_V(i,\mu)$, and by \ref{pf:liftinj:N2}, \ref{pf:liftinj:N5} the
theta function $\Phi_P$ occurs for $u_y$ along $\omega_P(i,\mu)$, whose characters are $(\chi_W(i,\mu),\chi_V(i,\mu))$. By (DN) (with $A=V_i$,
$W=W_i$, and $\chi_V(i,\mu)$ in the role of $\chi_W$) there are a theta function $f$ and $g\in G_{W_i}$ with $\liftinjup_{\omega_P}(f,u_y)(g)\ne0$.
Since $u_y$ is left $\Gamma_{V_i}$-invariant and $\mathfrak F_V(i,\mu)$ is a fundamental domain for the Haar measure $m$, the integral over $[G_{V_i}]$
in \eqref{eq:pf:liftinj:up} is the integral over $\mathfrak F_V(i,\mu)$, which is $\mathrm{down}(y,f,\ell,g)$. Hence $y$ has some lift.
\end{proof}

\subsubsection{The main theorem}

Besides $\liftinjP{NV}$, the proof uses the following properties of $(i,\mu)$, all proved in the development from (G) and (F) and, where indicated,
inputs of \cref{pf:liftinj:lit}.
\begin{itemize}
\item \emph{Eigen-compatibility} $\liftinjP{EC}$: for every Hecke family and system $s$ there is a system $c$ on the $W$-side such that $\theta^F(y)$
lies in the $c$-part $\mathrm{part}_W(c)\subseteq\mathrm{MMap}(i)$ for every eigenvector $y$ with system $s$. From (e)
(\lref{GlobalTheory.DownEigen.downEig}).
\item \emph{Generation at finite places} $\liftinjP{Gen}_f$: if $y$ has no component along the level characters at $P_g$ and
$\mathrm{down}(y,f,\ell,g)=0$ for all $f$ of \emph{finite type} (theta functions killed by every finite-adelic Heisenberg combination that kills the
distinguished vector of $\omega_P(i,\mu)$) and all $\ell,g$, then $y$ has no lift. From (e) (\lref{GlobalTheory.DownLift.GenFinite.genFin}).
\item \emph{Generation at real places} $\liftinjP{Gen}_\infty$: if $\theta^F(y)=0$, then $\mathrm{down}(y,f,\ell,g)=0$ for every $f$ of finite type
(\lref{GlobalTheory.GenReal.genInf}).
\item $\liftinjP{GM}$: $e_W(\theta^G(y))=\theta^F(y)$ for good $y$ (\lref{GlobalTheory.GoodLift.goodMem_of_given}).
\item $\liftinjP{GC}$: good vectors have no component along the level characters at $P_g$ (\lref{GlobalTheory.GoodTypeNoChi.goodNoChi}).
\item \emph{Existence of parameters} $\liftinjP{PE}$: systems of non-zero eigenvectors have Satake families of rank $\dim V_i$, and every $c$ with
$\mathrm{part}_W(c)\ne0$ has a Satake family $B'$ of rank $r$ with $|b|\le\Nm w^{(r-1)/2}$ for $b\in B'(w)$, $w$ good. From (a)
(\lref{GlobalTheory.DownLift.Params.paramsEx}).
\item The \emph{unramified correspondence} $\liftinjP{UC}(\mathrm{ok})$, for a predicate $\mathrm{ok}$ on places: if $y$ is an eigenvector with system $s$,
$0\ne\theta^F(y)\in\mathrm{part}_W(c)$, $B$ is a Satake family of rank $\dim V_i$ for $s$ and $B'$ one of rank $r$ for $c$, then for every good $w$
with $\mathrm{ok}(w)$,
\begin{equation}\label{eq:pf:liftinj:unr}
B(w)=\chi_V(w)\chi_W(w)^{-1}B'(w)+\liftinjstr\bigl(\Nm w,\chi_V(w),t\bigr).
\end{equation}
At self-conjugate (inert) places this follows from (e) (\lref{GlobalTheory.DownInert.unrCorr_inert}) and at the other places from the
correspondence for $\GL_n$ (\lref{GlobalTheory.SplitCorr.unrCorr_split}).
\end{itemize}
The derivation of $\liftinjP{TNZ}$ and $\liftinjP{INJ}$ from these properties is carried out in the development by two short lemmas; we recall their
arguments.
\begin{enumerate}[label=(T\arabic*)]
\item\label{pf:liftinj:T1} (\lref{GlobalTheory.DownLift.typeNZ_of_nonVan_gen}) $\liftinjP{NV}$, $\liftinjP{Gen}_\infty$, $\liftinjP{Gen}_f$,
$\liftinjP{GC}$ and $\liftinjP{GM}$ imply $\liftinjP{TNZ}$: let $y$ be a good eigenvector with a segment and $\theta^G(y)=0$, and suppose $y\ne0$. Then
$\theta^F(y)=e_W(\theta^G(y))=0$, so all downward integrals of $y$ against functions of finite type vanish; $y$ has no component along the level
characters, so $y$ has no lift by $\liftinjP{Gen}_f$, contradicting $\liftinjP{NV}$.
\item\label{pf:liftinj:T2} (\lref{GlobalTheory.LiftInjProof.liftInj_of_typeNZ}) If $\liftinjP{TNZ}$, $\liftinjP{EC}$, $\liftinjP{UC}$ for the
predicates ``$w$ is self-conjugate'' and ``$w$ is not self-conjugate'' and $\liftinjP{PE}$ hold at $(i,\mu_0)$ for every $i$, then $\liftinjP{INJ}(R)$:
given $x,y$ as in \cref{pf:liftinj:realdefs}(g), the core condition gives $e_W\circ\theta^G=\theta^F$ with $e_W$ injective, so $\theta^F(x)\ne0$; by
$\liftinjP{EC}$ and $\liftinjP{PE}$ there are $c$ and Satake families $B$, $B'$ with $\theta^F(x)\in\mathrm{part}_W(c)$, and \eqref{eq:pf:liftinj:unr}
at every good $w$ (using the two predicates according to whether $w$ is self-conjugate) shows that $s$ has a segment, since
$|\chi_V(w)|=|\chi_W(w)|=1$ (\lref{GlobalTheory.LiftInjProof.hasString_of_lift}); $\liftinjP{TNZ}$ applied to $y$ gives $y=0$.
\end{enumerate}

\begin{lemma}[The unramified correspondence at all good places]\label{pf:liftinj:ucall}
Under inputs (a) and (e) of \cref{pf:liftinj:lit}, (G) and (F), $\liftinjP{UC}(\mathrm{ok})$ holds for every $i$, every tuple $\mu$ and every predicate
$\mathrm{ok}$.
\formallabel{GlobalTheory.T.unramifiedThetaCorrespondence_all_goodPrimes}
\formallabel{GlobalTheory.UnrAll.unrCorr_all}
\end{lemma}

\begin{proof}
For $\mathrm{ok}\equiv\mathrm{true}$: given a good place $w$, apply the self-conjugate case or the other case of the unramified correspondence cited
above, according to whether $w$ is self-conjugate. Since the conclusion of $\liftinjP{UC}(\mathrm{ok})$ only concerns places $w$ with $\mathrm{ok}(w)$,
the case $\mathrm{ok}\equiv\mathrm{true}$ implies all others.
\end{proof}

\begin{theorem}[\cref{prop:families:liftinj}]\label{pf:liftinj:main}
Assume the inputs (a)--(e) of \cref{pf:liftinj:lit}. Let $D$ be a global datum, $R$ a real datum over $D$ and $F$ a face such that (G) and (F) hold.
Then $\liftinjP{INJ}(R)$ holds, and $\liftinjP{NV}(i,\mu)$ and $\liftinjP{TNZ}(i,\mu)$ hold for all $i$ and all tuples $\mu$.
\formallabel{GlobalTheory.T.liftInj_and_nonVan_typeNZ_of_given}
\formallabel{GlobalTheory.Core.Rk.LiftOnto.liftInj_nonVan_typeNZ_of_hypotheses_alt}
\end{theorem}

\begin{proof}
By \cref{pf:liftinj:euler} (inputs (a), (c)) (ER) holds for $D$; by \cref{pf:liftinj:pole} (input (d)) (SP) holds; by \cref{pf:liftinj:downne} (input
(b)) (DN) holds for $D$; by \cref{pf:liftinj:nonvan}, $\liftinjP{NV}(i,\mu)$ holds for all $i,\mu$. The properties $\liftinjP{EC}$,
$\liftinjP{Gen}_f$ (input (e)), $\liftinjP{Gen}_\infty$, $\liftinjP{GM}$, $\liftinjP{GC}$ and $\liftinjP{PE}$ (input (a)) hold for all $i,\mu$ by the
results cited above, and $\liftinjP{UC}(\mathrm{ok})$ holds for all predicates by \cref{pf:liftinj:ucall}. Now \ref{pf:liftinj:T1} gives
$\liftinjP{TNZ}(i,\mu)$ for all $i,\mu$, and \ref{pf:liftinj:T2} (with $\mu=\mu_0$) gives $\liftinjP{INJ}(R)$.
\end{proof}

\begin{corollary}\label{pf:liftinj:nvcor}
Under inputs (a)--(e), $\liftinjP{NV}(i,\mu)$ holds for every real datum satisfying (G) and (F), every $i$ and every tuple $\mu$.
\formallabel{GlobalTheory.T.nonGalois_downLift_nonvanishing}
\formallabel{GlobalTheory.DownLift.nonvanishing_of_down_lift}
\end{corollary}

\begin{proof}
This is part of \cref{pf:liftinj:main}.
\end{proof}

\begin{remark}\label{pf:liftinj:mainrem}
(a) In \cref{prop:families:liftinj}, statement (a) is $\liftinjP{NV}$, (b) is $\liftinjP{TNZ}$ and (c) is $\liftinjP{INJ}$. The splitting characters of
$\omega_P(i,\mu)$ are $(\chi_W(i,\mu),\chi_V(i,\mu))$ in the order (definite side, other side), and the segment in the parameters is twisted by
$\chi_V(i,\mu)$, the character attached to the $W_i$-side; with this convention (DN) is applied with $\chi_V(i,\mu)$ in the role of $\chi_W$.
(b) The inequality $r+5\le t$ of the segment hypothesis is supplied by the ranks of the datum (\ref{pf:liftinj:N1}). It is used only through its
consequence $r<t$: \cref{pf:liftinj:pole,pf:liftinj:resne} use $r<t$, and \cref{pf:liftinj:downne} uses $\dim A\le\dim H$, $r<t$ and $r+2<2t$,
which follow from $r<t$ and $r\ge1$.
\end{remark}

% ---------------------------------------------------------------------------
\subsection{Type non-vanishing for the real data at the real places}\label{sec:pf:liftinj:B}

The construction attaches to every global datum $D$, character frame $\mathit{cf}$, norm choice $\mathit{nch}$ and tied tuple datum $\mathit{tp}$ a
type of \emph{real-place data}. An element $X$ is a large record; the components used here are a real datum $X.R$ over $D$ (with its finite group,
index types, family input and number field), a face $X.F$, and proofs of (G) for $(X.R,X.F)$ and of (F) for $X.R$. The other components (row and
lift identifications, column, division and eigen data, recursion data, the tie with $\mathit{cf}$ and $\mathit{tp}$) play no role in this subsection.

The \emph{geometric compatibility package} $\mathcal X$ is a conjunction of twelve propositions of the development: existence of Weil
representations for unitary pairs, existence and classification of discrete series, Harish-Chandra analyticity, trace class of cusp forms,
Dixmier--Malliavin and Flath factorisation, the Weil--Siegel formula for unitary groups in the convergent range, inputs (b), (c), (d) of
\cref{pf:liftinj:lit}, and a height-counting statement. Only (b), (c), (d) and the Weil--Siegel formula are used below.

\begin{definition}[Saturation and the Rallis statement]\label{pf:liftinj:satdef}
\begin{enumerate}[label=(\alph*)]
\item \emph{Saturation} $\liftinjP{Sat}(i,\mu)$: for every Hecke family $h$ and every system $s$ with a segment there is a system $c$ such that every good
eigenvector $y$ with system $s$ is of the form $(\theta^G)^*(x)$ for some $x\in H_W(i)$ with $e_W(x)\in\mathrm{part}_W(c)$ lying in the image of the
idempotent $e_H$ of $H_W(i)$ attached to the face; here $(\theta^G)^*$ is the adjoint of $\theta^G_{i,\mu}$ for the Hilbert space structures.
\item The \emph{Rallis statement} $\liftinjP{RI}(D)$ asserts, for a positive definite $V$ with consistent signature data, a hermitian space $W$ with
$2\dim W<\dim V$, a Weil datum for $(V,W)$, theta functions $f_1,f_2$ and smooth cusp forms $\varphi_1,\varphi_2$ on $G_W$, the convergence of the
theta lifts $\Theta(f_j,\varphi_j)$ of $\varphi_j$ to $G_V$ and the Rallis inner product formula
\[
\langle\Theta(f_1,\varphi_1),\Theta(f_2,\varphi_2)\rangle_{[G_V]}=\vol([G_V])\int_{G_W}\langle\omega(1,g)f_2,f_1\rangle\,\langle R(g)\varphi_2,\varphi_1\rangle\,dg
\]
with absolute convergence.
\end{enumerate}
\formal{GlobalTheory.DownLift.Saturated}
\formal{GlobalTheory.RawLift.RallisConv}
\end{definition}

\begin{lemma}[Adjoint surjectivity]\label{pf:liftinj:adj}
Let $f:X\to Y$ be a linear map between finite-dimensional Hilbert spaces and $N\subseteq X$ a subspace with $f^*f(N)\subseteq N$ such that $f$ is
injective on $N$. Then every $y\in N$ is of the form $f^*f(y')$ with $y'\in N$.
\formallabel{GlobalTheory.DownLift.Sat.exists_adjoint_comp}
\end{lemma}

\begin{proof}
If $u\in N$ and $f^*f(u)=0$, then $\|f(u)\|^2=\langle f^*f(u),u\rangle=0$, so $u=0$. Hence $f^*f|_N:N\to N$ is injective, hence bijective.
\end{proof}

\begin{lemma}[Saturation]\label{pf:liftinj:sat}
Assume inputs (a) and (e), (G) and (F). If $\liftinjP{TNZ}(i,\mu)$, $\liftinjP{EC}(i,\mu)$ and $\liftinjP{GM}(i,\mu)$ hold, then $\liftinjP{Sat}(i,\mu)$
holds.
\formallabel{GlobalTheory.DownLift.Sat.saturated}
\end{lemma}

\begin{proof}
Let $h$ be a Hecke family, $s$ a system with a segment, and let $N\subseteq H_V(i,\mu)$ be the subspace of good eigenvectors with system $s$. Write
$\Theta=\theta^G_{i,\mu}$. By $\liftinjP{EC}$ there is $c$ with $\theta^F(y)\in\mathrm{part}_W(c)$ for every eigenvector $y$ with system $s$. Two facts proved
in the development show that $\Theta^*\Theta$ maps $N$ to $N$: $\Theta^*\Theta$ maps eigenvectors with system $s$ to eigenvectors with system $s$
(\lref{GlobalTheory.DownLift.Sat.eig_keep}, from input (e) at self-conjugate places and the split partners elsewhere), and $\Theta^*(x)$ is good
whenever $e_Hx=x$ (\lref{GlobalTheory.DownLift.Sat.good_adjoint}); and $e_H\Theta(y)=\Theta(y)$ for all $y$
(\lref{GlobalTheory.Hk3.eH_laws}: the laws of the idempotent $e_H$). By $\liftinjP{TNZ}$, $\Theta$ is injective on $N$. By \cref{pf:liftinj:adj}
every $y\in N$ is $\Theta^*(x)$ with $x=\Theta(y')$, $y'\in N$; then $e_W(x)=\theta^F(y')\in\mathrm{part}_W(c)$ by $\liftinjP{GM}$ and the choice of $c$,
and $x$ lies in the image of $e_H$.
\end{proof}

\begin{proposition}[Lift rows from non-vanishing]\label{pf:liftinj:rallis}
Assume inputs (a) and (e) of \cref{pf:liftinj:lit}. Let $R$ be a real datum satisfying (G) and (F), assume $\liftinjP{RI}(D)$ and $\liftinjP{NV}(i,\mu)$ for
all $i,\mu$. Then $\liftinjP{INJ}(R)$ holds, and for all $i,\mu$: $\liftinjP{TNZ}(i,\mu)$, $\liftinjP{GM}(i,\mu)$, $\liftinjP{EC}(i,\mu)$,
$\liftinjP{UC}(\mathrm{true})$ and $\liftinjP{Sat}(i,\mu)$.
\formallabel{GlobalTheory.T.liftInj_and_properties_of_rallis_of_nonVan}
\formallabel{GlobalTheory.DownLift.liftRows_injective_and_typeProperties_of_nonvanishing}
\end{proposition}

\begin{proof}
The conclusions are obtained as in \cref{pf:liftinj:main}: $\liftinjP{EC}$, $\liftinjP{Gen}_f$, $\liftinjP{Gen}_\infty$, $\liftinjP{GM}$,
$\liftinjP{GC}$ and $\liftinjP{PE}$ hold by the results cited in \cref{sec:pf:liftinj:realdata}; $\liftinjP{UC}$ holds for the predicate ``$w$ is
self-conjugate'' and for its negation, hence for $\mathrm{true}$ (\cref{pf:liftinj:ucall}); \ref{pf:liftinj:T1} gives $\liftinjP{TNZ}$ from $\liftinjP{NV}$,
and \ref{pf:liftinj:T2} gives $\liftinjP{INJ}(R)$. Finally \cref{pf:liftinj:sat} gives $\liftinjP{Sat}(i,\mu)$. The hypothesis $\liftinjP{RI}(D)$ is not
used (\cref{pf:liftinj:rallisrem}).
\end{proof}

\begin{remark}\label{pf:liftinj:rallisrem}
In the formal proof, $\liftinjP{RI}(D)$, input (a) and the laws of $R$ are first used to derive a statement on the theta maps of $R$ and, for all $i,\mu$,
the injectivity of the adjoints of all raw theta maps, together with a positivity and a spanning statement (a derivation of about $10\,000$ lines). These
are passed only to the saturation lemma, whose formal statement takes three further hypotheses (the injectivity of the adjoints of raw theta maps, the
unramified correspondence and the existence of parameters) which it does not use; the proof of \cref{pf:liftinj:sat} above does not use them either.
Hence neither $\liftinjP{RI}(D)$ nor that derivation affects any conclusion of \cref{pf:liftinj:rallis}, and we give no proof of the derivation. The same
holds for the derivation of $\liftinjP{RI}(D)$ from input (a) and the Siegel--Weil formula contained in $\mathcal X$ (about $4\,000$ lines), which the
formal proof of \cref{pf:liftinj:real} performs in order to apply \cref{pf:liftinj:rallis}: we apply instead the proof of \cref{pf:liftinj:rallis}, which
does not use $\liftinjP{RI}(D)$.
\end{remark}

\begin{theorem}[Type non-vanishing at the real places]\label{pf:liftinj:real}
Assume inputs (a) and (e) of \cref{pf:liftinj:lit} and the geometric compatibility package $\mathcal X$. Then for every global datum $D$, all
$\mathit{cf},\mathit{nch},\mathit{tp}$, all real-place data $X$ over them, every $j$ and every tuple $\mu$, $\liftinjP{TNZ}(j,\mu)$ holds for the real
datum $X.R$.
\formallabel{GlobalTheory.T.galoisRep_descends_at_real_places}
\formallabel{GlobalTheory.DownLift.AtReal.lift_splits_locally}
\end{theorem}

\begin{proof}
Inputs (b), (c), (d) contained in $\mathcal X$, together with (a) and (e), are the hypotheses of \cref{pf:liftinj:nvcor}, which gives
$\liftinjP{NV}(i,\mu)$ for $X.R$ and all $i,\mu$. The proof of \cref{pf:liftinj:rallis}, which does not use $\liftinjP{RI}(D)$ (\cref{pf:liftinj:rallisrem}),
applied to $X.R$, $X.F$ and the proofs of (G) and (F) contained in $X$, gives $\liftinjP{TNZ}(j,\mu)$.
\end{proof}

\begin{remark}\label{pf:liftinj:realrem}
The conclusion is the case of \cref{pf:liftinj:main} (statement $\liftinjP{TNZ}$) for the real datum contained in $X$, and by \cref{pf:liftinj:rallisrem} the
route through $\liftinjP{RI}$ is not needed for it. The formal name refers to the role of the statement in the construction, not to its content.
\end{remark}

% ---------------------------------------------------------------------------
\subsection{The third-layer packages}\label{sec:pf:liftinj:C}

This subsection proves an assembly statement which regroups literature statements into the packages used at the third layer of the construction.
The objects involved are propositions of the development; we introduce short names for them and say where each component comes from.

\begin{definition}[The packages]\label{pf:liftinj:packdef}
\begin{enumerate}[label=(\alph*)]
\item $\mathcal R$ is the statement that for every number field $E$, algebraic closure $\Omega$, valuation subring $A$ of $\Omega$ and place $w$ of $E$
below $A$, a local reciprocity datum at $w$ exists (the local Langlands correspondence for $\GL_n$ in the form used by the development).
\item The \emph{non-endoscopic hypotheses} $\mathcal N$ have five components: a local representation theory package $\mathcal N_{\mathrm{loc}}$ (one of whose
fields, which we denote $h_{\mathrm{loc}}$, is a proof of $\mathcal R$; the others include the local converse theorem, depth-zero induction, Bernstein--Zelevinsky
theory, the geometric lemma, Zelevinsky's classification, Green's cuspidal characters and the generic unitary dual of $\GL_n$); a global automorphic
and Galois package $\mathcal N_{\mathrm{aut}}$ (global Artin reciprocity, rapid decay, Harish-Chandra finiteness, ray class fields, base change for
definite unitary groups, Galois representations of polarized cuspidal representations, purity, \dots); a Galois cohomology package
$\mathcal N_{\mathrm{coh}}$ (Poitou--Tate duality, local finiteness, local Tate duality and Euler characteristics, global finiteness and Euler
characteristic, unramified orthogonality, Greenberg--Wiles, local Bloch--Kato statements); a join package $\mathcal N_{\mathrm{join}}$ (containing a
toric denominator statement and a toric transport statement); and the functional equation of Hecke--Tate $L$-functions.
\item The \emph{third-layer literature package} $\Pi$ has fourteen fields: the theta-local inputs, rapid decay, Harish-Chandra finiteness, Green's
cuspidal characters (twice), Zelevinsky's polynomial ring, two ray class field statements, ray characters from Artin reciprocity, base change for
definite unitary groups, the generic unitary dual and the standard cosocle of $\GL_n$, Chebotarev's theorem and Howe's unramified generation.
\item The \emph{series package} $\mathcal S$ has six fields: rationality and integrality of twist series, the three Hecke statements of
\cref{pf:liftinj:hecke}, and a conjugation statement for toric ratios.
\item The \emph{elliptic package} $\mathcal E$ has twenty-five fields: Chebotarev statements, Poitou--Tate for totally complex fields, cohomological
dimension two, tame local theory, local finiteness, local Tate duality, local Euler characteristic, global finiteness and Euler characteristic,
unramified orthogonality, archimedean signature matching, Greenberg--Wiles, ray characters, Chebotarev for degree one primes, toric denominators,
the five series and Hecke fields of $\mathcal S$, finiteness of automorphic forms on $\GL_n$, two statements on Galois representations attached to
cuspidal representations, de Rham $\Rightarrow$ potentially semistable, and the local Bloch--Kato conclusion at all primes.
\item For $h\in\mathcal R$, the \emph{local Langlands package} $\mathcal L(h)$ has six fields: $\Pi$, $\mathcal S$, an archimedean comparison, a
parameter statement, a point statement, and a cuspidal parameter statement at members.
\end{enumerate}
\formal{UpperTheory.Zero.LLC.RecTExists}
\formal{GlobalTheory.Mem.LiteralAutomorphy.NonEndoscopicHypothesesAlt}
\formal{GlobalTheory.HeightCounting.Up.PublishedInputsThirdLayerAlt}
\formal{GlobalTheory.HeightCounting.Up.SerN}
\formal{GlobalTheory.Mem.Top.Bl.EllipticHypothesis}
\formal{GlobalTheory.Mem.Top.Bl.LocalLanglandsHypothesis}
\end{definition}

The bulk of the proof (about $22\,500$ lines) is the following result on Hecke operators.

\begin{proposition}[Hecke integrality, transform and reading]\label{pf:liftinj:hecke}
The following three statements hold:
\begin{enumerate}[label=(\alph*)]
\item (Hecke integrality) for every CM field $E$, host, datum, level and moduli space of the development and every coefficient datum, the frame
and value forms of the Hecke operators are integral, both for the good classes and for all classes;
\item (Hecke transform) for primes split off the level, the Hecke operators of the development transform sections of the moduli spaces, and keep
cusp sections of the cores;
\item (Hecke reading) for real data tied to a core, the frame pin at the distinguished embedding implies that the theta values read the
Hecke operators at split primes: the function on $G_{W_i}$ attached to the image of a section under a geometric Hecke correspondence equals a
scalar times a sum of right translates of the function attached to the section.
\end{enumerate}
\formallabel{GlobalTheory.HeckeMember.heckeInt_transform_readTie}
\end{proposition}

\begin{proof}
The three statements are proved separately in \cref{sec:pf:liftinj:heckesec}: (a) is \cref{pf:liftinj:hint} (integrality of frames and values of Hecke
operators, which rests on the unit statement \cref{pf:liftinj:thetaunit}); (b) is \cref{pf:liftinj:htr}, applied with the point laws of Hecke
operators, which are obtained from the theta transport law
(\lref{GlobalTheory.HeckePointLaw.heckePointLaws_of_transport}: point laws follow from transport laws;
\lref{GlobalTheory.ThetaTransport.thetaTransportLaw}: the theta transport law holds); and (c) is \cref{pf:liftinj:hread}, applied with the split
linear model and the fixedness of the vacuum stabiliser (\lref{GlobalTheory.SplitLinear.splitLinearModel},
\lref{GlobalTheory.VacStabFixed.vacFixed}); it rests on the matching of Hecke kernels with cosets (\cref{pf:liftinj:hker}).
\end{proof}

\begin{theorem}[The third-layer packages]\label{pf:liftinj:third}
$\mathcal N$ and the geometric compatibility package $\mathcal X$ imply $\mathcal E\wedge\Pi\wedge\mathcal S\wedge\exists h\in\mathcal R:\mathcal L(h)$.
\formallabel{GlobalTheory.T.hyps_e_third_layer_ser_n_of_geometric_galois}
\formallabel{GlobalTheory.MemberClose.exists_hyp_n_of_geometric_galois_compat}
\end{theorem}

\begin{proof}
\emph{Local Langlands and the published inputs.} By \lref{GlobalTheory.Mem.LiteralAutomorphy.localLanglandsGL_and_purity_at_member_alt}
($\mathcal N_{\mathrm{loc}}$ and $\mathcal N_{\mathrm{aut}}$ imply $\mathcal R$ and $\Pi$), we obtain $h\in\mathcal R$ and $\Pi$. In this lemma $h$ is the field
$h_{\mathrm{loc}}$ of $\mathcal N_{\mathrm{loc}}$, and the fields of $\Pi$ are fields of $\mathcal N_{\mathrm{loc}}$ and $\mathcal N_{\mathrm{aut}}$, the
theta-local inputs assembled from them (\lref{GlobalTheory.T.published_inputs_combined_alt_alt}), Green's theorem, base change, the generic unitary
dual of $\GL_n$ (\lref{PositivityTheory.Gud.generic_unitary_dual_GL}), Chebotarev's theorem and Howe's unramified generation
(\lref{GlobalTheory.Hk3.howeUnram_holds}).

\emph{The series package.} Rationality and integrality of twist series follow from global Artin reciprocity
(\lref{GlobalTheory.seriesRatS_and_serInt_of_globalArtin} and \lref{GlobalTheory.SerStdFit.seriesRat_of_seriesRatS}); the three Hecke fields are
\cref{pf:liftinj:hecke}; and the toric ratio statement follows from the toric transport statement of $\mathcal N_{\mathrm{join}}$
(\lref{GlobalTheory.ToricTransport.toricRatioConjugates_of_transport}: transport gives conjugation of toric sums;
\lref{GlobalTheory.ToricRatioBridge.ratioConj_of_toricRatioConjugates}: a Galois descent argument showing that the values of toric sums are
algebraic with Galois conjugates computed by the Artin symbol).

\emph{The elliptic package.} The two Chebotarev statements, Poitou--Tate for totally complex fields, cohomological dimension two, tame local theory,
archimedean signature matching and ray characters are obtained from Poitou--Tate duality (\lref{GlobalTheory.PoitouTateConj.poitouTateDuality_of_normalised},
from its normalised form in $\mathcal N_{\mathrm{coh}}$), the local statements of $\mathcal N_{\mathrm{coh}}$ and Artin reciprocity in
$\mathcal N_{\mathrm{aut}}$, by \lref{GlobalTheory.twist_front_inputs}; the other cohomological fields are fields of $\mathcal N_{\mathrm{coh}}$; the
toric denominator statement is a field of $\mathcal N_{\mathrm{join}}$; the series and Hecke fields are those of $\mathcal S$; finiteness of automorphic forms,
the Galois statement for cuspidal representations and de Rham $\Rightarrow$ potentially semistable are fields of the theta-local inputs in $\Pi$; the
local-global statement above $\ell$ is obtained by combining two fields of the theta-local inputs; and the local Bloch--Kato conclusion at all primes
is \lref{GlobalTheory.BlochKatoLocalOfHyp.blochKatoLocal_allPrimes_of_blochKato}.

\emph{The local Langlands package.} For $h$ as above: $\Pi$ and $\mathcal S$; the archimedean comparison, whose first field is
\lref{GlobalTheory.PhtCompare.phtCmp} (a comparison of Hodge--Tate type data) and whose two other fields are inhabited by trivial elements; the
parameter statement (\lref{GlobalTheory.HeightCounting.Up.conjDualParam_stable_of_prime_not_dividing_level_alt}, from $\Pi$ and the multiplicity
statement \lref{GlobalTheory.MtInMain.prN_sMt_alt_alt}); the point statement (\cref{pf:liftinj:ptst},
from $\Pi$, $\mathcal X$ and $h$); and the cuspidal parameter statement (\lref{GlobalTheory.HeightCounting.Up.cuspParameter_at_member_alt}, from $\Pi$
and $h$).
\end{proof}

\begin{remark}\label{pf:liftinj:thirdrem}
(a) The element $h\in\mathcal R$ produced by the theorem is the field $h_{\mathrm{loc}}$ of $\mathcal N_{\mathrm{loc}}$: the intermediate lemma producing it takes
eleven further local statements as hypotheses and returns its first argument. (b) The archimedean comparison contains two fields which assert the
non-emptiness of types inhabited by a trivial element, and the proof uses that element; the hypotheses of the lemma producing the archimedean
comparison are not used. (c) The functional equation component of $\mathcal N$ is not used, and $\mathcal X$ is used only in the point statement.
\end{remark}

% ---------------------------------------------------------------------------
\subsection{The point conditions}\label{sec:pf:liftinj:D}

\begin{definition}[Point conditions, linked inputs and the jump condition]\label{pf:liftinj:pointdef}
Let $D$ be a global datum.
\begin{enumerate}[label=(\alph*)]
\item The \emph{point conditions} $\mathcal P(D)$ are the conjunction of:
(A) the label tie, a split Hecke statement, a tie at $0$, two admissibility conditions, strong multiplicity one for cuspidal sums and isobaric
multiplicity one;
(B) a row size condition, a third admissibility condition, archimedean signature matching and a statement on the first theta lift;
(C) a joint occurrence condition, conditions on residual representations and on depth zero, and three recursion statements for each local index $i$
(for the parameters of lifts, for the membership of downward lifts, and for good parts);
(T) for every local index $i$ of $D$: every non-zero element $F$ of the eigenspaces attached to a good vector $x$ of index $i$ has non-zero theta lift
$\theta_P(\mu_i,j_i)(F)$;
(Adm) a rank condition for all size data of $D$.
\item The \emph{linked inputs} $\mathcal I$ for $h\in\mathcal R$ and $D$ are a record containing a character choice $P$ (with character frame, norm choice
and tuple datum), real-place data $X$ agreeing with $P$, a reading condition for the good idempotent, and named conditions: a hypothesis on sums of
representations, a global compatibility, a base change conclusion, a statement on the localisation scalars, an isotropy condition, a local series
condition, and $\liftinjP{TNZ}(j,\mu)$ for the real datum $X.R$ and all $j,\mu$. From these the development derives further conditions (among them a
first admissibility condition, a multiplicity tie and local readings) and a map producing a first ``provision'' statement from the summand statement
and the conditions (A), (B).
\item The \emph{jump condition} (J) asserts: for every $a$-dimensional $\ell_i$-adic representation of the absolute Galois group of $E$ ($\ell_i$ an
auxiliary prime of the datum), every place $w$ above $\ell_i$ at which it is de Rham and every embedding $\tau$ inducing $w$, the multiset of jumps of
the Hodge filtration at $\tau$ has $a$ elements.
\item The \emph{theta-local inputs} $\mathcal T$ are a conjunction of local and global statements (Bernstein--Zelevinsky finiteness, irreducibility of
unitary products, the geometric lemma, irreducibility of unlinked products, depth-zero Jacquet injectivity, the Jacquet--Shalika theorem, finiteness
of automorphic forms on $\GL_n$, \dots).
\end{enumerate}
\formal{GlobalTheory.NRows.PointConditions'}
\formal{GlobalTheory.NRows.RowsAAlt}
\formal{GlobalTheory.NRows.RowsB}
\formal{GlobalTheory.NRows.RowsC}
\formal{GlobalTheory.RowCopy.TypeNZRec}
\formal{GlobalTheory.NRows.LinkedInputs}
\formal{Formalism.MainNotions.N.JumpsCard}
\formal{GlobalTheory.NRows.ThetaLocalInputsAlt}
\end{definition}

\begin{theorem}[The point conditions]\label{pf:liftinj:point}
The theta-local inputs $\mathcal T$, Harish-Chandra finiteness and input (a) of \cref{pf:liftinj:lit} imply: for every $h\in\mathcal R$ and every global
datum $D$ with linked inputs $\mathcal I$ satisfying (J), the point conditions $\mathcal P(D)$ hold.
\formallabel{GlobalTheory.T.theta_conditions_of_inputs_of_jumpsCard}
\formallabel{GlobalTheory.RowsAll.locComp_of_isRc_classMapD_cenLoc}
\end{theorem}

\begin{proof}
Let $P$, $X$ be the character choice and real-place data of $\mathcal I$.
\begin{enumerate}[label=\arabic*.]
\item Condition (A) follows from $\mathcal T$ and (J) (\lref{GlobalTheory.T.thetaLocal_a_conditions_of_jumpsCard}, which is
\lref{UpperTheory.GalS.a_close_main}: the theta-local inputs and the jump condition give (A)).
\item Harish-Chandra finiteness and rapid decay give a join $\mathcal J$ of the character choice $P$ with $X$
(\lref{GlobalTheory.T.nonempty_join_of_agree}: such a join exists when $X$ agrees with $P$).
\item The isotropy condition and the statement on localisation scalars give positivity of the localisation scalars and condition (Adm)
(\lref{GlobalTheory.T.thetaLiftSpaces_rank_const_and_localisation_scalars}).
\item From $\mathcal T$, $\mathcal J$, (A), the first admissibility condition and the positivity of step~3: condition (B) and, for every local index, the
pair lift, split base change and local representation statements (\cref{pf:liftinj:bcond}).
\item The linked inputs give the first provision statement from the summand statement (\lref{UpperTheory.GalS.Parts.main_alt_alt}: the summands
suffice for every global datum), (A) and (B).
\item Condition (C) follows from step 4, (A), (B), the first admissibility condition, the isotropy condition and the provision statement
(\lref{GlobalTheory.T.c_statement_of_pair_lift_splitBC_repOfLoc}).
\item For each local index $i$: every non-zero good vector of index $i$ for $X$ is an eigenvector whose system has a segment
(\cref{pf:liftinj:eigstr}, from finiteness of automorphic forms on $\GL_n$, the provision statement, an
admissibility condition and the label tie), and these vectors are good in the sense of \cref{pf:liftinj:realdefs}(d)
(\lref{GlobalTheory.T.eigGood_of_archimedean_signature_adm5}). Together with $\liftinjP{TNZ}(j,\mu)$ for $X.R$ (a field of $\mathcal I$), the
transfer lemma \lref{GlobalTheory.T.type_nonzero_recursion_of_eigen'} (type non-vanishing of the real data, applied to these eigenvectors, gives (T)
at $i$) gives (T) at $i$.
\end{enumerate}
\end{proof}

\begin{remark}\label{pf:liftinj:pointrem}
The theorem does not use \cref{pf:liftinj:main} directly: the type non-vanishing of the real-place data enters as a field of the linked inputs, which
is filled elsewhere in the construction (by \cref{pf:liftinj:real}). The formal name refers to the role of the statement (``theta conditions'') in the
construction.
\end{remark}

% ---------------------------------------------------------------------------
\subsection{Further ingredients}\label{sec:pf:liftinj:further}

This subsection proves results which were cited above: ingredients of the residue identity (\ref{pf:liftinj:R7}, and six results on which the
identification of residue measures in \cref{sec:pf:liftinj:ingr} rests), the three statements on Hecke operators of \cref{pf:liftinj:hecke}, two steps
of the proof of \cref{pf:liftinj:point}, and the point statement used in the proof of \cref{pf:liftinj:third}. Deeper ingredients are again cited,
with a statement of what is used.

\subsubsection{A cut with non-zero eigenvalue}\label{sec:pf:liftinj:jp}

We prove \ref{pf:liftinj:R7}. The cut is built at a split place from data over the residue field, and its eigenvalue is a
finite exponential sum whose value at $t/2$ is a positive real number divided by a positive real number.

Let $O$ be a discrete valuation ring with fraction field $L$, uniformiser $\varpi$ and finite residue field $k$ with $q$ elements. We
write elements of $\GL_{2n}$ in $n\times n$ blocks and put
\[
p(a,d,b)=\begin{pmatrix}a&b\\0&d\end{pmatrix},\qquad n(y)=\begin{pmatrix}1&y\\0&1\end{pmatrix},\qquad \tau_\varpi=p(1,\varpi,0)
\]
for $a,d\in\GL_n$ and $b,y\in\mathrm{Mat}_n$; $P$ is the group of the matrices $p(a,d,b)$, $K=\GL_{2n}(O)$, and $K(\varpi^2)$ is the kernel of
the reduction map $K\to\GL_{2n}(O/\varpi^2)$. For $g\in\GL_{2n}(k)$ let $i(g)=n-\rank d(g)$, where $d(g)$ is the lower right block of $g$; for
$\kappa\in K$, $\bar\kappa\in\GL_{2n}(k)$ is its reduction. A \emph{local height} is a function $h:\GL_{2n}(L)\to\R_{\ge0}$ with
$h(p(a,d,b)\kappa)=\nu(\det a)/\nu(\det d)$ for all $a,d,b$ and all $\kappa\in K$, where $\nu:L^\times\to\R_{>0}$ is a homomorphism which is
trivial on $O^\times$ and satisfies $\nu(\varpi)=q^{-1}$. For a non-trivial character $\psi_0$ of $k$ and $\beta\in\mathrm{Mat}_n(k)$ put
\[
c_\beta(Y)=q^{-n^2}\,\overline{\psi_0(\tr(Y\beta))}\qquad(Y\in\mathrm{Mat}_n(k)).
\]
The local computation rests on five statements, each proved in the development:
\begin{enumerate}[label=(J\arabic*)]
\item\label{pf:liftinj:J1} (\lref{GlobalTheory.T.iwasawa_decomposition_statement}) Every $g\in\GL_{2n}(L)$ is of the form $p(a,d,b)\kappa$
with $\kappa\in K$.
\item\label{pf:liftinj:J2} (\lref{GlobalTheory.T.height_bound_statement}) For every local height $h$ and every $\kappa\in K$:
$h(\kappa\tau_\varpi^{-1})=q^{2i(\bar\kappa)-n}$.
\item\label{pf:liftinj:J3} (\lref{GlobalTheory.T.uniform_boundedness_statement}) If $X$ is a group with a left invariant probability measure
$\nu_X$ and $\rho:X\to K$ is a surjective homomorphism such that $\rho^{-1}(K(\varpi^2))$ is measurable, then for every right
$K(\varpi^2)$-invariant function $\phi$ on $K$ the integral $\int_X\phi(\rho(x))\,d\nu_X(x)$ equals the average of $\phi$ over $K/K(\varpi^2)$.
\item\label{pf:liftinj:J4} (\lref{GlobalTheory.T.compactness_statement}) Let $T:K\to\GL_{2n}(k)$ be a map such that for every $\kappa\in K$ there
are $a,d,b$ and $\kappa'\in K$ with $\kappa\tau_\varpi=p(a,d,b)\kappa'$ and $T(\kappa)=\bar\kappa'$, and let $\Psi$ be a function on $\GL_{2n}(k)$ which
is left invariant under $P(k)$. Then the average of $\Psi\circ T$ over $K/K(\varpi^2)$ equals
$\sum_gq^{2n\,i(g)}\Psi(g)\big/\sum_gq^{2n\,i(g)}$, the sums running over $g\in\GL_{2n}(k)$.
\item\label{pf:liftinj:J5} (\lref{GlobalTheory.T.positivity_statement}) For every finite field $k$, every non-trivial $\psi_0$, every
$\beta\in\mathrm{Mat}_n(k)$ of rank $r$ and every $m\ge r$, the sum
$\sum_{g\in\GL_{2n}(k)}\sum_{Y\in\mathrm{Mat}_n(k)}q^{r\,i(g)}\,c_\beta(Y)\,q^{m\,i(g\,n(Y))}$ is a positive real number.
\end{enumerate}

\begin{lemma}[The local eigenvalue]\label{pf:liftinj:jploc}
Let $n=r+t$ with $r\ge1$, let $\psi_0$ be a non-trivial character of $k$ and $\beta\in\mathrm{Mat}_n(k)$ of rank $r$. For $z\in\C$ and
$g\in\GL_{2n}(k)$ put
\begin{equation}\label{eq:pf:liftinj:cz}
W_z(g)=\sum_{Y\in\mathrm{Mat}_n(k)}c_\beta(Y)\,q^{(i(g\,n(Y))-i(g))(2z+n)},\qquad
c(z)=\frac{\sum_{g}q^{2n\,i(g)}\,W_z(g)}{\sum_{g}q^{2n\,i(g)}} .
\end{equation}
Then $c$ is entire, $c(t/2)\ne0$, and for every uniformiser $\varpi$, every map $\ell:k\to O$ with $\ell(y)\equiv y$ modulo $\varpi$, every
local height $h$, every $(X,\nu_X,\rho)$ as in \ref{pf:liftinj:J3} and every $z\in\C$,
\begin{equation}\label{eq:pf:liftinj:jploc}
\int_X\sum_{Y\in\mathrm{Mat}_n(k)}c_\beta(Y)\,h\bigl(\rho(x)\,n(\varpi^{-1}\ell(Y))\bigr)^{z+n/2}\,d\nu_X(x)=c(z),
\end{equation}
where $\ell(Y)$ is obtained by applying $\ell$ to the entries of $Y$.
\formallabel{GlobalTheory.T.alternative_local_computation_of_components}
\formallabel{GlobalTheory.CompanionTheory.Jp.jpAlt_of_steps}
\end{lemma}

\begin{proof}
The function $c$ is a finite linear combination of the functions $z\mapsto q^{2mz}$, $m\in\Z$, divided by a positive number; hence it is entire.
At $z=t/2$ we have $2z+n=n+t$ and $2n-(n+t)=r$, so that
\[
\sum_gq^{2n\,i(g)}W_{t/2}(g)=\sum_g\sum_Yq^{r\,i(g)}\,c_\beta(Y)\,q^{(n+t)\,i(g\,n(Y))},
\]
which is a positive real number by \ref{pf:liftinj:J5} with $m=n+t\ge r$. Hence $c(t/2)\ne0$.

Fix $\varpi,\ell,h$. By \ref{pf:liftinj:J1} and the defining property, a local height satisfies $h(gk)=h(g)$ for $k\in K$ and
$h(p(a,d,b)g)=\nu(\det a)\nu(\det d)^{-1}h(g)$ for all $g$, since $p(a,d,b)p(a',d',b')=p(aa',dd',ab'+bd')$; in particular $h(\kappa)=1$ for
$\kappa\in K$. For $\kappa\in K$ choose, by \ref{pf:liftinj:J1}, a decomposition $\kappa\tau_\varpi=p(a,d,b)\kappa'$ with $\kappa'\in K$, and put
$T(\kappa)=\bar\kappa'$. Since $\tau_\varpi^{-1}n(\varpi^{-1}\ell(Y))\tau_\varpi=n(\ell(Y))$, we have
$\kappa\,n(\varpi^{-1}\ell(Y))=p(a,d,b)\,\kappa'n(\ell(Y))\,\tau_\varpi^{-1}$, and \ref{pf:liftinj:J2}, applied to $\kappa'n(\ell(Y))\in K$ whose
reduction is $T(\kappa)n(Y)$, gives
\[
h\bigl(\kappa\,n(\varpi^{-1}\ell(Y))\bigr)=\frac{\nu(\det a)}{\nu(\det d)}\,q^{2i(T(\kappa)n(Y))-n}.
\]
For $Y=0$ the left side is $h(\kappa)=1$. Dividing, $h(\kappa\,n(\varpi^{-1}\ell(Y)))=q^{2(i(T(\kappa)n(Y))-i(T(\kappa)))}$, and therefore
\begin{equation}\label{eq:pf:liftinj:hred}
\sum_Yc_\beta(Y)\,h\bigl(\kappa\,n(\varpi^{-1}\ell(Y))\bigr)^{z+n/2}=W_z(T(\kappa))\qquad(\kappa\in K).
\end{equation}
The left side of \eqref{eq:pf:liftinj:hred} is right $K(\varpi^2)$-invariant in $\kappa$: if $\gamma=1+\gamma'\in K(\varpi^2)$ and
$x=\varpi^{-1}\ell(Y)$, then $n(x)^{-1}\gamma\,n(x)=1+n(-x)\gamma'n(x)$ has integral entries, because $\gamma'\in\varpi^2\mathrm{Mat}_{2n}(O)$ and
the entries of $n(\pm x)$ lie in $\varpi^{-1}O$, and its determinant $\det\gamma$ is a unit; so $n(x)^{-1}\gamma\,n(x)\in K$ and
$h(\kappa\gamma\,n(x))=h(\kappa\,n(x))$. By \ref{pf:liftinj:J3} and \eqref{eq:pf:liftinj:hred}, the left side of \eqref{eq:pf:liftinj:jploc}
equals the average of $W_z\circ T$ over $K/K(\varpi^2)$. Finally $W_z$ is left $P(k)$-invariant: the lower right block of $p(a,d,b)g$ is
$d\,d(g)$, so $i(pg)=i(g)$ and $i(pg\,n(Y))=i(g\,n(Y))$ for $p\in P(k)$. The map $T$ has the property required in \ref{pf:liftinj:J4}, and
\ref{pf:liftinj:J4} shows that the average equals $c(z)$.
\end{proof}

\begin{proposition}[A cut with non-zero eigenvalue]\label{pf:liftinj:jpos}
Statement \ref{pf:liftinj:R7} holds. That is, let $n=r+t$ with $r\ge1$, $\psi$ standard, $\chi$ an $r$-adapted Hecke character and $n(\cdot)$ a
unipotent coordinate. For every finite set $S$ of finite places of $F$ there are a finite place $v\notin S$, a cut $\liftinjcut$ at $v$ killing
rank $<r$ and an entire function $c$ with $c(t/2)\ne0$ such that for every Haar probability measure $\nu_v$ on $K_v$, every $z\in\C$, every
continuous right $K_v$-invariant function $f$ on $G_n$ of law $(\chi,z)$ and every $g\in G_n$:
$\int_{K_v}(\liftinjcut f)(gk)\,d\nu_v(k)=c(z)f(g)$.
\formallabel{GlobalTheory.CompanionTheory.Jp.jPos_of_idemEx}
\end{proposition}

\begin{proof}
\emph{Step 1: a good place.} By \lref{GlobalTheory.T.jacquet_piatetski_bound_holds}, for every finite set $S$ there are a finite place
$v\notin S$ such that $\psi_v$ is trivial on $O_v$ but not on $\varpi_v^{-1}O_v$, and an $F$-embedding $\tau:E\to F_v$ (so that $v$ splits in $E$)
such that $\chi_v(a)=1$ for every $a\in(E\otimes_FF_v)^\times$ whose images under $\tau$ and under $\tau$ composed with complex conjugation have the
same absolute value.

\emph{Step 2: the local eigenvalue.} \Cref{pf:liftinj:jploc} holds for every discrete valuation ring with finite residue field. By
\lref{GlobalTheory.T.split_place_rank_cut_of_alternative}, this implies: at every place $v$ such that $\psi_v$ is trivial on $O_v$ but not on
$\varpi_v^{-1}O_v$, and for every $F$-embedding $\tau:E\to F_v$, there are a cut $\liftinjcut$ at $v$ killing rank $<r$ and an entire function $c$ with
$c(t/2)\ne0$ such that
\[
\int_{K_v}\sum_{x\in\liftinjcut}c_x\,\phi_z\bigl(k\,n(x)\bigr)\,d\nu_v(k)=c(z)
\]
for every Haar probability measure $\nu_v$ on $K_v$ and every $z$, where $\phi_z=H_v^{z+n/2}$ and $H_v$ is the height of the local group
$G_{n,v}$ relative to $K_v$. In that proof $c$ is the function \eqref{eq:pf:liftinj:cz} for $O=O_v$, the matrix
$\beta=\mathrm{diag}(1,\dots,1,0,\dots,0)$ of rank $r$ and a non-trivial character $\psi_0$ of the residue field attached to $\psi_v$; the cut consists
of the hermitian matrices whose only non-zero component is at $v$, where it corresponds under $\tau$ to $\varpi_v^{-1}\ell(Y)$, with the
coefficients $c_\beta(Y)$; and $X=K_v$ is mapped onto $\GL_{2n}(O_v)$ by means of $\tau$.

\emph{Step 3: eigenfunctions.} By \lref{GlobalTheory.T.section_eigenfunction_maximal_compact}, if $v,\tau,\chi$ are as in Step 1, $\liftinjcut$
is a cut at $v$, $\nu_v$ a Haar measure on $K_v$ and $c'=\int_{K_v}\sum_{x}c_x\phi_z(k\,n(x))\,d\nu_v(k)$, then every right $K_v$-invariant
function $f$ on $G_n$ of law $(\chi,z)$ satisfies $\int_{K_v}(\liftinjcut f)(gk)\,d\nu_v(k)=c'f(g)$ for all $g$.

The place $v$ of Step 1, with the cut and the function $c$ given by Step 2 for this $v$ and $\tau$, has the required properties by Step 3.
\end{proof}

\begin{remark}\label{pf:liftinj:jposrem}
The formal statement of \cref{pf:liftinj:jpos} takes the existence of $S$-idempotents (\ref{pf:liftinj:R3}) as a hypothesis; it is not used.
\end{remark}

\subsubsection{A bound for spherical Eisenstein series}\label{sec:pf:liftinj:sphsec}

\begin{proposition}[Spherical bound]\label{pf:liftinj:sph}
Let $D$ be a global datum, $p\ge1$, $H$ a height on $G_p$ (\cref{pf:liftinj:schrdef}) and $C\subseteq G_p$ compact. There is $B\in\R$ such that for
every real $\sigma\in(p/2,p/2+1]$ and every $h\in C$ the series $\sum_{\gamma}H(\gamma h)^{\sigma+p/2}$, over a set of representatives $\gamma$
of $(P_p\cap\Gamma_p)\backslash\Gamma_p$, converges and
\[
(\sigma-\tfrac p2)\sum_{\gamma}H(\gamma h)^{\sigma+p/2}\le B .
\]
\formallabel{PositivityTheory.Rg.Rsw.sphBound}
\formallabel{GlobalTheory.T.spherical_bound_of_heights}
\formallabel{GlobalTheory.CompanionTheory.SphPole.sphBound_of_htExist}
\end{proposition}

\begin{proof}
We use two results of the development. First (\lref{PositivityTheory.Rg.Rsw.htExist}), there is a compact set $C_0\subseteq G_p$ with
$G_p=P_pC_0$. Second (\lref{GlobalTheory.CompanionTheory.SphPole.height_bound}), there are a height $H'$ on $G_p$ and $B'\in\R$ such that
for $0<\eps\le1$ the series $\sum_\gamma H'(\gamma)^{p+\eps}$ converges and $\eps\sum_\gamma H'(\gamma)^{p+\eps}\le B'$. This is the bound at
$h=1$ for one particular height; it carries the analytic content (convergence of the spherical Eisenstein series to the right of $p/2$, with
at most a simple pole at $p/2$).

The function $(u_1,u_2)\mapsto H(u_1u_2)/H'(u_1)$ is continuous on $G_p\times G_p$, hence bounded on the compact set $C_0\times C$ by some $c$;
put $c_1=\max(c,1)$. Let $\gamma\in\Gamma_p$ and $h\in C$, and write $\gamma=q\,c_0$ with $q\in P_p$ and $c_0\in C_0$; let $m(a)$ be the Levi
component of $q$. By the transformation law of heights,
\[
H(\gamma h)=|\det a|\,H(c_0h)\le|\det a|\,c_1H'(c_0)=c_1H'(\gamma).
\]
Let $\eps=\sigma-p/2\in(0,1]$, so that $\sigma+p/2=p+\eps$. Since $c_1\ge1$, we get $H(\gamma h)^{p+\eps}\le c_1^{p+1}H'(\gamma)^{p+\eps}$. Hence
the series for $H$ converges, and
$(\sigma-p/2)\sum_\gamma H(\gamma h)^{\sigma+p/2}\le c_1^{p+1}\eps\sum_\gamma H'(\gamma)^{p+\eps}\le c_1^{p+1}B'$. The proposition holds with
$B=c_1^{p+1}B'$.
\end{proof}

\subsubsection{The Siegel--Weil formula in the convergent range}\label{sec:pf:liftinj:swsec}

Let $E/F$ be a CM extension, $J\in\mathrm{Mat}_M(E)$ a non-degenerate hermitian matrix and $\Uni(J)$ its unitary group. For
$X=(X_1,\dots,X_n)\in(\A_E^M)^n$ let $\mathrm{gr}(X)=({}^t\bar X_aJX_b)_{a,b}$ be its Gram matrix; $\Uni(J)(\A_F)$ acts on $n$-tuples
componentwise. For a Schwartz--Bruhat function $\Phi$ on $(\A_E^M)^n$, measures $dX$ on $(\A_E^M)^n$ and $db$ on the space
$\mathrm{Herm}_n(\A_F)$ of hermitian $n\times n$ matrices over $\A_E$, and $T\in\mathrm{Herm}_n(E)$, put
\begin{equation}\label{eq:pf:liftinj:coefN}
\Phi^\star(b)=\int\Phi(X)\,\psi\bigl(\tfrac12\tr(b\,\mathrm{gr}(X))\bigr)\,dX,\qquad
\mathrm{co}_T(\Phi)=\int_{\mathrm{Herm}_n(\A_F)}\Phi^\star(b)\,\psi\bigl(-\tfrac12\tr(Tb)\bigr)\,db,
\end{equation}
where $\tr$ is the trace of matrices followed by the trace from $\A_E$ to $\A_F$. Thus $\Phi^\star(b)$ is the integral of the function
obtained from $\Phi$ by the action of $n(b)$ in the Schr\"odinger model \eqref{eq:pf:liftinj:schr}, and $\mathrm{co}_T(\Phi)$ has the shape of
the $T$-th Fourier coefficient of the Siegel Eisenstein series of $\Phi$. For a Haar measure $\mu$ on $\Uni(J)(\A_F)$ with a fundamental domain
$\mathcal F$ for $\Uni(J)(F)$, the \emph{theta measure} of $T$ is the measure on $(\A_E^M)^n$
\begin{equation}\label{eq:pf:liftinj:thetameas}
\vartheta_T=\mu(\mathcal F)^{-1}\sum_{X\in(E^M)^n,\ \mathrm{gr}(X)=T}\bigl(\text{image of }\mu|_{\mathcal F}\text{ under }h\mapsto h^{-1}X\bigr),
\end{equation}
so that $\int\Phi\,d\vartheta_T=\mu(\mathcal F)^{-1}\sum_{\mathrm{gr}(X)=T}\int_{\mathcal F}\Phi(h^{-1}X)\,d\mu(h)$ for $\Phi\ge0$. For each place
$v$ of $F$ (including the archimedean place $\infty$, for which $F_\infty=F\otimes\R$) and each non-singular hermitian $T'$ over $F_v$, the
\emph{Gram fibre measure} $m_v(T')$ is a measure on $\{X\in((E\otimes F_v)^M)^n:\mathrm{gr}(X)=T'\}$; these measures disintegrate the standard
measure on the set of $X$ with non-singular Gram matrix along $\mathrm{gr}$, with respect to the standard measure on hermitian matrices, and
depend continuously on $T'$. They exist at every place (\lref{GlobalTheory.CompanionTheory.SW.Step.loc_families}). $L_v$ denotes the
standard lattice of $((E\otimes F_v)^M)^n$.

\begin{lemma}[Pure tensors]\label{pf:liftinj:swpure}
Let $2n<M$, let $\psi$ be standard, $dX$ and $db$ the self-dual measures for $\psi$, $\mu$ a Haar measure on $\Uni(J)(\A_F)$ with a
fundamental domain $\mathcal F$ for $\Uni(J)(F)$ and $T\in\mathrm{Herm}_n(E)$ non-singular. Let $\Phi(X)=\phi(X_\infty)\,\mathbf 1_A(X_f)$, where $\phi$
is a Schwartz function of the archimedean components and $A=\prod_vA_v$ is a product of compact open sets $A_v$ with $A_v=L_v$ for all $v$
outside a finite set. Then $\Phi$ is $\vartheta_T$-integrable and $\mathrm{co}_T(\Phi)=\int\Phi\,d\vartheta_T$.
\formallabel{GlobalTheory.CompanionTheory.SiegelND.Coef.pure_identity}
\end{lemma}

\begin{proof}
We use four facts.
\begin{enumerate}[label=(\alph*)]
\item (\lref{GlobalTheory.T.siegel_weil_measure_of_rank_small_or_zero}) Since $n=0$ or $n+2\le M$, $\vartheta_T$ is the restricted product
of the measures $m_v(T)$ with constant $1$: $m_v(T)(L_v)<\infty$ for all $v$, $\sum_v|m_v(T)(L_v)-1|<\infty$, and there is a finite set $S_0$
such that for all finite $S\supseteq S_0$ and all non-negative measurable $\phi_\infty$ and $\phi_v$ ($v\in S$)
\[
\int\phi_\infty\otimes\bigotimes_{v\in S}\phi_v\otimes\bigotimes_{v\notin S}\mathbf 1_{L_v}\,d\vartheta_T
=\int\phi_\infty\,dm_\infty(T)\,\prod_{v\in S}\int\phi_v\,dm_v(T)\,\prod_{v\notin S}m_v(T)(L_v).
\]
\item (\lref{GlobalTheory.T.siegel_weil_convergent_range}) For every Schwartz--Bruhat function $\Phi$ the function $\Phi^\star$ is
integrable on $\mathrm{Herm}_n(\A_F)$, since $2n<M$ (the convergence estimate of Weil \cite{Weil1965}).
\item (\lref{GlobalTheory.CompanionTheory.SiegelND.Coef.coefS_pure}) For a finite set $S$ let $\mathrm{co}_{T,S}(\Phi)$ be the integral
\eqref{eq:pf:liftinj:coefN} restricted to the $b$ whose components at the finite places $v\notin S$ lie in the standard lattice of hermitian
matrices. There is a finite set $S_1$ such that for every finite $S\supseteq S_1$ and every $\Phi$ as in the lemma with $A_v=L_v$ for
$v\notin S$: $\mathrm{co}_{T,S}(\Phi)=\int\phi\,dm_\infty(T)\,\prod_{v\in S}m_v(T)(A_v)$.
\item (\lref{GlobalTheory.T.archimedean_gram_fibre_measure_formula}, \lref{GlobalTheory.T.local_gram_fibre_measure_standard_volume})
$\phi$ is $m_\infty(T)$-integrable, and $m_v(T)(A_v)<\infty$ for every compact open $A_v$.
\end{enumerate}
Applying (a) to $|\phi|$ and to the positive and negative parts of the real and imaginary parts of $\phi$, and using (d), we find that for every
finite $S\supseteq S_0$ containing the places where $A_v\ne L_v$ the function $\Phi$ is $\vartheta_T$-integrable and
\begin{equation}\label{eq:pf:liftinj:swe}
\int\Phi\,d\vartheta_T=\int\phi\,dm_\infty(T)\,\prod_{v\in S}m_v(T)(A_v)\,\prod_{v\notin S}m_v(T)(L_v)
\end{equation}
(\lref{GlobalTheory.T.euler_product_measure_eventually_standard}). By (b) and dominated convergence, $\mathrm{co}_{T,S}(\Phi)\to\mathrm{co}_T(\Phi)$
as $S$ increases through the finite sets of places, because every finite adele is integral at almost all places. By the summability in (a),
$\prod_{v\notin S}m_v(T)(L_v)\to1$. By (c) and \eqref{eq:pf:liftinj:swe}, $\mathrm{co}_{T,S}(\Phi)\prod_{v\notin S}m_v(T)(L_v)=\int\Phi\,d\vartheta_T$
for all sufficiently large $S$, while the left side tends to $\mathrm{co}_T(\Phi)$. Hence $\mathrm{co}_T(\Phi)=\int\Phi\,d\vartheta_T$.
\end{proof}

\begin{theorem}[Siegel--Weil formula in the convergent range]\label{pf:liftinj:sw}
Let $2n<M$, let $\psi$ be standard, $dX$ and $db$ the self-dual measures for $\psi$, $\mu$ a Haar measure on $\Uni(J)(\A_F)$ with a fundamental
domain $\mathcal F$ for $\Uni(J)(F)$, $\Phi$ a Schwartz--Bruhat function on $(\A_E^M)^n$ and $T\in\mathrm{Herm}_n(E)$ non-singular. Then
$\mathrm{co}_T(\Phi)=\int\Phi\,d\vartheta_T$.
\formallabel{GlobalTheory.T.siegel_weil_hermitian_formula}
\formallabel{GlobalTheory.CompanionTheory.SiegelND.Coef.coefN_eq_theta}
\end{theorem}

This is the coefficientwise form of the Siegel--Weil formula of Weil \cite{Weil1965} for unitary groups in the range of absolute
convergence.

\begin{proof}
By \lref{GlobalTheory.T.schwartz_bruhat_finite_sum_of_products}, $\Phi=\sum_{i=1}^N\Phi_i$ with $\Phi_i$ as in \cref{pf:liftinj:swpure}: a
Schwartz--Bruhat function is a finite sum of products of a Schwartz function of the archimedean components with the indicator function of a
product of compact open sets, almost all of which are standard lattices. Such products are Schwartz--Bruhat functions
(\lref{GlobalTheory.T.factorizable_schwartz_orbital_integral}). By \cref{pf:liftinj:swpure}, each $\Phi_i$ is $\vartheta_T$-integrable with
$\mathrm{co}_T(\Phi_i)=\int\Phi_i\,d\vartheta_T$. The map $\Phi\mapsto\Phi^\star$ is linear and each $\Phi_i^\star$ is integrable (fact (b) in the
proof of \cref{pf:liftinj:swpure}); hence $\mathrm{co}_T(\Phi)=\sum_i\mathrm{co}_T(\Phi_i)=\sum_i\int\Phi_i\,d\vartheta_T=\int\Phi\,d\vartheta_T$. In
particular $\Phi$ is $\vartheta_T$-integrable.
\end{proof}

\begin{proposition}[Fibration over a Gram matrix]\label{pf:liftinj:fib}
Let $D$ be a global datum, $V$ a hermitian space and $r\ge0$ with $2r<\dim V$, $\psi$ standard, $\theta$ a family of archimedean embeddings, $\mu$
a measure of covolume one on $\mathrm{Mat}_{\dim V\times r}(\A_E)$, $\lambda$ a measure of covolume one on $\mathrm{Herm}_r(\A_F)$ and
$b_0\in\mathrm{Herm}_r(E)$ non-singular. For $\varphi\in\mathcal S(V^r)$ put
\[
\mathrm W(\varphi)=\int_{\mathrm{Herm}_r(\A_F)}\Bigl(\int\varphi(y)\,\psi\bigl(\tfrac12\tr(b\langle y,y\rangle_V)\bigr)\,d\mu(y)\Bigr)\,
\psi\bigl(-\tfrac12\tr(b_0b)\bigr)\,d\lambda(b).
\]
Then there is a measure $\mathfrak f$ on $\mathrm{Mat}_{\dim V\times r}(\A_E)$ which is finite on compact sets, vanishes outside
$\{y:\langle y,y\rangle_V=b_0\}$ and is invariant under $G_V$, such that every $\varphi\in\mathcal S(V^r)$ is $\mathfrak f$-integrable with
$\mathrm W(\varphi)=\int\varphi\,d\mathfrak f$.
\formallabel{PositivityTheory.Rg.Rsw.L.fibEx}
\end{proposition}

\begin{proof}
Let $J$ be the Gram matrix of $V$, so that $G_V=\Uni(J)(\A_F)$. There is a Haar measure $\mu_V$ on $G_V$ with a fundamental domain $\mathcal F$ for
$\Gamma_V$ of measure $1$, since $[G_V]$ has finite volume \cite{Borel1963} (\lref{GlobalTheory.Published.SiegelOneVector.exists_normalised}).
By \lref{GlobalTheory.T.exists_self_dual_siegel_weil_measures}, there are self-dual measures $dX$, $db$ for $\psi$ such that
$\mathrm W(\varphi)=\mathrm{co}_{b_0}(\tilde\varphi)$ for every $\varphi$, where $\tilde\varphi$ is $\varphi$ viewed as a function of $r$-tuples of
columns. Let $\mathfrak f$ be the image of $\vartheta_{b_0}$ (formed with $\mu_V$ and $\mathcal F$) under the identification of $r$-tuples with
matrices. By \lref{PositivityTheory.PlacePermOperatorBicommutantSpan.theta_inv}, $\vartheta_{b_0}$ is invariant under $\Uni(J)(\A_F)$ and vanishes
outside $\{X:\mathrm{gr}(X)=b_0\}$; this gives the invariance and the support property of $\mathfrak f$. For $\varphi\in\mathcal S(V^r)$ the function
$\tilde\varphi$ is Schwartz--Bruhat (\lref{GlobalTheory.T.schwartz_function_schwartz_bruhat}), and Schwartz--Bruhat functions are
$\vartheta_{b_0}$-integrable by the proof of \cref{pf:liftinj:sw}; since every compact set lies in the set where some non-negative
Schwartz--Bruhat function is at least $1$, $\mathfrak f$ is finite on compact sets. Finally \cref{pf:liftinj:sw}, applied with $M=\dim V$, $n=r$ and
$T=b_0$, gives $\mathrm W(\varphi)=\mathrm{co}_{b_0}(\tilde\varphi)=\int\tilde\varphi\,d\vartheta_{b_0}=\int\varphi\,d\mathfrak f$.
\end{proof}

\subsubsection{Fourier--Jacobi coefficients of Eisenstein series}\label{sec:pf:liftinj:fjsec}

\begin{definition}[Fourier--Jacobi data]\label{pf:liftinj:fjdef}
Let $n=\rho+p$, let $\psi$ be standard and $b_0\in\mathrm{Herm}_\rho(E)$; put $\psi_{b_0}(z)=\psi(\frac12\tr(b_0z))$ for $z\in\mathrm{Herm}_\rho(\A_F)$ and
$\langle x,y\rangle_{b_0}={}^t\bar x\,b_0\,y$ for $x,y\in\mathrm{Mat}_{\rho\times p}(\A_E)$.
\begin{enumerate}[label=(\alph*)]
\item \emph{Fourier--Jacobi maps} consist of a continuous homomorphism $\mathrm{emb}:G_p\to G_n$, continuous maps
$\mathrm{nz}:\mathrm{Herm}_\rho(\A_F)\times\mathrm{Mat}_{\rho\times p}(\A_E)\to G_n$ and $\mathrm{sh}:\mathrm{Mat}_{\rho\times p}(\A_E)\to G_n$, and an
element $w'\in G_n$, whose matrices in the block decomposition attached to $n=\rho+p$ are
\[
\mathrm{emb}\begin{pmatrix}a&b\\c&d\end{pmatrix}=\begin{pmatrix}\mathrm{diag}(1,a)&\mathrm{diag}(0,b)\\\mathrm{diag}(0,c)&\mathrm{diag}(1,d)\end{pmatrix},\qquad
\mathrm{nz}(z,x)=\begin{pmatrix}1&\bigl(\begin{smallmatrix}z&x\\{}^t\bar x&0\end{smallmatrix}\bigr)\\0&1\end{pmatrix},
\]
$\mathrm{sh}(y)=\mathrm{diag}\bigl(\bigl(\begin{smallmatrix}1&y\\0&1\end{smallmatrix}\bigr),\bigl(\begin{smallmatrix}1&0\\-{}^t\bar y&1\end{smallmatrix}\bigr)\bigr)$
and $w'=\bigl(\begin{smallmatrix}\mathrm{diag}(0,1)&\mathrm{diag}(1,0)\\\mathrm{diag}(-1,0)&\mathrm{diag}(0,1)\end{smallmatrix}\bigr)$.
\emph{Fourier--Jacobi measures} are a measure $\lambda_z$ of covolume one on $\mathrm{Herm}_\rho(\A_F)$ with a fundamental domain $\mathcal G$ for
$\mathrm{Herm}_\rho(E)$ and a measure $\mu$ of covolume one on $\mathrm{Mat}_{\rho\times p}(\A_E)$ with a fundamental domain $\mathcal F$.
\item For functions $F,f$ on $G_n$, $g\in G_n$, a function $\phi$ on $\mathrm{Mat}_{\rho\times p}(\A_E)$ and a function $\Omega$ on $G_p$ with values in
functions on $\mathrm{Mat}_{\rho\times p}(\A_E)$, put
\begin{align*}
F_{b_0}(g)&=\int_{\mathcal G}F(\mathrm{nz}(z,0)g)\,\overline{\psi_{b_0}(z)}\,d\lambda_z(z),\\
\mathrm{FJ}(F,g,\phi)&=\int\Bigl(\int_{\mathcal F}F_{b_0}(\mathrm{nz}(0,x)\,\mathrm{sh}(y)\,g)\,d\mu(x)\Bigr)\overline{\phi(y)}\,d\mu(y),\\
f^{w}(g)&=\int_{\mathrm{Herm}_\rho(\A_F)}f(w'\,\mathrm{nz}(z,0)\,g)\,\overline{\psi_{b_0}(z)}\,d\lambda_z(z),\\
\mathrm R(f,g,\Omega)(h)&=\int f^{w}(\mathrm{sh}(y)\,\mathrm{emb}(h)\,g)\,\overline{\Omega(h)(y)}\,d\mu(y)\qquad(h\in G_p).
\end{align*}
$\mathrm{FJ}(F,g,\phi)$ is the \emph{Fourier--Jacobi coefficient} of $F$ and $\mathrm R(f,g,\Omega)$ the \emph{Fourier--Jacobi section} of $f$.
\item A \emph{theta section} for $(\psi,\chi',\theta,b_0,\mu)$ is a jointly continuous function $\Omega:G_p\times\mathrm{Mat}_{\rho\times p}(\A_E)\to\C$
with $\Omega(h)\in\mathcal S(\mathrm{Mat}_{\rho\times p})$ for all $h$, locally uniformly dominated by Schwartz functions, such that
\[
\Omega(m(a)h)(y)=\chi'(\det a)|\det a|^{\rho/2}\,\Omega(h)(ya),\quad \Omega(n(b)h)(y)=\psi\bigl(\tfrac12\tr(b\langle y,y\rangle_{b_0})\bigr)\Omega(h)(y),
\]
and $\Omega(w_ph)(x)=\int\Omega(h)(y)\,\psi(\tr\langle x,y\rangle_{b_0})\,d\mu(y)$; its \emph{Heisenberg theta function} is
$\mathrm{th}(\phi)(z,x,y)=\psi_{b_0}(z)\sum_{\xi\in\mathrm{Mat}_{\rho\times p}(E)}\psi(\tr\langle x,\xi\rangle_{b_0})\,\phi(\xi+y)$.
\item For a function $f$ on $G_p$, $E_p(f)(h)=\sum_{\gamma}f(\gamma h)$, the sum over representatives of $(P_p\cap\Gamma_p)\backslash\Gamma_p$.
\end{enumerate}
\formal{PositivityTheory.Rg.Rsw.Maps}
\formal{PositivityTheory.Rg.Rsw.Meas}
\formal{PositivityTheory.Rg.Rsw.zc}
\formal{PositivityTheory.Rg.Rsw.pc}
\formal{PositivityTheory.Rg.Rsw.FJ}
\formal{PositivityTheory.Rg.Rsw.Wz}
\formal{PositivityTheory.Rg.Rsw.Rsec}
\formal{PositivityTheory.Rg.Rsw.IsTh}
\formal{PositivityTheory.Rg.Rsw.th}
\formal{PositivityTheory.Rg.Rsw.eisS}
\end{definition}

\begin{proposition}[Fourier--Jacobi coefficients of Siegel Eisenstein series]\label{pf:liftinj:fj}
Let $n=\rho+p$, fix Fourier--Jacobi maps and measures, let $\psi$ be standard, $b_0\in\mathrm{Herm}_\rho(E)$ non-singular, $\chi,\chi'$ Hecke characters,
$\theta$ a family of archimedean embeddings and $\Omega$ a theta section for $(\psi,\chi',\theta,b_0,\mu)$. Let $s\in\C$ with $\Re s>n/2$ and let $f$
be a continuous function on $G_n$ of law $(\chi,s)$. Then for every $g\in G_n$
\[
\mathrm{FJ}\bigl(E(\cdot,f),g,\Omega(1)\bigr)=E_p\bigl(\mathrm R(f,g,\Omega)\bigr)(1).
\]
\formallabel{PositivityTheory.Rg.Rsw.fjEis}
\formallabel{GlobalTheory.T.fourier_jacobi_eisenstein}
\formallabel{GlobalTheory.CompanionTheory.FJE.fjEis}
\end{proposition}

That is, the Fourier--Jacobi coefficient of the Siegel Eisenstein series on $G_n$ is the Siegel Eisenstein series on $G_p$ of the Fourier--Jacobi
section, evaluated at $1$.

\begin{proof}
We use the following results of the development.
\begin{enumerate}[label=(\alph*)]
\item (Unfolding, \lref{GlobalTheory.T.fourier_jacobi_unfolding}.) For every $g'\in G_n$ the family
$(q,\eta)\mapsto f^{w}(\mathrm{sh}(\eta)\,\mathrm{emb}(q)\,g')$, indexed by representatives $q$ of $(P_p\cap\Gamma_p)\backslash\Gamma_p$ and
$\eta\in\mathrm{Mat}_{\rho\times p}(E)$, is summable with sum $E(\cdot,f)_{b_0}(g')$, and for every compact set of $g'$ it is dominated by a summable family
independent of $g'$.
\item (Theta twist, \lref{GlobalTheory.T.theta_twist_identity}.) If $\gamma\in\Gamma_p$ and
$\mathrm{emb}(\gamma)\,\mathrm{nz}(z,x)\,\mathrm{sh}(y)=\mathrm{nz}(z',x')\,\mathrm{sh}(y')\,\mathrm{emb}(\gamma)$, then
$\mathrm{th}(\Omega(1))(z,x,y)=\mathrm{th}(\Omega(\gamma))(z',x',y')$.
\item (Folding, \lref{GlobalTheory.T.fourier_jacobi_expansion_bounded}.) If $F$ is left $\Gamma_n$-invariant and
$(y,x)\mapsto F_{b_0}(\mathrm{nz}(0,x)\,\mathrm{sh}(y)\,g)$ is bounded and measurable, then for every Schwartz function $\phi$
\[
\mathrm{FJ}(F,g,\phi)=\int_{\mathcal F\times\mathcal F}F_{b_0}(\mathrm{nz}(0,x)\,\mathrm{sh}(y)\,g)\,\overline{\mathrm{th}(\phi)(0,x,y)}\,d\mu(y)\,d\mu(x).
\]
\item (Terms, \lref{GlobalTheory.T.fourier_jacobi_terms_summable} and \lref{GlobalTheory.T.theta_fourier_jacobi_lattice_swap}.)
Write $S_q(y,x)=\sum_\eta f^{w}(\mathrm{sh}(\eta)\,\mathrm{emb}(q)\,\mathrm{nz}(0,x)\,\mathrm{sh}(y)\,g)$. There are $t_q\ge0$ with $\sum_qt_q<\infty$ and
$|S_q|\le t_q$. Moreover, for every $q$,
\[
\int_{\mathcal F\times\mathcal F}S_q(y,x)\,\overline{\mathrm{th}(\Omega(1))(0,x,y)}\,d\mu(y)\,d\mu(x)=\mathrm R(f,g,\Omega)(q);
\]
this is proved by rewriting the integrand, by means of (b), in the coordinates of $\mathrm{emb}(q)\,\mathrm{nz}(0,x)\,\mathrm{sh}(y)\,\mathrm{emb}(q)^{-1}$
(the lattice swap), by the invariance of the integral over $\mathcal F\times\mathcal F$ under the lattice transformations which occur
(\lref{GlobalTheory.T.lattice_transform_invariance}), and by the evaluation of the resulting twisted lattice sums
(\lref{GlobalTheory.T.twisted_lattice_sums_bounded_product}).
\end{enumerate}
By the coset decomposition of the Eisenstein series, $E(\cdot,f)$ is left $\Gamma_n$-invariant. By (a), applied at
$g'=\mathrm{nz}(0,x)\mathrm{sh}(y)g$, we have $E(\cdot,f)_{b_0}(\mathrm{nz}(0,x)\mathrm{sh}(y)g)=\sum_qS_q(y,x)$, which is bounded by $\sum_qt_q$ and
measurable. By (c) with $\phi=\Omega(1)$,
\[
\mathrm{FJ}\bigl(E(\cdot,f),g,\Omega(1)\bigr)=\int_{\mathcal F\times\mathcal F}\sum_qS_q(y,x)\,\overline{\mathrm{th}(\Omega(1))(0,x,y)}\,d\mu(y)\,d\mu(x).
\]
The measure of $\mathcal F\times\mathcal F$ is $1$ and the theta function of the Schwartz function $\Omega(1)$ is bounded
(\lref{GlobalTheory.T.exists_bound_lattice_sum_norm_schwartzA}), say by $B_\theta$; since $\sum_qt_qB_\theta<\infty$, the sum over $q$ may be
taken outside the integral, and (d) gives $\sum_q\mathrm R(f,g,\Omega)(q\cdot1)=E_p(\mathrm R(f,g,\Omega))(1)$.
\end{proof}

\subsubsection{Residue measures on hermitian cells}\label{sec:pf:liftinj:cellsec}

\begin{definition}[Residue measures, sections on cells and orbital measures]\label{pf:liftinj:resmudef}
Let $p\ge1$, and let $w_p$, $m(a)$, $n(b)$ be as in \cref{sec:pf:liftinj:schr}.
\begin{enumerate}[label=(\alph*)]
\item A function on $G_p$ is \emph{of type $\rho$} if it is continuous and of law $(1,p/2)$. A \emph{residue measure} is a map $\mathfrak m$ from
functions on $G_p$ to $\C$ which is additive and homogeneous on functions of type $\rho$, takes non-negative real values on non-negative
functions of type $\rho$, and satisfies: for every height $H$ there is $C$ with $|\mathfrak m(f)|\le CB$ whenever $f$ is of type $\rho$ and
$|f|\le BH^p$; $\mathfrak m(f(\cdot\,x))=\mathfrak m(f)$ for $x=n(b)$ and for $x\in\Gamma_p$; and for every jointly continuous family $(f_\sigma)$,
$\sigma\in[p/2,p/2+1]$, with $f_\sigma$ of law $(1,\sigma)$, if $(\sigma-p/2)E(1,f_\sigma)\to a$ as $\sigma\to p/2^+$ then $a=\mathfrak m(f_{p/2})$.
\item For a finite set $S$ and an $S$-idempotent $e$, $\mathrm{Herm}_S(e)=\{x\in\mathrm{Herm}_p(\A_F):ex=x\}$, and the \emph{cell} of $e$ is the
open set $\Omega_e$ of the $g\in G_p$ for which $(1-e)+e\,c(g)$ is invertible, where $c(g)$ is the lower left block of $g$.
\item Let $G$ be a topological group, $P$ a subgroup and $h$ a function on $G$. A \emph{section} for $(P,h)$ is a continuous $f$ with
$f(qg)h(g)=f(g)h(qg)$ for $q\in P$. A \emph{setup} $(P,h,\mathfrak m)$ consists of a closed subgroup $P$ such that $G=PC$ for a compact set $C$, a
continuous positive function $h$ with $h(qg)h(g')=h(g)h(qg')$ for $q\in P$, and a map $\mathfrak m$ which is linear and positive on sections. A
finite measure $\nu$ on $G/P$ \emph{represents} $\mathfrak m$ if $\mathfrak m(f)=\int(f/h)\,d\nu$ for all sections $f$ (with integrability); $\mathfrak m$ is
\emph{invariant under $x$} if $\mathfrak m(f(\cdot\,x))=\mathfrak m(f)$ for all sections $f$. An open set $\Omega\subseteq G$ has \emph{few exceptions} for
$\mathfrak m$ if every $g\in G$ has a left $P$-stable open neighbourhood $U$ and an integer $d$ such that for every neighbourhood $V$ of $1$ and every
$M$ there are $\gamma_1,\dots,\gamma_M\in V$ under which $\mathfrak m$ is invariant with $\#\{j:x\gamma_j\notin\Omega\}\le d$ for all $x\in U$. A
\emph{cell} $(\Omega,\tau,w,n)$ consists of an open left $P$-stable set $\Omega$, a continuous additive map $n$ from a locally compact group $N$ to $G$
and a continuous map $\tau:\Omega\to N$ with $\tau(qx)=\tau(x)$ for $q\in P$, $xn(y)\in\Omega$ and $\tau(xn(y))=\tau(x)+y$.
\item For $B$, $b_0\in\mathrm{Herm}_r(E)$ and $z=(z_1,z_2)\in\mathrm{Mat}_{\dim B\times r}(\A_E)\times\mathrm{Mat}_{\dim B\times t}(\A_E)$ put
$y(z)=b_0^{-1}\langle z_1,z_2\rangle_B$ and $q(z)=\langle z_2,z_2\rangle_B-{}^t\bar y(z)\,b_0\,y(z)$. A measure $\nu$ on such pairs is
\emph{orbital} for $(B,b_0,e)$ if it is finite on compact sets, vanishes outside the set of $z$ with $ez_1=z_1$, $ez_2=z_2$, $\langle z_1,z_1\rangle_B=eb_0$ and
$q(z)=0$, is invariant under $z\mapsto(hz_1,hz_2)$ for $h\in G_B$ with $e(h-1)=h-1$ and under $z\mapsto(z_1,z_2+z_1y)$ for $y$ with $ey=y$, and
satisfies $\nu\circ(z\mapsto(z_1,z_2g))^{-1}=|\det g|^{-(r+t)}\nu$ for $g\in\GL_t(\A_E)$ with $e(g-1)=g-1$.
\end{enumerate}
\formal{PositivityTheory.Rg.Rsw.IsRho}
\formal{PositivityTheory.Rg.Rsw.IsResMu}
\formal{PositivityTheory.Rg.Rsw.L.hermS}
\formal{PositivityTheory.Rg.Rsw.Cell.cell}
\formal{PositivityTheory.Rg.Rsw.Cell.Sec}
\formal{PositivityTheory.Rg.Rsw.Cell.Setup}
\formal{PositivityTheory.Rg.Rsw.Cell.Repr}
\formal{PositivityTheory.Rg.Rsw.Cell.Inv}
\formal{PositivityTheory.Rg.Rsw.Cell.Few}
\formal{PositivityTheory.Rg.Rsw.Cell.IsCell}
\formal{PositivityTheory.Rg.Rsw.L.IsOrb}
\end{definition}

\begin{proposition}[Residue measures on cells]\label{pf:liftinj:mucell}
Let $\mathfrak m$ be a residue measure on $G_p$, $n(\cdot)$ a unipotent coordinate and $w=w_p$. For every finite set $S$, every $S$-idempotent $e$ and
every Haar measure $\lambda$ on $\mathrm{Herm}_S(e)$, either $\mathfrak m(f)=0$ for every right $K^S$-invariant function $f$ of type $\rho$, or there is
$c>0$ such that for every right $K^S$-invariant function $f$ of type $\rho$ the function $x\mapsto f(w\,n(x))$ is $\lambda$-integrable and
\[
\mathfrak m(f)=c\int_{\mathrm{Herm}_S(e)}f(w\,n(x))\,d\lambda(x).
\]
\formallabel{PositivityTheory.Rg.Rsw.Cell.muCell}
\formallabel{PositivityTheory.Rg.Rsw.Cell.muCell_of}
\end{proposition}

\begin{proof}
We use the following results of the development; $c(g)$ and $d(g)$ denote the lower blocks of $g\in G_p$.
\begin{enumerate}[label=(C\arabic*)]
\item\label{pf:liftinj:C1} (\lref{GlobalTheory.CompanionTheory.SetCell.setC}, from the existence of heights.) There is a function $h$ on
$G_p$ such that $(P_p,h,\mathfrak m)$ is a setup, the sections for $(P_p,h)$ are exactly the functions of type $\rho$, $h$ is right $K^S$-invariant
for every $S$, and $\mathfrak m$ is invariant under $n(y)$ and under $w\,n(y)\,w^{-1}$ for all $y$.
\item\label{pf:liftinj:C2} (\lref{GlobalTheory.CompanionTheory.RswCell.riesz}.) Every setup has a representing measure.
\item\label{pf:liftinj:C3} (\lref{GlobalTheory.CompanionTheory.CellTop.cellAlg}.) There is $\tau:G_p\to\mathrm{Herm}_S(e)$ such that
$(\Omega_e,\tau,w,n)$ is a cell, $\Omega_e$ is right $K^S$-stable, and $e\,c(g)\,\tau(g)=e\,d(g)$ for $g\in\Omega_e$.
\item\label{pf:liftinj:C4} (\lref{GlobalTheory.CompanionTheory.CellFew.fewC}.) If $\mathfrak m$ is invariant under $w\,n(y)\,w^{-1}$ for all $y$,
then $\Omega_e$ has few exceptions for $\mathfrak m$.
\item\label{pf:liftinj:C5} (\lref{GlobalTheory.CompanionTheory.RswCell.NullCompl.nullC_of_invM}, using
\lref{GlobalTheory.CompanionTheory.RswCellInv.invM}.) If $\nu$ represents $\mathfrak m$ and $\Omega$ is open, left $P$-stable and has few
exceptions, then $\nu$ vanishes on the image of the complement of $\Omega$.
\item\label{pf:liftinj:C6} (\lref{GlobalTheory.CompanionTheory.CellDecomp.decomp}.) If $g\in\Omega_e$, $y\in\mathrm{Herm}_S(e)$ and
$e\,c(g)\,y=e\,d(g)$, then $g=q\,w\,n(y)\,u$ with $q\in P_p$ and $u\in K^S$ commuting with $n(x)$ for all $x\in\mathrm{Herm}_S(e)$.
\item\label{pf:liftinj:C7} (\lref{GlobalTheory.CompanionTheory.RswCell.cellFormula_of_invM}.) Let $\nu$ represent $\mathfrak m$, let
$(\Omega,\tau,w,n)$ be a cell with $\nu$ vanishing on the complement of $\Omega$, and assume that $\mathfrak m$ is invariant under all $n(y)$ and that
\begin{equation}\label{eq:pf:liftinj:cocycle}
h(x\,n(y))\,h(w\,n(\tau x))=h(w\,n(\tau x+y))\,h(x)\qquad(x\in\Omega).
\end{equation}
Then there is $c\ge0$ such that $\mathfrak m(f)=c\int f(w\,n(y))\,d\lambda(y)$, with integrability, for every section $f$ satisfying
$f(x)\,h(w\,n(\tau x))=f(w\,n(\tau x))\,h(x)$ for $x\in\Omega$.
\end{enumerate}
Let $h$ be as in \ref{pf:liftinj:C1}, $\nu$ a representing measure \ref{pf:liftinj:C2} and $\tau$ as in \ref{pf:liftinj:C3}. By \ref{pf:liftinj:C4}
and \ref{pf:liftinj:C5}, $\nu$ vanishes on the complement of $\Omega_e$. For $x\in\Omega_e$ write, by \ref{pf:liftinj:C6} applied with $y=\tau(x)$,
$x=q\,w\,n(\tau x)\,u$. We verify \eqref{eq:pf:liftinj:cocycle}: for $y\in\mathrm{Herm}_S(e)$, since $n(\tau x+y)=n(\tau x)n(y)$ and $u$ commutes with
$n(y)$, we have $x\,n(y)=q\,w\,n(\tau x+y)\,u$; by right $K^S$-invariance of $h$, $h(xn(y))=h(q\,w\,n(\tau x+y))$ and $h(x)=h(q\,w\,n(\tau x))$; and
the law of $h$ gives $h(q\,w\,n(\tau x+y))\,h(w\,n(\tau x))=h(w\,n(\tau x+y))\,h(q\,w\,n(\tau x))$. Hence \ref{pf:liftinj:C7} applies and gives $c\ge0$.
If $f$ is of type $\rho$ and right $K^S$-invariant, then $f(x)=f(q\,w\,n(\tau x))$ and $h(x)=h(q\,w\,n(\tau x))$, and the section property of $f$ gives
$f(x)\,h(w\,n(\tau x))=f(w\,n(\tau x))\,h(x)$; so the formula of \ref{pf:liftinj:C7} holds for $f$. If $c=0$ we are in the first case; otherwise
$c>0$, and dividing the integrability statement for the measure $c\lambda$ by $c$ gives the second case.
\end{proof}

\begin{remark}\label{pf:liftinj:mucellrem}
In the formal statement the conclusion is asserted for all finite sets $S$ containing some finite set $S_0$, and the proof takes $S_0=\emptyset$; the
existence of heights and of $S$-idempotents are hypotheses of the formal statement and enter through \ref{pf:liftinj:C1} and \ref{pf:liftinj:C3},
\ref{pf:liftinj:C4}, \ref{pf:liftinj:C6}.
\end{remark}

\begin{proposition}[Eisenstein measures]\label{pf:liftinj:eismeas}
Assume the following statements, each of which is a result of the development: \ref{pf:liftinj:R3} (idempotents), \ref{pf:liftinj:R5} (invariance of good
functions), \ref{pf:liftinj:R6} (majorants), \cref{pf:liftinj:mucell}, \cref{pf:liftinj:fib}, and:
\begin{enumerate}[label=(\roman*)]
\item (bumps) for every $S$-idempotent $e$ and every compact set $C$ of matrices in $\mathrm{Mat}_{a\times b}(\A_E)$ which are integral outside $S$,
there are compactly supported $(S,e)$-good Schwartz functions $\eta_1,\dots,\eta_a$ on row vectors with values in $[0,1]$ such that $\eta_i$ is $1$ on
the $i$-th rows of the elements of $C$;
\item (invariance of theta sections) for large $S$, every theta section $\Omega$ with $(S,e)$-good $\Omega(1)$ satisfies $\Omega(hu)=\Omega(h)$ for
$u\in K^S$;
\item (product theta sections) for $b_0$ diagonal and non-singular, $\chi$ a $\rho$-adapted Hecke character with $\rho\ge1$ and Schwartz functions
$\eta_1,\dots,\eta_\rho$ on row vectors, there is a theta section $\Omega$ with $\Omega(1)(y)=\prod_i\eta_i(y_i)$, $y_i$ the rows of $y$;
\item (the cell identity) in the setting of (RC), with Fourier--Jacobi maps and measures, $b_0$ non-singular, measures $\mu_r,\mu_t$ of covolume one,
a measure $\mathfrak f$ as in \cref{pf:liftinj:fib} for $(B,b_0,\mu_r,\lambda_z)$, a theta section $\Omega$ and $\Phi\in\mathcal S(B^n)$:
$z\mapsto\Phi(j(z))$ is $\mathfrak f\times\mu_t$-integrable and for every $x\in\mathrm{Herm}_t(\A_F)$
\[
\mathrm R(\varphi_\Phi,1,\Omega)(w_t\,n(x))=\int\Phi(j(z))\,\overline{\Omega(1)(y(z))}\,\psi\bigl(\tfrac12\tr(x\,q(z))\bigr)\,d(\mathfrak f\times\mu_t)(z),
\]
where $\varphi_\Phi(g)=(\omega_B(g)\Phi)(0)$ and $j(z_1,z_2)$ is the matrix with column blocks $z_1,z_2$;
\item (positive functionals) on a locally compact second countable space, every positive linear functional on a point-separating $*$-algebra of
compactly supported continuous functions containing, for every compact set, a non-negative function equal to $1$ on it, is the integral against a
measure finite on compact sets, which is determined by its integrals on the algebra;
\item (approximate identities) for every $S$-idempotent $e$ and Haar measure $\lambda$ on $\mathrm{Herm}_S(e)$ there are integrable functions
$g_k$ with $|g_k|\le1$ and $g_k\to1$ pointwise, whose transforms $G_k(T)=\int g_k(x)\psi(\frac12\tr(xT))\,d\lambda(x)$ are non-negative on hermitian
$T$ and tend to $0$ uniformly on every closed set of non-zero hermitian $T$ with $eT=T$.
\end{enumerate}
Then in the setting of (RC), for every residue measure $\mathfrak m$ on $G_t$, Fourier--Jacobi maps and measures and $b_0\in\mathrm{Herm}_{\dim W}(E)$
diagonal and non-singular such that $\mathrm R(\varphi_\Phi,1,\Omega)$ is of type $\rho$ for every theta section $\Omega$ and every
$\Phi\in\mathcal S(B^n)$, there is a finite set $S_0$ such that for every finite $S\supseteq S_0$ and every $S$-idempotent $e$ there is an orbital measure
$\nu$ for $(B,b_0,e)$ with
\[
\mathfrak m\bigl(\mathrm R(\varphi_\Phi,1,\Omega)\bigr)=\int\Phi(j(z))\,\overline{\Omega(1)(y(z))}\,d\nu(z)
\]
for every $(S,e)$-good $\Phi\in\mathcal S(B^n)$ and every theta section $\Omega$ for which $\Omega(1)$ is $(S,e)$-good and compactly supported.
\formallabel{PositivityTheory.Rg.Rsw.Em.eisMeas}
\formallabel{GlobalTheory.CompanionTheory.EisE.eisMeas}
\end{proposition}

\begin{proof}
From \ref{pf:liftinj:R3}, (i), (v) and (vi) the development first derives auxiliary statements: topological properties of the spaces involved (local compactness, second countability), the richness of an algebra $\mathcal A_{S,e}$ of test functions on pairs $z$ generated by the
functions $z\mapsto\Phi(j(z))$ with $\Phi$ good and compactly supported (\lref{GlobalTheory.CompanionTheory.AlgRich.algRich}), a law for the
measure $\sigma_S(\mathfrak f\times\mu_t)$ on the pairs $z$ with $ez_1=z_1$, $ez_2=z_2$, obtained from $\mathfrak f\times\mu_t$ by restriction to the pairs
which are integral outside $S$ and multiplication by $e$ (\lref{GlobalTheory.CompanionTheory.SigL.main}),
and, from (v) and (vi) and two further lemmas (\lref{GlobalTheory.CompanionTheory.EisDens.dens_main},
\lref{GlobalTheory.CompanionTheory.AlgStab.algStab}, \lref{GlobalTheory.CompanionTheory.HermSc.hermScale_main}), the existence of orbital
measures (\lref{GlobalTheory.CompanionTheory.OrbExist.orbEx}): if $x\mapsto\int f(z)\psi(\frac12\tr(x\,q(z)))\,d\sigma_S(z)$ is $\lambda$-integrable for every
$f\in\mathcal A_{S,e}$, then there is an orbital measure $\nu_0$ with
\[
\int_{\mathrm{Herm}_S(e)}\int f(z)\,\psi\bigl(\tfrac12\tr(x\,q(z))\bigr)\,d\sigma_S(z)\,d\lambda(x)=\int f\,d\nu_0
\]
for every continuous $f$ which is integrable in the same sense and can be approximated by compactly supported continuous functions up to
$\eps f_0$, $f_0\ge0$ another such function.

Now let the data be given, let $w=w_t$ and $n(\cdot)$ a unipotent coordinate. Let $S_2$ be given by \ref{pf:liftinj:R5} for $\omega_B$, $S_3$ by (ii),
and $S_4$ by an integrality statement for $y(z)$ (\lref{GlobalTheory.CompanionTheory.ColumnSplitting.YvInt.intOff_yv}); put
$S_0=S_2\cup S_3\cup S_4$. Let $S\supseteq S_0$ and $e$ an
$S$-idempotent, and for $(S,e)$-good $\Phi$ and $\Omega(1)$ put $f_{\Phi,\Omega}=\mathrm R(\varphi_\Phi,1,\Omega)$. By \ref{pf:liftinj:R5} and (ii),
$f_{\Phi,\Omega}$ is right $K^S$-invariant: for $u\in K^S$ the element $\mathrm{emb}(u)$ of $G_n$ lies in $K^S$, so that
$\omega_B(\mathrm{emb}(u))\Phi=\Phi$ and hence $\varphi_\Phi(x\,\mathrm{emb}(u))=\varphi_\Phi(x)$, while $\Omega(hu)=\Omega(h)$; inserting this in the definition of
$\mathrm R(\varphi_\Phi,1,\Omega)(hu)$ gives $\mathrm R(\varphi_\Phi,1,\Omega)(h)$. It is of type $\rho$ by assumption.

If $\mathfrak m$ vanishes on right $K^S$-invariant functions of type $\rho$ (first case of \cref{pf:liftinj:mucell}), the zero measure, which is orbital,
has the required property. Otherwise there is $c>0$ with $\mathfrak m(f_{\Phi,\Omega})=c\int f_{\Phi,\Omega}(w\,n(x))\,d\lambda(x)$. Choose $\mu_r,\mu_t$ of covolume
one and $\mathfrak f$ by \cref{pf:liftinj:fib}. By (iv) and a computation of the restriction to $\mathrm{Herm}_S(e)$
(\lref{GlobalTheory.CompanionTheory.EisE.cell_id}),
\[
f_{\Phi,\Omega}(w\,n(x))=\int\xi_{\Phi,\Omega}(z)\,\psi\bigl(\tfrac12\tr(x\,q(z))\bigr)\,d\sigma_S(z),\qquad
\xi_{\Phi,\Omega}(z)=\Phi(j(z))\,\overline{\Omega(1)(y(z))},
\]
for $x\in\mathrm{Herm}_S(e)$, so that $\mathfrak m(f_{\Phi,\Omega})=c\int\int\xi_{\Phi,\Omega}(z)\psi(\frac12\tr(xq(z)))\,d\sigma_S(z)\,d\lambda(x)$. The integrability
hypothesis of the existence of orbital measures is verified for the generators of $\mathcal A_{S,e}$ by the same identity, applied to the theta sections
given by (i) and (iii) which are equal to $1$ on the relevant compact sets; and $\xi_{\Phi,\Omega}$ is approximated by the functions $\xi_{\Phi_1,\Omega}$
with $\Phi_1$ good and compactly supported, up to $\eps\,\xi_{\Phi_0,\Omega_0}$ for a non-negative pair $(\Phi_0,\Omega_0)$ dominating $(\Phi,\Omega)$. Hence
$\mathfrak m(f_{\Phi,\Omega})=c\int\xi_{\Phi,\Omega}\,d\nu_0$, and $\nu=c\,\nu_0$ is orbital and has the required property.
\end{proof}

\subsubsection{Hecke operators on the moduli spaces}\label{sec:pf:liftinj:heckesec}

We prove the three statements of \cref{pf:liftinj:hecke}. They concern the moduli spaces of abelian varieties with action of a CM field which
the development attaches to a CM field $E$, a host, a datum, a level and a coefficient datum; we write $X$ for the moduli space, $A\to X$ for the
universal abelian scheme with its $E$-action and its line bundle $\mathcal L$, $F_c$ for the coefficient field, and $m_\tau$ ($\tau:E\to\C$) for the
multiplicities of the datum. Let $I$ be the set consisting of a symbol $\infty$ and the pairs $(\tau,i)$ with $1\le i\le m_\tau$.

\begin{definition}[Framed points, sections and Hecke transforms]\label{pf:liftinj:frdef}
\begin{enumerate}[label=(\alph*)]
\item Let $S$ be an $F_c$-algebra. A \emph{framed point} over $S$ is a point $x\in X(S)$ together with tangent vectors $u(\tau,i)$ of $A_x$ along the
zero section ($S[\eps]$-points, $\eps^2=0$) on which $E$ acts through $\tau$, such that every such tangent vector is a unique $S$-linear combination
of the $u(\tau,i)$, $1\le i\le m_\tau$, and a generator $s$ of the sections of $\mathcal L$ at $x$ (every section is a unique multiple of $s$). A
homomorphism $\phi:S\to S'$ carries framed points over $S$ to framed points over $S'$. For $t\in(S^\times)^I$, the \emph{Borel frame change} $y\cdot t$ of
$y=(x,u,s)$ is the framed point with the same $x$, with $u'(\tau,i)=t(\tau,i)u(\tau,i)+\sum_{j<i}c_ju(\tau,j)$ for some $c_j\in S$, and with
$s'=t(\infty)s$.
\item A \emph{test ring} is a quotient of a polynomial ring over $F_c$ in finitely many variables by a radical ideal. For $e:I\to\N$, a \emph{section of
weight $e$} is a rule $f$ assigning to every framed point $y$ over a test ring $R$ an element $f(y)\in R$, compatibly with homomorphisms of test
rings, such that $f(y\cdot t)=\prod_{\alpha\in I}t(\alpha)^{e(\alpha)}f(y)$. If a framed point $y$ over an $F_c$-algebra is the image of a framed
point $y'$ over a test ring, $f(y)$ denotes the image of $f(y')$. The weight attached to $(a,n)$ is $e(\infty)=n$, $e(\tau,i)=a(\tau,i)$.
\item A \emph{Hecke prime} is a maximal ideal $P$ of $O_E$ with $P+\bar P=O_E$, lying over a prime $p$ which does not divide the level. For a framed
point $y$ over an algebraically closed field and $i\ge0$, the development defines the \emph{Hecke kernels} $C$ of type $(P,p,i)$ at $y$ (certain
subgroups of $A_y$) and the \emph{Hecke quotients} $z$ of $y$ by $C$ (framed points with an isogeny $A_y\to A_z$ of kernel $C$, compatible with the
frames, the generators of $\mathcal L$ being related by the theta trace). For $c\in F_c$ the normalised scalar $\kappa_{P,i,n}(c)\in F_c$ satisfies
$\kappa_{P,i,n}(c)=c\,\kappa_{P,i,n}(1)$. For sections $f,g$ of the same weight, $g$ is the \emph{Hecke transform} $g=c\,T_{P,i}f$ if for every
algebraically closed field $\Omega$ and every framed point $y$ over $\Omega$
\[
g(y)=c\sum_Cf(z_C),
\]
the sum over the Hecke kernels $C$ of type $(P,p,i)$ at $y$, where $z_C$ is a Hecke quotient of $y$ by $C$ (and the term is $0$ if there is none).
\item Let $T$ be a complete discrete valuation ring of characteristic $0$ with fraction field $T_{\mathrm{fr}}$ and an embedding $F_c\to T_{\mathrm{fr}}$. A
framed point over $T_{\mathrm{fr}}$ is \emph{integral} if it is the image of a framed point over a test ring and $A_x$ extends to a semi-abelian scheme
over $T$ with a $T$-basis of the $\tau$-eigenspaces of its tangent space and a generator over $T$ of the sections of $\mathcal L$, compatible with
$u$ and $s$; it is \emph{good} if this semi-abelian scheme is proper, and \emph{algebraic} if the residue field of $T$ is finite. For a prime $\lambda$
of $F_c$, $T$ lies \emph{over $\lambda$} if $v_T(a)=e\,v_\lambda(a)$ for $a\in F_c$, for some integer $e>0$.
\item A framed point $z$ over a field is \emph{Hecke-$p$ over} a framed point $y$ if all points of order $p^2$ of $A_y$ and all points of order $p^3$
of $A_z$ are rational and $y,z$ are related by a pair of isogenies of the type of the Hecke correspondences at $p$, compatible with the frames.
\end{enumerate}
\formal{Formalism.L2.Pj.FrPt}
\formal{Formalism.L2.Pj.FrMap}
\formal{Formalism.L2.Pj.FrBor}
\formal{Formalism.L2.Pj.TRing}
\formal{Formalism.L2.Pj.secW}
\formal{Formalism.L2.Pj.evalAt}
\formal{Formalism.L2.Pj.ex}
\formal{GlobalTheory.IndefiniteRowSideFullHecke.UnusedSlot39.IsHeckePrime}
\formal{Formalism.L2.Pj.hkKer}
\formal{Formalism.L2.Pj.hkRelFr}
\formal{Formalism.L2.Pj.HkRel}
\formal{Formalism.L2.Pj.IntFr}
\formal{GlobalTheory.GoodPts.Good}
\formal{Formalism.L2.Pj.OverLa}
\formal{GlobalTheory.HkIntegral.HkOver}
\end{definition}

\begin{lemma}[Theta unit]\label{pf:liftinj:thetaunit}
Let $T$ be a complete discrete valuation ring with finite residue field in which the prime $p$ is invertible, let $y,y'$ be integral framed points
over $T_{\mathrm{fr}}$, $z$ a framed point over $T_{\mathrm{fr}}$ and $t\in(T_{\mathrm{fr}}^\times)^I$. If $z$ is Hecke-$p$ over $y$ and $z=y'\cdot t$,
then $v_T(t(\infty))=0$.
\formallabel{GlobalTheory.HkIntegral.theta_unit}
\end{lemma}

\begin{proof}
We use three results. By \lref{GlobalTheory.IntegralPull.exists_ext_all_models_uniformised}, for every algebraic integral framed point $y_0$
over $T$ there are a finite extension $T'$ of $T$ with $v_{T'}|_{T_{\mathrm{fr}}}=e\,v_T$ for some $e>0$ and an integral framed point $y_1$ over $T'$
extending $y_0$ all of whose integral models admit a $p$-adic uniformisation compatible with torsion. By
\lref{GlobalTheory.HkIntegral.theta_high}, if $z_1$ is Hecke-$p$ over an integral framed point $y_1$ and $z_1=y_1'\cdot t_1$ for an integral
framed point $y_1'$ all of whose integral models are uniformised in this sense, then $v(t_1(\infty))\le0$. By
\lref{GlobalTheory.HkIntegral.theta_low}, under the same relations, if all integral models of $y_1$ are uniformised, then $t_1(\infty)$ lies in the
valuation ring.

For the inequality $v_T(t(\infty))\le0$, apply the first result to $y'$ and extend $y$ to an integral framed point over $T'$ by base change; the
images $z_1$ of $z$ and $t_1$ of $t$ over $T'$ satisfy the same relations, because the Hecke-$p$ relation and Borel frame changes are preserved under
extension of scalars (\lref{GlobalTheory.HkIntegral.frBor_map}), and $p$ remains invertible. The second result gives
$e\,v_T(t(\infty))=v_{T'}(t_1(\infty))\le0$. For the inequality $v_T(t(\infty))\ge0$, apply the first result to $y$ instead and the third result; then
$e\,v_T(t(\infty))\ge0$.
\end{proof}

For a class $\mathcal C$ of algebraic integral framed points (a property of such points over all complete discrete valuation rings), a \emph{Hecke point}
of $\mathcal C$ is a framed point $z$ over an algebraically closed field $\Omega$ for which there are: a prime $\lambda$ of $F_c$, a Hecke prime $P$ over
$p$ with $p\notin\lambda$, an integer $i$, a complete discrete valuation ring $T$ over $\lambda$, an algebraic integral framed point $y$ over $T$ in
$\mathcal C$, an embedding $k:T_{\mathrm{fr}}\to\Omega$ over $F_c$ such that $\Omega$ is algebraic over $T_{\mathrm{fr}}$, and a Hecke kernel $C$ of type
$(P,p,i)$ at the image $y_\Omega$ of $y$ such that $z$ is a Hecke quotient of $y_\Omega$ by $C$. We consider two properties of $\mathcal C$:
\begin{itemize}
\item \emph{frame integrality} $\mathrm{FI}(\mathcal C)$: for every Hecke point $z$ of $\mathcal C$, with data as above, there are a finite extension $T'$
of $T$, an extension $j:T'_{\mathrm{fr}}\to\Omega$ of $k$, an algebraic integral framed point $y'$ over $T'$ in $\mathcal C$ and units $u\in(T'^\times)^I$
with $z=j(y')\cdot j(u)$;
\item \emph{value integrality} $\mathrm{VI}(\mathcal C)$: every Hecke point $z$ of $\mathcal C$ is the image of a framed point over a test ring, and for every
section $f$ of weight $(a,n)$ whose normalised valuations at all algebraic integral framed points over $\lambda$ are non-negative, $f(z)$ is integral over
$T$ (via $k$).
\end{itemize}

\begin{lemma}[Values from frames]\label{pf:liftinj:hvalint}
For every class $\mathcal C$, $\mathrm{FI}(\mathcal C)$ implies $\mathrm{VI}(\mathcal C)$.
\formallabel{GlobalTheory.HkIntegral.heckeValueInt}
\end{lemma}

\begin{proof}
Let $z$ be a Hecke point of $\mathcal C$ with data as above, $T$ lying over $\lambda$ with index $e$. Every framed point over a local ring, in particular $z$,
is the image of a framed point over a test ring (\lref{GlobalTheory.FlagTan.fromTest_of_isLocalRing}). Let $f$ be a section of weight $(a,n)$ with
non-negative valuations, let $e_{a,n}$ be the weight attached to $(a,n)$, and let $T',j,y',u$ be given by $\mathrm{FI}(\mathcal C)$, with
$v_{T'}|_{T_{\mathrm{fr}}}=r\,v_T$. Since $y'$ is an algebraic integral framed point over $T'$, which lies over $\lambda$ with index $re$, the normalised
valuation of $f$ at $y'$ is non-negative, so $f(y')\in T'$, say $f(y')=b$. By naturality $f(j(y'))=j(b)$, and by the weight law
$f(z)=f(j(y')\cdot j(u))=j\bigl(\prod_\alpha u(\alpha)^{e_{a,n}(\alpha)}\,b\bigr)$, with $\prod_\alpha u(\alpha)^{e_{a,n}(\alpha)}b\in T'$. As $T'$ is finite and
integral over $T$ (\lref{GlobalTheory.T.ringHom_finite_and_isIntegral_of_valuation_scaled}), this element is a root of a monic polynomial over $T$;
applying $j$, which extends $k$, shows that $f(z)$ is integral over $T$ via $k$.
\end{proof}

\begin{lemma}[Frame integrality]\label{pf:liftinj:hframeint}
Let $\mathcal C$ be a class such that $y\in\mathcal C$ and ($y$ good $\Rightarrow$ $y'$ good) imply $y'\in\mathcal C$. Then $\mathrm{FI}(\mathcal C)$ holds.
\formallabel{GlobalTheory.HkIntegral.heckeFrameInt}
\end{lemma}

\begin{proof}
Let $z$ be a Hecke point of $\mathcal C$ with data $\lambda,P,p,i,T,y,k,C$. The steps are as follows.
\begin{enumerate}[label=\arabic*.]
\item By \lref{GlobalTheory.T.exists_linear_heckeIsogeny_inverse_up_to_p_sq}, $y_\Omega$ and $z$ are related by a pair of isogenies of the type
of the Hecke correspondences at $p$ (as in \cref{pf:liftinj:frdef}(e)). By \lref{GlobalTheory.HkIntegral.exists_hkOver}, there are a finite
extension $L$ of $T_{\mathrm{fr}}$ with an
embedding $j_1:L\to\Omega$ extending $k$ and framed points $y_1,z_1$ over $L$, with $y_1$ the image of $y$ and $z$ the image of $z_1$, such that $z_1$
is Hecke-$p$ over $y_1$.
\item There is a complete discrete valuation ring $T_1$ with finite residue field and fraction field $L$, whose valuation restricts to a positive
multiple of $v_T$ (\lref{GlobalTheory.PlaceExtField.exists_cdvr_finite_extension}).
\item By \lref{GlobalTheory.HkIntegral.intFrExt}, after a further finite extension $T_2$ of $T_1$ (again with a positive multiple of the valuation),
there is an integral framed point $y_0$ over $T_2$ whose point is the image of the point of $z_1$. Let $j:T_{2,\mathrm{fr}}\to\Omega$ extend $j_1$
(an embedding of a field into an algebraically closed field extends to every finite extension), let $y_2$ be the integral
framed point over $T_2$ obtained from $y$ by base change and $z_2$ the image of $z_1$ over $T_2$.
\item The image $z_2$ is Hecke-$p$ over $y_2$, since these relations are preserved by the extension of scalars from $L$ to $T_{2,\mathrm{fr}}$. By
\lref{GlobalTheory.HkIntegral.exists_frBor_unit}, applied to $y_2$, $y_0$ and $z_2$, which rests on \cref{pf:liftinj:thetaunit}, there are an
integral framed point $y'$ over $T_2$ with ($y_2$ good $\Rightarrow$ $y'$ good) and units $u\in(T_2^\times)^I$ with $z_2=y'\cdot u$.
\end{enumerate}
Good reduction is preserved under the extension $T\to T_2$ (\lref{GlobalTheory.T.goodReduction_stable_valuationRing_extension}), so ($y$ good
$\Rightarrow$ $y_2$ good $\Rightarrow$ $y'$ good), and $y'\in\mathcal C$ by the hypothesis on $\mathcal C$; $y'$ is algebraic since $T_2$ has finite residue
field. Applying $j$ gives $z=j(z_2)=j(y')\cdot j(u)$. Since $T_2$ is a finite extension of $T$ and $j$ extends $k$, this is $\mathrm{FI}(\mathcal C)$ at $z$.
\end{proof}

\begin{proposition}[Hecke integrality]\label{pf:liftinj:hint}
For every CM field $E$ and all host, datum, level, moduli space and coefficient datum, frame integrality and value integrality hold for the class of
good algebraic integral framed points and for the class of all algebraic integral framed points.
\formallabel{GlobalTheory.HkIntegral.heckeInt}
\end{proposition}

\begin{proof}
Both classes satisfy the hypothesis of \cref{pf:liftinj:hframeint}: for the class of good points it is the implication itself, and for the class of
all points it is empty. Hence $\mathrm{FI}$ holds for both classes by \cref{pf:liftinj:hframeint}, and $\mathrm{VI}$ by \cref{pf:liftinj:hvalint}.
\end{proof}

\begin{proposition}[Hecke transform]\label{pf:liftinj:htr}
Assume the point laws of the development (finiteness of the Hecke kernels, existence and uniqueness of Hecke quotients, and detection of sections by
their values at points). Then:
\begin{enumerate}[label=(\alph*)]
\item for every CM field $E$, host, datum whose signature ranks satisfy $p_0+q_0\ge1$, level, moduli space and coefficient datum, every Hecke prime $P$
over $p$, every $i\le p_0+q_0$, every $c\in F_c$, every weight $a$ and every even $n$: for every section $f$ of weight $(a,n)$ there is a section $g$ of
weight $(a,n)$ with $g=\kappa_{P,i,n}(c)\,T_{P,i}f$;
\item for every global datum $D$ and core of $D$ satisfying its laws, every index $j$ in range, every Hecke prime $P$ over $p$, every $i\le r$, every
$c\in F_c$ and every dominant weight $(a,n)$: whenever the Hecke transforms of (a) exist for these data, sections of weight $(a,n)$ are detected by
their values, and the Hecke kernels are finite, the Hecke operator $f\mapsto g$ maps cusp sections (those whose pull-backs to the complex
uniformisations of the core are holomorphic cusp forms) to cusp sections.
\end{enumerate}
\formallabel{GlobalTheory.HeckeEx.heckeTransform_of_pointLaws}
\formallabel{GlobalTheory.HeckeEx.secTransform}
\formallabel{GlobalTheory.HeckeEx.cuspTransform}
\end{proposition}

\begin{proof}
(a) Let $f$ be a section of weight $(a,n)$. For a test ring $R$ and a framed point $y$ over $R$, the point laws and
\lref{GlobalTheory.HeckeEx.exists_sum} give an element $g_0(y)\in R$ such that for every algebraically closed field $\Omega$, every
$\phi:R\to\Omega$ and the image $y_\Omega$ of $y$, $\phi(g_0(y))=\sum_Cf(z_C)$, the sum over the Hecke kernels at $y_\Omega$ as in
\cref{pf:liftinj:frdef}(c). Since $R$ is reduced, an element of
$R$ is determined by its images under all homomorphisms to algebraically closed fields
(\lref{GlobalTheory.Frames.GeometricValues.eq_zero_of_forall_algClosed}). Hence $g_0$ is compatible with homomorphisms of test rings, because
both sides have the same images at all geometric points. For a Borel frame change $y\cdot t$, the sum over the Hecke kernels at $(y\cdot t)_\Omega$ is
$\prod_\alpha\phi(t(\alpha))^{e(\alpha)}$ times the sum at $y_\Omega$ (\lref{GlobalTheory.T.heckeKernel_value_sum_weight_triangular_frame}: the Hecke sum
transforms by the weight character under triangular frame changes), so $g_0$ has weight $(a,n)$. Thus $g_0$ is a section of weight $(a,n)$, and so is
$g=\kappa_{P,i,n}(c)g_0$, which satisfies $g=\kappa_{P,i,n}(c)T_{P,i}f$.

(b) Let $f$ be a cusp section and $g$ its Hecke transform. Then $g$ is a section of weight $(a,n)$, and cuspidality of $g$ is checked at each good
embedding of $F_c$ in $\C$ on each component of the complex uniformisation. By \lref{GlobalTheory.HeckeEx.rdU_mem_SS}, if the functions attached to
a section of dominant weight on the components of a complex uniformisation are holomorphic cusp forms of the corresponding weight, then so are the
functions attached to its Hecke transforms (for a Hecke prime and finite Hecke kernels). Hence $g$ is a cusp section.
\end{proof}

For the reading statement, let $D$ be a global datum, $R$ real data for $D$ and a core tied to $R$, with the standing conditions of
\cref{pf:liftinj:hecke}(c), and let $i$ be an index. The core provides the moduli space $X_{i+1}$ with coefficient field $F_c$ and a distinguished
embedding $\sigma_0:F_c\to\C$; at $\sigma_0$ the complex uniformisation attaches to every $h\in G_{W_i}$ in a certain domain a complex framed point
$x_h$, and to every section $f$ of weight $a$ a function $\varphi_W(f)$ on $G_{W_i}$, with
\begin{equation}\label{eq:pf:liftinj:phiW}
\varphi_W(f)(h)=\ell_a\,\Lambda(h)\,f(x_h)
\end{equation}
for $h$ in the domain, where $\ell_a$ is a constant and $\Lambda(h)\in\C$ a constant of the frame at $h$
(\lref{GlobalTheory.HeckeReadProof.phiW_eq}). \emph{Hecke data at a split prime} are $(P,p,k,\alpha,\mathcal R)$ where $P$ is a place of $E$ split
over $F$ outside the bad set of $R$, $p$ is the prime below $P$ and does not divide the level of the core, $k\le r$, $\alpha\in G_{W_i}$ is a Hecke
element of type $k$ at $P$ and $\mathcal R$ is a set of representatives of $K\alpha K/K$ for the level group $K$ of $W_i$; $\mu(P)\in\C^\times$ is the value at
$P$ of the character $\mu_0$ of $R$. The \emph{frame pin} at $(i,\sigma_0)$ is the identity, for a constant of absolute value one and an orthonormal
lattice basis of $V_i$, between the theta values of the vacuum vector of $R$ at the real points and products of Riemann theta functions of the complex
uniformisation, multiplied by the frame constants. We write $\kappa_i$ for the weight of the datum at $i$ and $n_i$ for its component at $\infty$.

\begin{proposition}[Hecke kernels and cosets]\label{pf:liftinj:hker}
Assume the point laws, the split linear model of $D$ and the fixedness of the vacuum stabiliser of $D$. Let $R$, the core and $i$ be as above, and
assume the frame pin at $(i,\sigma_0)$. Let $(P,p,k,\alpha,\mathcal R)$ be Hecke data at a split prime and $h\in G_{W_i}$ with $x_h$ defined. Then there is
a bijection $e$ from the set of Hecke kernels of type $(P,p,k)$ at $x_h$ onto $\mathcal R$ such that every kernel $C$ has a Hecke quotient, and for every
Hecke quotient $z$ of $x_h$ by $C$ and every section $f$ of weight $\kappa_i$
\[
\ell_{\kappa_i}\,\Lambda(h)\,\sigma_0\bigl(\kappa_{P,k,n_V}(1)\bigr)\,f(z)=\mu(P)^{-k}\,\varphi_W(f)\bigl(h\,e(C)\bigr),
\]
where $n_V$ is the rank of the datum.
\formallabel{GlobalTheory.T.heckeReading_kernels_match_cosets_through_frame}
\formallabel{GlobalTheory.HeckeTerm.heckeTerm}
\end{proposition}

\begin{proof}
By \lref{GlobalTheory.T.heckeKernel_complex_point_matches_cosets}, there is a bijection $e$ from the Hecke kernels at $x_h$ onto $\mathcal R$
such that for every kernel $C$ and every Hecke quotient $z$ of $x_h$ by $C$ the point $x_{h\,e(C)}$ is defined and $z=x_{h\,e(C)}\cdot t$ for some $t$ with
$t(\tau,j)=1$ for all $(\tau,j)$; this is the comparison of the Hecke correspondence on the complex uniformisation with the coset decomposition of the
adelic Hecke operator. Every kernel has a Hecke quotient by the point laws. Let $C$, $z$ and $t$ be as above, and $f$ a section of weight $\kappa_i$. By
\lref{GlobalTheory.T.heckeTerm_generator_constant_via_frame},
\[
t(\infty)^{n_i}\,\Lambda(h)\,\sigma_0\bigl(\kappa_{P,k,n_V}(1)\bigr)=\mu(P)^{-k}\,\Lambda\bigl(h\,e(C)\bigr),
\]
which computes the scaling of the generator of $\mathcal L$ along the Hecke isogeny in terms of the frame constants. By the weight law, since
$t(\tau,j)=1$, $f(z)=t(\infty)^{n_i}f(x_{h\,e(C)})$ (\lref{GlobalTheory.FrameEval.evR_frBor}). Multiplying by $\ell_{\kappa_i}\Lambda(h)\sigma_0(\kappa_{P,k,n_V}(1))$
and using \eqref{eq:pf:liftinj:phiW} at $h\,e(C)$ gives the assertion.
\end{proof}

\begin{proposition}[Hecke reading]\label{pf:liftinj:hread}
Assume the split linear model and the fixedness of the vacuum stabiliser for every global datum. Let $D$, $R$, the core and $i$ be as above, with the frame
pin at $(i,\sigma_0)$. Then for all Hecke data $(P,p,k,\alpha,\mathcal R)$ at a split prime, all $c\in F_c$ and all sections $f,g$ of weight $\kappa_i$ with
$g=\kappa_{P,k,n_V}(c)\,T_{P,k}f$, and for all $h\in G_{W_i}$,
\[
\varphi_W(g)(h)=\sigma_0(c)\,\mu(P)^{-k}\sum_{a\in\mathcal R}\varphi_W(f)(h\,a).
\]
\formallabel{GlobalTheory.HeckeReadProof.heckeReadTie}
\end{proposition}

\begin{proof}
Fix $h$. By the laws of the core and since $\sigma_0$ is a good embedding, $x_h$ is defined
(\lref{GlobalTheory.IndefiniteRowSideFullHecke.OrthonormalLatticeThetaValues.exists_witness}). The point laws hold
(\lref{GlobalTheory.HeckePointLaw.heckePointLaws}), so \cref{pf:liftinj:hker} applies and gives the bijection $e$; in particular the set of Hecke
kernels at $x_h$ is finite. For each kernel $C$ choose a Hecke quotient $z_C$. By \eqref{eq:pf:liftinj:phiW} and the Hecke relation at $y=x_h$ over
$\Omega=\C$ (via $\sigma_0$),
\[
\varphi_W(g)(h)=\ell_{\kappa_i}\Lambda(h)\,g(x_h)=\ell_{\kappa_i}\Lambda(h)\,\sigma_0(c)\,\sigma_0\bigl(\kappa_{P,k,n_V}(1)\bigr)\sum_Cf(z_C),
\]
using $\kappa_{P,k,n_V}(c)=c\,\kappa_{P,k,n_V}(1)$. By \cref{pf:liftinj:hker}, for every kernel $C$
\[
\ell_{\kappa_i}\Lambda(h)\,\sigma_0(c)\,\sigma_0\bigl(\kappa_{P,k,n_V}(1)\bigr)f(z_C)=\sigma_0(c)\,\mu(P)^{-k}\,\varphi_W(f)\bigl(h\,e(C)\bigr),
\]
and re-indexing the sum by the bijection $e$ gives the assertion.
\end{proof}

\subsubsection{Two steps of the point conditions}\label{sec:pf:liftinj:ptsteps}

We prove the statements used in steps 4 and 7 of the proof of \cref{pf:liftinj:point}. Let $D$ be a global datum and $h\in\mathcal R$. For a local index
$i$ of $D$, the automorphic data of $D$ provide a set $\mathcal A_V(i)$ of automorphic representations $\pi$ on the $V$-side, with local components
$\mathrm{loc}_V(i,\pi,w)$ at the places $w$ of $E$ split over $F$ and Satake parameters $\mathrm{sat}_V(i,\pi,w)$ at the other places, a notion of base
change of $\pi$ (a list of representations of general linear groups), and, for every type $\sigma$ on the $W$-side, local components
$\mathrm{loc}_W(i,\sigma,w)$ at split places; there are a pairing relation between types $\sigma$ and representations $\pi$ and a local theta lift
$\mathrm{lift}_w$ of parameters at split places. The theta data of $D$ provide isotypic parts $\mathrm{isot}(j,\sigma)$ and maps $\theta_T,\theta_P$.

\begin{definition}[Local statements and rows]\label{pf:liftinj:rowdef}
Let $i$ be a local index of $D$, with tuple $\mu_i$ and index $j_i$.
\begin{enumerate}[label=(\alph*)]
\item \emph{Pair lift} at $i$: $\mathrm{loc}_V(i,\pi,w)=\mathrm{lift}_w(\mathrm{loc}_W(i,\sigma,w))$ for all paired $(\sigma,\pi)$ and all split $w$.
\item \emph{Split base change} at $i$: if $\pi,\pi'\in\mathcal A_V(i)$ have the same local components at almost all split places, then any base changes of
$\pi$ and of $\pi'$ agree up to order.
\item \emph{Local representation} at $i$: if $\pi,\pi'\in\mathcal A_V(i)$ have the same local components at all split places and the same unramified base
change of their Satake parameters at all other places, then $\pi=\pi'$.
\item The \emph{row size} condition: for every local index $i$ and every type $\sigma$ with $\mathrm{isot}(j_i,\sigma)\ne0$, at every split place $w$ at which
$\mathrm{loc}_W(i,\sigma,w)$ is unramified, its Frobenius eigenvalues have absolute value less than $\Nm w^{(t-1)/2}$. Condition (B) of
\cref{pf:liftinj:pointdef} is the conjunction of the row size condition, the third admissibility condition (a split theta correspondence statement),
archimedean signature matching and the statement on the first theta lift.
\item \emph{Positivity of the localisation scalars} at $(\mu,j)$: for every type $\sigma$, $\theta_P\circ\theta_T$ acts on $\mathrm{isot}(j,\sigma)$ by a real
scalar $\lambda(\mu,j,\sigma)$, which is positive if $\mathrm{isot}(j,\sigma)\ne0$.
\end{enumerate}
\formal{GlobalTheory.RowTie.PairLift}
\formal{GlobalTheory.RowTie.SplitBC}
\formal{GlobalTheory.RowTie.RepOfLoc}
\formal{GlobalTheory.RowCopy.RowSize}
\formal{GlobalTheory.NRows.RowsB}
\formal{GlobalTheory.RowCopy.LamPos}
\end{definition}

\begin{proposition}[Condition (B) and the local statements]\label{pf:liftinj:bcond}
Assume the theta-local inputs $\mathcal T$. Let $h\in\mathcal R$, let $D$ be a global datum with a character choice $P$ and a join of $P$, and assume
condition (A), the first admissibility condition and the positivity of the localisation scalars at $(\mu_i,j_i)$ for every local index $i$. Then (B)
holds, and for every local index $i$ the pair lift, split base change and local representation statements hold.
\formallabel{GlobalTheory.T.b_conditions_pairLift_splitBC_repOfLoc_of_a}
\formallabel{UpperTheory.One.M.latticesB_closing_alt}
\end{proposition}

\begin{proof}
The join provides row modules for every local index. We use:
\begin{enumerate}[label=(\arabic*)]
\item a faithful reading $S$ of the tuple datum of $P$ together with its laws, from the finiteness statement in $\mathcal T$
(\lref{GlobalTheory.T.exists_true_reading_with_laws_suppLaws}); its unit property
(\lref{GlobalTheory.RowUnitProof.rowUnit_of_isTrue_alt}); and the pair reading obtained from the join and Bernstein--Zelevinsky theory in $\mathcal T$
(\lref{GlobalTheory.T.pairRead_of_join_isTrue_unit});
\item the admissibility interface (\lref{GlobalTheory.AdmIface.main_alt_alt}): the first admissibility condition implies the existence of base
changes for all $\pi\in\mathcal A_V(i)$, and the fifth admissibility condition contained in (A) implies a Chebotarev splitting statement.
\end{enumerate}
\emph{Pair lift.} By Howe duality and Bernstein's theorem contained in $\mathcal T$, the laws, the unit property and the pair reading, together with the
goodness of the representations read by $S$ (\lref{UpperTheory.One.M.repGood_of_isTrue'}), the pair lift statement holds at every $i$
(\lref{UpperTheory.One.M.pairLift_of_howeDuality_bernstein_alt}).

\emph{Split base change.} The tie at $0$, the second admissibility condition and the temperedness condition in (A) give an isomorphism of Galois
representations at $i$ (\lref{UpperTheory.GalS.galIsoN_of_tie0_adm2_adm7}); with the Chebotarev splitting statement of (2) and a further condition of
(A), this gives split base change at every $i$ (\lref{UpperTheory.GalS.splitBC_of_galIso}).

\emph{Local representation.} This holds at every $i$ by \lref{GlobalTheory.T.repOfLoc_of_isTrue_laws} (from the faithful reading and its laws).

\emph{Condition (B).} The row size condition holds because the local components read by $S$ are unitary and spherical
(\lref{GlobalTheory.SplitRead.unitarySphericalComponent_of_reading_alt}, using Bernstein--Zelevinsky theory, and
\lref{GlobalTheory.RowSizeClosing.rowSize_of_unitarySphericalComponent}). The third admissibility condition follows from the pair lift statements
(\lref{GlobalTheory.RowTie.split_theta_correspondence_of_local_pairings}). Archimedean signature matching follows from $\mathcal T$, the row modules and
(A) (\lref{GlobalTheory.T.archimedeanSignatureMatching_of_modules}). The statement on the first theta lift follows from the positivity of the
localisation scalars, the pair lift statements, the orthogonality of types given by the form of the join
(\lref{GlobalTheory.T.type_orthogonality_of_form'}), the existence of base changes, split base change and the local representation statements
(\lref{GlobalTheory.Proof.ThetaOneGenusOne.thetaOne_of_ties}).
\end{proof}

\begin{proposition}[Eigen-segments]\label{pf:liftinj:eigstr}
Assume finiteness of automorphic forms on $\GL_n$. Let $h\in\mathcal R$, let $D$ be a global datum with a character choice $P$ and real-place data $X$
agreeing with $P$, and assume the provision statement, the seventh admissibility condition and the label tie for $D$. Let $i$ be a local index of $D$
whose tuple $\mu_i$ is used, $j$ the index with $j+1=j_i$, $x$ a good vector of index $i$, $F$ an element of the eigenspace attached to $x$, and
$y\in H_V(j,\mu)$ an element for the real datum $X.R$ and the tuple $\mu$ of $X$ attached to $\mu_i$, such that $\mathrm{ext}_V(y)$ is the function attached
to $F$. If $y\ne0$, then for every Hecke family $h'$ there is a system $s$ with a segment such that $y$ is an eigenvector of $h'$ with system $s$
(\cref{pf:liftinj:realdefs}).
\formallabel{GlobalTheory.T.eigString_of_agree_aux41_adm7_tieLab}
\formallabel{GlobalTheory.EigStr.eigString_all_of_aux41_adm7_tieLab}
\end{proposition}

\begin{proof}
Let $h'$ be a Hecke family. We use six compatibility statements of the development, all under the hypotheses of the proposition; each place is split or
inert over $F$ (\lref{GlobalTheory.LPositivity.split_or_inert}).
\begin{enumerate}[label=(Q\arabic*)]
\item\label{pf:liftinj:Q1} At every good place $w$ above a split place of $F$ and for every Hecke label $d$ at $w$, $h'(w,\delta_d)y$ is a multiple of $y$
(\lref{GlobalTheory.T.theta_eigen_compatibility_first_alt}); the same holds at good places above inert places
(\lref{GlobalTheory.T.theta_eigen_compatibility_second_alt}).
\item\label{pf:liftinj:Q2} If $y$ is an eigenvector with system $s$, then at every place $w$ above a place of $F$ outside the bad set of $X.R$ there is a
multiset $B_w$ of rank $\dim V_j$ whose Satake transform gives the values of $s$ at $w$
(\lref{GlobalTheory.T.theta_eigen_compatibility_third_alt} for split, \lref{GlobalTheory.T.theta_eigen_compatibility_fourth_alt} for inert
places).
\item\label{pf:liftinj:Q3} If moreover $B$ is a Satake family of rank $\dim V_j$ for $s$ off the bad set, then at every good $w$,
$B(w)=B'(w)+\liftinjstr(\Nm w,\chi_V(w),t)$ with all elements of $B'(w)$ of absolute value at most $\Nm w^{(r-1)/2}$
(\lref{GlobalTheory.T.theta_eigen_compatibility_fifth_alt} for split, \lref{GlobalTheory.T.theta_eigen_compatibility_sixth_alt} for inert places).
\end{enumerate}
Define $s(w,d)$ to be the scalar $c$ with $h'(w,\delta_d)y=cy$ if it exists, and $0$ otherwise. By \ref{pf:liftinj:Q1}, $y$ is an eigenvector with system
$s$ (the scalar is unique since $y\ne0$). Define $B(w)$ to be a multiset as in \ref{pf:liftinj:Q2} if one exists, and the empty multiset otherwise; by
\ref{pf:liftinj:Q2}, $B$ is a Satake family of rank $\dim V_j$ for $s$ off the bad set. By \ref{pf:liftinj:Q3}, $B$ has the segment decomposition at every good
place. Hence $s$ has a segment.
\end{proof}

\subsubsection{The point statement}\label{sec:pf:liftinj:ptstsec}

The construction of the development proceeds through \emph{stages}; at each stage, every level $l$ over it (a strictly henselian acyclic level, with its
automorphic data, multiplicity spaces, character choice and tuple datum) determines global data $D'$ built from $l$ with a bound; a property of global data
holds \emph{eventually at $l$} if there is a bound such that it holds for all global data built from $l$ with this bound. The global data built at a stage
have the theta data of the datum of the stage.

\begin{lemma}[Consequences of the point conditions]\label{pf:liftinj:ptaux}
Assume the theta-local inputs $\mathcal T$, Harish-Chandra finiteness and rapid decay of cusp forms. Let $h\in\mathcal R$ and let $D$ be a global datum with
linked inputs $\mathcal I$ satisfying (J). Then the provision statement (41), residual irreducibility, conclusion (42) and conclusion (Z) hold for $D$.
\formallabel{GlobalTheory.HeightCounting.Up.PointStatements.aux_residual_aux42_prvZ_of_inputs}
\end{lemma}

\begin{proof}
By \cref{pf:liftinj:point}, the point conditions $\mathcal P(D)$ hold. The linked inputs produce the first provision statement from the summand statement
(\lref{UpperTheory.GalS.Parts.main_alt_alt}) and the conditions (A), (B) of $\mathcal P(D)$; it contains (41) and (42). Residual irreducibility follows from
$\mathcal I$ and $\mathcal P(D)$ (\lref{GlobalTheory.HeightCounting.Up.PointStatements.residualIrreducibleGood_of_inputs}). Given (42) and the isotypic
condition contained in $\mathcal I$, conclusion (Z) is equivalent to a reduced multiplicity statement
(\lref{Formalism.MainNotions.SumRepresentations.multiplicity_statement_iff_reduced}), which holds by
\lref{GlobalTheory.T.reducedMultiplicity_of_inputs}.
\end{proof}

\begin{theorem}[The point statement]\label{pf:liftinj:ptst}
The third-layer literature package $\Pi$ and the geometric compatibility package $\mathcal X$ imply, for every $h\in\mathcal R$: for every stage and every level
$l$ over it, eventually at $l$ the built global data $D'$ satisfy
\begin{enumerate}[label=(\roman*)]
\item if the largeness hypothesis (O) holds for $D'$, then (41), residual irreducibility, (42) and (Z) hold for $D'$;
\item the Hecke eigenspaces of $D'$ span its spaces (the eigen-span condition);
\item the isotypic decomposition of the theta data of $D'$ is internal (a direct sum).
\end{enumerate}
\formallabel{GlobalTheory.HeightCounting.Up.ptSt_of_inputs_geometric_compat}
\end{theorem}

\begin{proof}
We combine four statements which hold at every stage: (1) eventually at every level, (O) implies the existence of linked inputs for $D'$
(\lref{GlobalTheory.T.eventually_nonempty_inputs_of_level_c}, from $\Pi$, $\mathcal X$ and $h$); (2) eventually at every level, (O) implies (J) for $D'$
(\lref{GlobalTheory.HeightCounting.JumpsSt.locComp_of_isRc_classMapD_cenLoc}, from the Hodge--Tate cardinality statement contained in $\mathcal X$);
(3) eventually at every level, the eigen-span condition holds (\lref{GlobalTheory.HeightCounting.Up.eventually_eigSpan_of_levelC}); (4) the datum
of the stage satisfies the eighth admissibility condition, which is the internality (iii) of its isotypic decomposition, together with a condition on the
localisation scalars (\lref{GlobalTheory.T.adm8_and_lamUsed_of_geometric_galois_compat}). Two properties which hold eventually at $l$ hold jointly
eventually at $l$ (take the larger bound). Let $D'$ be built from $l$ with
a bound for which (1)--(3) hold. If (O) holds for $D'$, choose linked inputs by (1); by (2), (J) holds, and \cref{pf:liftinj:ptaux}, with $\mathcal T$,
Harish-Chandra finiteness and rapid decay taken from $\Pi$, gives (i). (ii) is (3), and (iii) follows from (4) since $D'$ has the theta data of the datum of
the stage.
\end{proof}
